% Propriétés chimiques des alcools — Chimie Tle E — Cours signé OnBuch+
% Programme officiel MINESEC, Terminale C, D, E. Compiler avec : tectonic L01-proprietes-chimiques-des-alcools.tex
\newcommand{\DOCMATIERE}{Chimie}
\newcommand{\DOCNIVEAU}{Tle E}
\newcommand{\DOCMODULE}{Module 1 · Chimie organique}
\newcommand{\DOCLECON}{01}
\newcommand{\DOCTITRE}{Propriétés chimiques des alcools}
% OnBuch+ — préambule « cours signés » ENRICHI (compilé avec tectonic / XeLaTeX)
% Système visuel « Pop » d'OnBuch+ : fond crème (jamais blanc), bordures encre
% épaisses, ombres « dures » décalées, titres Archivo Black en capitales.
% Version enrichie pour les cours de chimie : chimie (mhchem, chemfig),
% graphes (pgfplots), boîtes pédagogiques supplémentaires (définition,
% propriété, méthode, attention, le savais-tu, exercice résolu, exercice,
% TP, bilan), page de garde refaite et signature OnBuch+ ancrée en bas.
%
% Macros attendues dans main.tex AVANT \input{preamble} :
%   \DOCMATIERE  \DOCNIVEAU  \DOCMODULE  \DOCLECON (numéro)  \DOCTITRE
\documentclass[11pt]{article}

\usepackage{fontspec}
\usepackage[french]{babel}
\usepackage[a4paper,margin=2cm,top=3.4cm,bottom=2.8cm,headheight=1.4cm,footskip=1.3cm]{geometry}
\usepackage{amsmath,amssymb,amsfonts}
\usepackage[table]{xcolor} % [table] : fournit aussi \rowcolors (alternance de lignes)
\usepackage{pagecolor}
\usepackage{tikz}
\usetikzlibrary{arrows.meta,positioning,calc,shapes.geometric,shapes.misc,shapes.arrows,decorations.pathmorphing,decorations.markings,patterns,shadows,mindmap,trees,fit,backgrounds,matrix,intersections}
\usepackage{pgfplots}
\pgfplotsset{compat=1.18}
\usepackage[version=4]{mhchem}
\usepackage{chemfig}
\usepackage{enumitem}
\usepackage{fancyhdr}
% [most] charge toutes les bibliothèques tcolorbox (listings, minted, raster,
% documentation...) et gonfle inutilement la mémoire TeX sur un document long
% (~300 boîtes) : on ne charge que ce qui est réellement utilisé.
\usepackage[skins,breakable]{tcolorbox}
\usepackage{graphicx}
\usepackage{colortbl}
\usepackage{booktabs}
\usepackage{tabularx}
\usepackage{array}
\usepackage{multicol}
\usepackage{siunitx}
% parse-numbers=false : accepte aussi \SI{2,5 \times 10^{-3}}{...}
\sisetup{locale=FR,output-decimal-marker={,},group-minimum-digits=4,parse-numbers=false}
\usepackage{qrcode}
% \degree : commande fréquemment utilisée par les modèles pour un angle ou une
% température en dehors de \SI{}{} (non fournie par les paquets chargés).
\providecommand{\degree}{^{\circ}}
\usepackage{lastpage}
\usepackage{hyperref}
\hypersetup{hidelinks}

% ============================================================
% Palette « Pop » OnBuch+ — mêmes hex que la classe P (plus_ui.dart)
% ============================================================
\definecolor{poporange}{HTML}{F59321}
\definecolor{popdark}{HTML}{E07A0C}
\definecolor{popcream}{HTML}{FFF6E8}
\definecolor{popink}{HTML}{1C1714}
\definecolor{poppurple}{HTML}{7A5AE0}
\definecolor{popgreen}{HTML}{2E9E6A}
\definecolor{poppink}{HTML}{E0457B}
\definecolor{popblue}{HTML}{2F7DE1}
\definecolor{popgold}{HTML}{C98A12}
\definecolor{popmuted}{HTML}{8A7F76}
\definecolor{popline}{HTML}{EADFD2}
\definecolor{popgray}{HTML}{8A7F76}
\colorlet{popgrey}{popgray}
% Alias tolérants (le modèle nomme parfois les couleurs différemment)
\colorlet{popwhite}{white}
\colorlet{popblack}{popink}
\colorlet{popred}{poppink}
\colorlet{popyellow}{popgold}
% Teintes claires pour fonds de tableaux / zones
\colorlet{popblueL}{popblue!10!white}
\colorlet{poporangeL}{poporange!14!white}
\colorlet{popgreenL}{popgreen!12!white}
\colorlet{poppurpleL}{poppurple!12!white}
\colorlet{poppinkL}{poppink!12!white}
\colorlet{popgoldL}{popgold!14!white}

\pagecolor{popcream}

% ============================================================
% Typographie
% ============================================================
\defaultfontfeatures{Ligatures=TeX}
\setmainfont{PlusJakartaSans-Medium.ttf}[Path=../../../fonts/,BoldFont=PlusJakartaSans-Bold.ttf]
% Maths en sans-serif (Fira Math) avec chiffres et lettres droites de Plus
% Jakarta, pour que les formules se fondent dans le texte.
\usepackage[math-style=ISO,bold-style=ISO]{unicode-math}
\setmathfont{FiraMath-Regular.otf}[Path=../../../fonts/]
\setmathfont{PlusJakartaSans-Medium.ttf}[Path=../../../fonts/,range={up/num,up/{Latin,latin},\mathup}]
\usepackage{icomma} % virgule décimale sans espace en mode math
% Caractères absents de la police de texte (sinon carré vide) : rendus par
% leur équivalent mathématique, en texte comme en formule.
\usepackage{newunicodechar}
\newunicodechar{α}{\ensuremath{\alpha}}
\newunicodechar{Α}{\ensuremath{\alpha}}
\newunicodechar{β}{\ensuremath{\beta}}
\newunicodechar{δ}{\ensuremath{\delta}}
\newunicodechar{✓}{\popcheck{popgreen}}
\newfontfamily\popdisplay{ArchivoBlack-Regular.ttf}[Path=../../../fonts/]
\newfontfamily\poplogo{SpaceGrotesk-Bold.ttf}[Path=../../../fonts/]
\newfontfamily\popsemi{PlusJakartaSans-SemiBold.ttf}[Path=../../../fonts/]
\newfontfamily\popxbold{PlusJakartaSans-ExtraBold.ttf}[Path=../../../fonts/]
\newfontfamily\popxboldls{PlusJakartaSans-ExtraBold.ttf}[Path=../../../fonts/,LetterSpace=3.0]

\setlist[enumerate]{leftmargin=1.7em,itemsep=5pt,topsep=4pt}
\setlist[itemize]{leftmargin=1.7em,itemsep=4pt,topsep=4pt,label={\color{poporange}\textbullet}}
\setlength{\parindent}{0pt}
\setlength{\parskip}{5pt}
\renewcommand{\arraystretch}{1.25}

% chemfig : liaisons un peu plus épaisses, encre Pop
\setchemfig{atom sep=2em,bond style={line width=0.9pt,popink},cram width=3pt}

% pgfplots : style Pop par défaut
\pgfplotsset{
  every axis/.append style={axis line style={popink,line width=1pt},tick style={popink},
    grid style={popline,line width=0.5pt},label style={font=\small},tick label style={font=\footnotesize}},
  popcurve/.style={poporange,line width=1.8pt,no markers,smooth},
}

% ============================================================
% Pill / squiggle
% ============================================================
\newcommand{\poppill}[3]{%
  \tikz[baseline=-0.55ex]\node[draw=popink,line width=1.2pt,rounded corners=2.6pt,%
    fill=#2,inner xsep=6.5pt,inner ysep=3pt]{%
    \popxboldls\fontsize{7}{7}\selectfont\color{#3}\MakeUppercase{#1}};%
}
\newcommand{\popsquiggle}[1]{%
  \tikz[baseline=-0.4ex]{
    \draw[#1,line width=2.6pt,line cap=round]
      plot [smooth] coordinates {(0,0.09) (0.34,0.01) (0.68,0.21) (1.02,0.01) (1.36,0.10)};
  }%
}
% Mot-clé surligné (marqueur orange sous le texte)
\newcommand{\cle}[1]{{\popxbold\color{popdark}#1}}

% ============================================================
% En-tête / pied
% ============================================================
\pagestyle{fancy}
\fancyhf{}
\renewcommand{\headrulewidth}{0pt}
\renewcommand{\footrulewidth}{0pt}
\lhead{\raisebox{-0.15ex}{\poplogo\fontsize{13}{13}\selectfont\color{popink}OnBuch\color{poporange}+}}
\rhead{\poppill{\DOCMATIERE\ \textbullet\ \DOCNIVEAU\ \textbullet\ Leçon \DOCLECON}{white}{popink}}
\lfoot{\popxbold\fontsize{7.6}{7.6}\selectfont\color{popmuted}\MakeUppercase{Cours signé OnBuch+}\ \textbullet\ {\popsemi\DOCTITRE}}
\rfoot{\tikz[baseline=-0.6ex]\node[rounded corners=8.5pt,fill=popink,minimum height=17pt,minimum width=17pt,inner xsep=5pt,inner sep=0]{\popxbold\fontsize{8}{8}\selectfont\color{popcream}\thepage\,/\,\pageref*{LastPage}};}

% ============================================================
% Titres
% ============================================================
\newcommand{\chaptitre}[2]{%
  \begin{tcolorbox}[enhanced,colback=poporange,colframe=popink,boxrule=2.6pt,arc=11pt,
      drop shadow=popink,left=15pt,right=15pt,top=12pt,bottom=11pt,before skip=3pt,after skip=18pt]
    \poppill{#2}{popink}{popcream}\par\vspace{9pt}
    {\raggedright\hyphenpenalty=10000\exhyphenpenalty=10000\popdisplay\fontsize{22}{24}\selectfont\color{popcream}\MakeUppercase{#1}\par}
  \end{tcolorbox}
}
% Section numérotée : pastille orange avec le numéro + titre Archivo
\newcounter{popsec}
\newcommand{\coursec}[1]{%
  \stepcounter{popsec}\setcounter{popsub}{0}%
  \par\vspace{16pt}\noindent
  \begin{minipage}{\linewidth}
  \tikz[baseline=-0.7ex]\node[draw=popink,line width=1.6pt,fill=poporange,rounded corners=4pt,minimum size=22pt,inner sep=0]{\popdisplay\fontsize{12}{12}\selectfont\color{popcream}\thepopsec};%
  \hspace{8pt}{\popdisplay\fontsize{15}{17}\selectfont\color{popink}\MakeUppercase{#1}}\par
  \vspace{4pt}\noindent{\color{poporange}\rule{36pt}{3.6pt}}
  \end{minipage}\par\nopagebreak\vspace{8pt}
}
\newcounter{popsub}
\newcommand{\courssub}[1]{%
  \stepcounter{popsub}%
  \par\vspace{9pt}\noindent{\popxbold\fontsize{11.5}{13}\selectfont\color{popink}{\color{poporange}\thepopsec.\thepopsub}\hspace{6pt}#1}\par\nopagebreak\vspace{4pt}
}

% ============================================================
% Boîtes pédagogiques — même recette : fond blanc, bordure épaisse colorée,
% ombre dure encre, badge plein. Titre optionnel [#1].
% ============================================================
\tcbset{popbase/.style={breakable,enhanced,colback=white,boxrule=2.3pt,arc=9pt,
  shadow={3pt}{-3pt}{0pt}{popink},left=14pt,right=14pt,top=11pt,bottom=11pt,before skip=11pt,after skip=15pt,
  fontupper=\popsemi\fontsize{10.6}{14.5}\selectfont\color{popink},
  fontlower=\popsemi\fontsize{10.6}{14.5}\selectfont\color{popink}}}
% Coche dessinée (le glyphe ✓ n'existe pas dans les polices du document)
\newcommand{\popcheck}[1]{\tikz[baseline=-0.5ex]\draw[#1,line width=1.6pt,line cap=round,line join=round](-0.13,0.0)--(-0.04,-0.09)--(0.13,0.1);}
\newcommand{\popboxhead}[3]{\poppill{#1}{#2}{white}\ifx\relax#3\relax\else\hspace{6pt}{\popxbold\fontsize{10.6}{12}\selectfont\color{popink}#3}\fi\par\vspace{7pt}}

\newtcolorbox{aretenir}[1][]{popbase,colframe=popgreen,before upper={\popboxhead{À retenir}{popgreen}{#1}}}
\newtcolorbox{exemplebox}[1][]{popbase,colframe=poporange,before upper={\popboxhead{Exemple}{poporange}{#1}}}
\newtcolorbox{definition}[1][]{popbase,colframe=popblue,before upper={\popboxhead{Définition}{popblue}{#1}}}
\newtcolorbox{propriete}[1][]{popbase,colframe=poppurple,before upper={\popboxhead{Propriété}{poppurple}{#1}}}
\newtcolorbox{methode}[1][]{popbase,colframe=popgold,before upper={\popboxhead{Méthode}{popgold}{#1}}}
\newtcolorbox{attention}[1][]{popbase,colframe=poppink,before upper={\popboxhead{Attention}{poppink}{#1}}}
\newtcolorbox{experience}[1][]{popbase,colframe=popblue,colback=popblueL,before upper={\popboxhead{Expérience}{popblue}{#1}}}
% « Le savais-tu ? » : carte violette pleine (contexte, culture, Cameroun)
\newtcolorbox{savaistu}[1][]{popbase,colback=poppurple,colframe=popink,
  fontupper=\popsemi\fontsize{10.4}{14.2}\selectfont\color{white},
  before upper={\poppill{Le savais-tu\,?}{white}{poppurple}\ifx\relax#1\relax\else\hspace{6pt}{\popxbold\color{white}#1}\fi\par\vspace{7pt}}}
% Exercice résolu : énoncé en haut, \tcblower puis la solution sur fond crème
\newcounter{popexo}\newcounter{popexr}
\newtcolorbox{exoresolu}[1][]{popbase,colframe=poporange,colbacklower=popcream,
  segmentation style={popink,line width=1.2pt,dashed},
  before upper={\stepcounter{popexr}\popboxhead{Exercice résolu \thepopexr}{poporange}{#1}},
  before lower={\poppill{Solution}{popink}{popcream}\par\vspace{6pt}}}
% Exercice (non résolu en place) — la correction est dans la partie Corrigés
\newtcolorbox{exercice}[1][]{popbase,colframe=popink,
  before upper={\stepcounter{popexo}\popboxhead{Exercice \thepopexo}{popink}{#1}}}
% Corrigé
\newtcolorbox{corrige}[1][]{popbase,colframe=popgreen,colback=popgreenL,
  before upper={\popboxhead{Corrigé}{popgreen}{#1}}}
% Objectifs / prérequis / bilan
\newtcolorbox{objectifs}{popbase,colframe=popink,colback=poporangeL,
  before upper={\popboxhead{Objectifs de la leçon}{popink}{}}}
\newtcolorbox{prerequis}{popbase,colframe=popink,colback=popblueL,
  before upper={\popboxhead{Prérequis}{popblue}{}}}
\newtcolorbox{bilan}{popbase,colframe=popink,colback=white,boxrule=2.8pt,
  before upper={\popboxhead{Fiche bilan}{poporange}{L'essentiel à connaître pour l'examen}}}
% Figure encadrée avec légende
\newtcolorbox{popfigure}[1][]{enhanced,breakable=false,colback=white,colframe=popline,boxrule=1.4pt,arc=8pt,
  left=10pt,right=10pt,top=10pt,bottom=8pt,before skip=10pt,after skip=12pt,halign=center,
  lower separated=false,halign lower=center,
  fontlower=\popsemi\fontsize{9}{11.5}\selectfont\color{popmuted},
  #1}
\newcommand{\legende}[1]{\par\vspace{6pt}{\popsemi\fontsize{9}{11.5}\selectfont\color{popmuted}#1}}

% ============================================================
% Signature OnBuch+
% ============================================================
\newcommand{\poplogoinline}[1]{{\poplogo\fontsize{#1}{#1}\selectfont\color{popink}OnBuch\color{poporange}+}}
\newcommand{\signatureonbuch}{%
  \begin{tcolorbox}[enhanced,colback=popink,colframe=popink,boxrule=2.2pt,arc=10pt,
      drop shadow=poporange,left=14pt,right=14pt,top=10pt,bottom=10pt,before skip=0pt,after skip=0pt]
    \tikz[baseline=-0.7ex]\node[circle,fill=popgreen,draw=popcream,line width=1.3pt,minimum size=22pt,inner sep=0]{\popcheck{white}};%
    \hspace{9pt}\begin{minipage}[c]{0.62\linewidth}
      {\popxbold\fontsize{12}{13}\selectfont\color{popcream}Signé\ \textbullet\ {\poplogo OnBuch\color{poporange}+}}\par\vspace{2pt}
      {\raggedright\popsemi\fontsize{8}{10.5}\selectfont\color{popline}\MakeUppercase{Équipe pédagogique OnBuch+}\\\MakeUppercase{Conforme au programme officiel MINESEC}\par}
    \end{minipage}\hfill
    {\poplogo\fontsize{22}{22}\selectfont\color{popcream}OnBuch\color{poporange}+}
  \end{tcolorbox}%
}

% ============================================================
% Page de garde
% #1 titre · #2 matière · #3 niveau · #4 module · #5 n° leçon · #6 durée ·
% #7 description / objectifs (texte ou itemize)
% ============================================================
\newcommand{\pagegarde}[7]{%
  \thispagestyle{empty}
  \newgeometry{a4paper,margin=1.7cm,top=1.6cm,bottom=1.4cm}%
  \begin{tikzpicture}[remember picture,overlay]
    \fill[poporange!18] ([xshift=-3.2cm,yshift=-3.6cm]current page.north east) circle (4.6cm);
    \fill[poppurple!14] ([xshift=2.4cm,yshift=7.6cm]current page.south west) circle (3.8cm);
  \end{tikzpicture}%
  {\poplogo\fontsize{30}{30}\selectfont\color{popink}OnBuch\color{poporange}+}\hfill
  \raisebox{0.6ex}{\poppill{Cours signé \textbullet\ Premium}{popink}{popcream}}\par
  \vspace{1pt}\popsquiggle{poppurple}\par
  \vspace{8pt}
  \begin{tcolorbox}[enhanced,colback=poporange,colframe=popink,boxrule=2.8pt,arc=13pt,
      drop shadow=popink,left=18pt,right=18pt,top=12pt,bottom=13pt,after skip=10pt]
    \poppill{#2 \textbullet\ #3}{popink}{popcream}\hspace{5pt}\poppill{#4}{white}{popink}\par\vspace{8pt}
    {\popdisplay\fontsize{14}{14}\selectfont\color{popink}LEÇON #5}\par\vspace{4pt}
    {\raggedright\hyphenpenalty=10000\exhyphenpenalty=10000\popdisplay\fontsize{25}{27}\selectfont\color{popcream}\MakeUppercase{#1}\par}
    \vspace{8pt}
    {\popxbold\fontsize{9.5}{9.5}\selectfont\color{popink}\MakeUppercase{Durée conseillée : #6}}
  \end{tcolorbox}
  \begin{tcolorbox}[enhanced,colback=white,colframe=popink,boxrule=2.2pt,arc=10pt,
      drop shadow=popink,left=16pt,right=16pt,top=10pt,bottom=10pt,after skip=8pt]
    \poppill{Dans cette leçon}{poppurple}{white}\par\vspace{5pt}
    \popsemi\fontsize{9.3}{12.6}\selectfont\color{popink}#7
  \end{tcolorbox}
  \noindent\begin{minipage}[t]{0.487\linewidth}
    \begin{tcolorbox}[enhanced,colback=popgreen,colframe=popink,boxrule=2.2pt,arc=10pt,
        drop shadow=popink,left=11pt,right=11pt,top=10pt,bottom=10pt,height=3.1cm,valign=center]
      \tcbox[colback=white,colframe=popink,boxrule=1.3pt,arc=5pt,boxsep=0pt,nobeforeafter,
        left=3pt,right=3pt,top=3pt,bottom=3pt,valign=center]{\qrcode[height=1.9cm]{https://play.google.com/store/apps/details?id=com.luvvix.onbuch&hl=fr}}%
      \hspace{8pt}\begin{minipage}[c]{0.5\linewidth}
        {\popdisplay\fontsize{11}{12.5}\selectfont\color{white}\MakeUppercase{Télécharge OnBuch}}\par\vspace{4pt}
        {\raggedright\popsemi\fontsize{8.4}{11}\selectfont\color{white}Gratuit \textbullet\ Google Play\\Tuteur IA, annales, concours, résultats\par}
      \end{minipage}
    \end{tcolorbox}
  \end{minipage}\hfill
  \begin{minipage}[t]{0.487\linewidth}
    \begin{tcolorbox}[enhanced,colback=poppurple,colframe=popink,boxrule=2.2pt,arc=10pt,
        drop shadow=popink,left=11pt,right=11pt,top=10pt,bottom=10pt,height=3.1cm,valign=center]
      \tikz[baseline=-0.7ex]\node[circle,fill=white,draw=popink,line width=1.3pt,minimum size=1.6cm,inner sep=0]{\popdisplay\fontsize{24}{24}\selectfont\color{poppurple}+};%
      \hspace{8pt}\begin{minipage}[c]{0.6\linewidth}
        {\popdisplay\fontsize{11}{12.5}\selectfont\color{white}\MakeUppercase{Passe à OnBuch+}}\par\vspace{4pt}
        {\raggedright\popsemi\fontsize{8.4}{11}\selectfont\color{white}Tous les cours signés, Tuteur IA illimité, fiches PDF — onglet OnBuch+ de l'app\par}
      \end{minipage}
    \end{tcolorbox}
  \end{minipage}
  \vspace{8pt}
  \popcontenu
  \vfill
  \signatureonbuch
  \restoregeometry
  \newpage
}

% Rangée « Ce cours contient » (page de garde)
\newcommand{\poptile}[3]{% #1 couleur #2 gros texte #3 légende
  \begin{tcolorbox}[enhanced,colback=#1,colframe=popink,boxrule=2pt,arc=9pt,shadow={2.5pt}{-2.5pt}{0pt}{popink},
      left=6pt,right=6pt,top=9pt,bottom=9pt,halign=center,valign=center,height=1.75cm,width=\linewidth,nobeforeafter]
    {\popdisplay\fontsize{12.5}{13}\selectfont\color{white}\MakeUppercase{#2}}\par\vspace{3pt}
    {\centering\popsemi\fontsize{7.8}{9.5}\selectfont\color{white}#3\par}
  \end{tcolorbox}}
\newcommand{\popcontenu}{%
  \par\noindent{\popxboldls\fontsize{7.5}{7.5}\selectfont\color{popmuted}CE COURS SIGNÉ CONTIENT}\par\vspace{7pt}
  \noindent\begin{minipage}[t]{0.235\linewidth}\poptile{popblue}{Cours}{complet, illustré, pas à pas}\end{minipage}\hfill
  \begin{minipage}[t]{0.235\linewidth}\poptile{poporange}{Méthodes}{et exercices résolus}\end{minipage}\hfill
  \begin{minipage}[t]{0.235\linewidth}\poptile{poppink}{Exercices}{type examen + corrigés}\end{minipage}\hfill
  \begin{minipage}[t]{0.235\linewidth}\poptile{popgreen}{Fiche bilan}{l'essentiel à retenir}\end{minipage}\par}

% Sommaire
\newtcolorbox{sommaire}{enhanced,colback=white,colframe=popink,boxrule=2.2pt,arc=10pt,
  drop shadow=popink,left=18pt,right=18pt,top=13pt,bottom=13pt,before skip=6pt,after skip=20pt,
  before upper={\poppill{Sommaire}{poppurple}{white}\par\vspace{9pt}\popsemi\fontsize{10.6}{14.5}\selectfont\color{popink}}}

% Signature de fin de cours — poussée tout en bas de la dernière page
\newcommand{\findecours}{%
  \par\vfill\nopagebreak
  \begin{center}\popsquiggle{poporange}\end{center}\vspace{4pt}
  \signatureonbuch
}

%% FIGFIX-BEGIN v4 — correctifs globaux des figures (docs/onbuch-generation/tools/figfix)
\usepackage{adjustbox}\usepackage{etoolbox}
% toute figure TikZ/pgfplots plus large que la ligne est réduite à la largeur disponible
\BeforeBeginEnvironment{tikzpicture}{\begin{adjustbox}{max width=\linewidth}}
\AfterEndEnvironment{tikzpicture}{\end{adjustbox}}
% légendes pgfplots lisibles par-dessus les courbes
\pgfplotsset{every axis/.append style={legend style={fill=white,fill opacity=0.92,text opacity=1,draw=black!15,inner sep=3pt}}}
% étiquettes d'arêtes (syntaxe edge["..."]) avec fond blanc
\tikzset{every edge quotes/.append style={fill=white,inner sep=1.2pt}}
%% FIGFIX-END
\begin{document}

\pagegarde{Propriétés chimiques des alcools}{Chimie}{Tle E}{Module 1 · Chimie organique}{01}{4 h (cours + TP)}{Cette leçon étudie les grandes réactions des alcools : action du sodium, déshydratation, oxydations (combustion, ménagée, catalytique) et estérification. Elle présente aussi les deux voies de préparation des alcools — hydratation des alcènes et fermentation alcoolique — et s'appuie sur un TP d'oxydation ménagée.
\vspace{4pt}
\begin{multicols}{2}\small
\begin{itemize}
\item À la fin de cette leçon, je sais écrire l'équation-bilan de la réaction d'un alcool avec le sodium et interpréter la mobilité de l'hydrogène du groupe hydroxyle
\item À la fin de cette leçon, je sais distinguer déshydratation intramoléculaire et intermoléculaire et écrire leurs équations-bilans avec les conditions expérimentales
\item À la fin de cette leçon, je sais écrire l'équation de combustion complète d'un alcool
\item À la fin de cette leçon, je sais prévoir les produits de l'oxydation ménagée d'un alcool selon sa classe (primaire, secondaire, tertiaire) avec KMnO4 ou K2Cr2O7 en milieu acide
\item À la fin de cette leçon, je sais écrire les demi-équations électroniques et l'équation-bilan d'une oxydation ménagée d'alcool ou d'aldéhyde
\item À la fin de cette leçon, je sais décrire l'expérience de la lampe sans flamme et identifier les produits par les tests caractéristiques (DNPH, Fehling, Schiff, Tollens)
\end{itemize}\end{multicols}}

\begin{sommaire}
\begin{multicols}{2}
\begin{enumerate}
\item Rappels et réaction avec le sodium
\item Déshydratation des alcools
\item Oxydation des alcools : combustion et oxydation ménagée
\item Lampe sans flamme et identification des produits
\item Estérification
\item Préparation des alcools
\item Travaux pratiques
\item Méthodes et exercices résolus
\item Exercices
\item Corrigés des exercices
\item Fiche bilan
\end{enumerate}
\end{multicols}
\end{sommaire}

\begin{objectifs}
\begin{itemize}
\item À la fin de cette leçon, je sais écrire l'équation-bilan de la réaction d'un alcool avec le sodium et interpréter la mobilité de l'hydrogène du groupe hydroxyle
\item À la fin de cette leçon, je sais distinguer déshydratation intramoléculaire et intermoléculaire et écrire leurs équations-bilans avec les conditions expérimentales
\item À la fin de cette leçon, je sais écrire l'équation de combustion complète d'un alcool
\item À la fin de cette leçon, je sais prévoir les produits de l'oxydation ménagée d'un alcool selon sa classe (primaire, secondaire, tertiaire) avec KMnO4 ou K2Cr2O7 en milieu acide
\item À la fin de cette leçon, je sais écrire les demi-équations électroniques et l'équation-bilan d'une oxydation ménagée d'alcool ou d'aldéhyde
\item À la fin de cette leçon, je sais décrire l'expérience de la lampe sans flamme et identifier les produits par les tests caractéristiques (DNPH, Fehling, Schiff, Tollens)
\item À la fin de cette leçon, je sais écrire l'équation d'estérification, citer ses caractéristiques et calculer un rendement
\item À la fin de cette leçon, je sais écrire les équations d'hydratation d'un alcène (règle de Markovnikov) et de fermentation alcoolique du glucose
\end{itemize}
\end{objectifs}

\begin{prerequis}
\begin{itemize}
\item Formules brutes, semi-développées et nomenclature des alcools ; classe d'un alcool (primaire, secondaire, tertiaire)
\item Couples oxydant/réducteur et demi-équations électroniques en milieu acide (Première)
\item Réactions acido-basiques : acide carboxylique + base (notion d'acide carboxylique)
\item Structure et réactivité des alcènes (addition électrophile, début d'année de Terminale)
\item Équilibrage d'équations-bilans et calculs de quantités de matière
\end{itemize}
\end{prerequis}

\newpage

\coursec{Rappels et réaction avec le sodium}

\courssub{Situation d'accroche : de Bafoussam à Douala, l'alcool est partout}

À Bafoussam comme à Douala, l'\emph{odontol} et les boissons alcoolisées artisanales sont obtenus par fermentation du maïs ou de la canne à sucre : comment le simple glucose se transforme-t-il en éthanol ? Pourquoi l'alcool à brûler vendu en pharmacie rougit-il un fil de cuivre chauffé, en dégageant une odeur piquante qui pique les yeux ? Pourquoi le vin de palme exposé à l'air se transforme-t-il, en quelques jours, en vinaigre ? Et comment les savonneries artisanales parfument-elles leurs savons avec des esters aux odeurs fruitées d'ananas ou de banane ?

Toutes ces questions, en apparence très différentes, trouvent leur réponse dans un même chapitre : les \cle{propriétés chimiques des alcools}. Au cours de cette leçon, tu vas découvrir que le petit groupe \ce{-OH}, en apparence modeste, cache un hydrogène \emph{mobile}, capable de réagir avec le sodium ; qu'il peut être arraché avec un hydrogène voisin (déshydratation) ; que le carbone qui le porte peut être oxydé de façon contrôlée ; et qu'il peut s'associer à un acide carboxylique pour fabriquer les esters des parfums.

\textbf{Problématique :} quelles transformations chimiques le groupe hydroxyle \ce{-OH} des alcools permet-il de réaliser, et comment les exploiter pour préparer ou identifier des composés oxygénés ?

Commençons par le commencement : rappeler ce qu'est un alcool, puis mettre en évidence expérimentalement la mobilité de l'hydrogène de son groupe fonctionnel.

\courssub{Rappels : structure, classe et nomenclature des alcools}

\begin{definition}[Alcool — groupe hydroxyle]
Un \cle{alcool} est un composé organique oxygéné dont la molécule possède un \cle{groupe hydroxyle} \ce{-OH} porté par un atome de carbone \emph{tétragonal} (hybridé $sp^3$, c'est-à-dire lié à quatre atomes par des liaisons simples). Sa formule générale s'écrit \ce{R-OH}, où R est un groupe alkyle (\ce{C_nH_{2n+1}} pour un alcool saturé acyclique, de formule brute \ce{C_nH_{2n+2}O}).
\end{definition}

\begin{attention}[Alcool ou phénol ?]
Si le groupe \ce{-OH} est porté directement par un carbone du cycle benzénique, le composé est un \emph{phénol}, pas un alcool : ses propriétés sont différentes. Au programme de Terminale, le carbone fonctionnel d'un alcool est toujours tétragonal.
\end{attention}

La \cle{classe} d'un alcool se détermine en comptant le nombre de groupes alkyles (ou d'atomes de carbone) directement liés au \emph{carbone fonctionnel}, c'est-à-dire le carbone qui porte le \ce{-OH} :

\begin{itemize}
\item \textbf{Alcool primaire} : le carbone fonctionnel est lié à \textbf{un seul} atome de carbone (ou à aucun pour le méthanol). Exemple : l'éthanol \chemfig{CH_3-CH_2-OH}.
\item \textbf{Alcool secondaire} : le carbone fonctionnel est lié à \textbf{deux} atomes de carbone. Exemple : le propan-2-ol \chemfig{H_3C-CH(-[2]OH)-CH_3}.
\item \textbf{Alcool tertiaire} : le carbone fonctionnel est lié à \textbf{trois} atomes de carbone. Exemple : le 2-méthylpropan-2-ol \chemfig{H_3C-C(-[2]OH)(-[6]CH_3)-CH_3}.
\end{itemize}

\begin{exemplebox}[Nommer rapidement un alcool]
Pour la nomenclature, on choisit la chaîne la plus longue contenant le carbone fonctionnel, on la numérote pour donner à ce carbone le plus petit indice possible, et on remplace le « e » final de l'alcane par le suffixe « -ol » précédé de l'indice de position. Ainsi \chemfig{CH_3-CH_2-CH(-[2]OH)-CH_3} est le \textbf{butan-2-ol} (alcool secondaire), et non « butan-3-ol » : on numérote toujours du côté qui donne l'indice le plus petit au \ce{-OH}.
\end{exemplebox}

\begin{aretenir}[Pourquoi la classe est-elle si importante ?]
La classe d'un alcool détermine son comportement en \emph{oxydation ménagée} (section 3) : un alcool primaire conduit à un aldéhyde puis à un acide carboxylique, un alcool secondaire à une cétone, et un alcool tertiaire ne subit pas d'oxydation ménagée. Savoir reconnaître la classe d'un alcool est donc un réflexe indispensable.
\end{aretenir}

\courssub{Polarisation de la liaison O--H : un hydrogène mobile}

Pourquoi l'hydrogène du groupe \ce{-OH} serait-il plus « réactif » que les hydrogènes de la chaîne carbonée ? L'explication tient en un mot : l'\cle{électronégativité}.

L'atome d'oxygène est beaucoup plus électronégatif que l'hydrogène ($\chi_O = 3{,}5$ contre $\chi_H = 2{,}2$ sur l'échelle de Pauling). Il attire donc fortement vers lui le doublet d'électrons de la liaison \ce{O-H} : la liaison est \emph{polarisée}. L'oxygène porte une charge partielle négative $\delta^-$ et l'hydrogène une charge partielle positive $\delta^+$ :

\[ \chemfig{C-[1]O^{\delta^-}-[-1]H^{\delta^+}} \]

Cet hydrogène, partiellement « dénudé », peut être arraché par un réactif suffisamment énergique : on dit qu'il est \cle{mobile}. L'alcool possède ainsi un (très faible) \emph{caractère acide} au sens de Brønsted : il peut céder un proton \ce{H+} pour donner l'ion \cle{alcoolate} \ce{R-O^-}.

En revanche, les liaisons \ce{C-H} de la chaîne carbonée sont très peu polarisées (le carbone et l'hydrogène ont des électronégativités proches) : leurs hydrogènes ne sont \emph{pas} mobiles. C'est toute la différence, et c'est ce que l'expérience avec le sodium va confirmer de manière quantitative.

\courssub{Expérience : action du sodium sur l'éthanol}

\begin{experience}[Le sodium réagit avec l'éthanol]
\textbf{Dispositif et protocole.} Dans un tube à essai sec, on verse environ \SI{5}{mL} d'éthanol absolu (éthanol anhydre). On y introduit, à l'aide d'une pince, un petit morceau de sodium de la taille d'un grain de riz, fraîchement sorti de l'huile (pétrole lampant) et essuyé. On recueille le gaz dégagé dans un tube retourné par déplacement, puis on présente une flamme à l'ouverture du tube.

\textbf{Observations.}
\begin{itemize}
\item Le sodium réagit avec un dégagement gazeux \emph{modéré} (effervescence douce) : il reste au fond du tube, fond lentement et disparaît progressivement.
\item Le gaz recueilli produit une \emph{détonation} caractéristique (« aboiement ») au contact de la flamme : c'est du \cle{dihydrogène} \ce{H2}.
\item Après évaporation de l'excès d'alcool, on obtient un solide blanc, l'\cle{éthanolate de sodium} \ce{C2H5O^-Na+}, qui, dissous dans l'eau, donne une solution fortement basique (le rouge de phénolphtaléine vire au rose foncé).
\end{itemize}

\textbf{Interprétation.} Le sodium, métal très réducteur, arrache l'hydrogène \emph{du groupe hydroxyle} : il se forme l'ion éthanolate \ce{C2H5O^-} et du dihydrogène. La réaction est nettement \emph{moins vive} qu'avec l'eau : avec l'eau, le sodium fond instantanément en une bille qui fuse à la surface, pouvant s'enflammer ; avec l'éthanol, tout se passe calmement.
\end{experience}

\begin{popfigure}
\begin{center}
\begin{tikzpicture}[scale=0.95, every node/.style={font=\small}]
% Tube à essai
\draw[thick, popdark] (0,0) -- (0,3.4);
\draw[thick, popdark] (1.4,0) -- (1.4,3.4);
\draw[thick, popdark] (0,0) arc (180:360:0.7);
% Liquide éthanol
\fill[popblueL] (0.03,-0.62) arc (180:360:0.67) -- (1.37,1.6) -- (0.03,1.6) -- cycle;
\draw[popblue, thick] (0.03,1.6) -- (1.37,1.6);
% Morceau de sodium
\fill[popmuted!60] (0.55,-0.35) rectangle (0.9,0.05);
\draw[popdark] (0.55,-0.35) rectangle (0.9,0.05);
\node[popdark, font=\footnotesize] at (2.6,-0.1) {sodium \ce{Na}};
\draw[-{Stealth}, popdark] (1.85,-0.12) -- (0.95,-0.12);
% Bulles
\foreach \x/\y/\r in {0.6/0.5/0.05, 0.85/0.9/0.07, 0.55/1.2/0.05, 0.95/1.35/0.06, 0.75/0.7/0.04}
  \draw[popblue] (\x,\y) circle (\r);
% Étiquette éthanol
\node[popdark, font=\footnotesize, align=left] at (-2.2,0.9) {éthanol absolu\\ \ce{C2H5OH}};
\draw[-{Stealth}, popdark] (-1.2,0.9) -- (-0.05,0.9);
% Tube de recueil du gaz
\draw[thick, popdark] (3.6,3.2) -- (3.6,5.6);
\draw[thick, popdark] (5.0,3.2) -- (5.0,5.6);
\draw[thick, popdark] (3.6,5.6) arc (180:0:0.7);
\fill[popgreenL] (3.63,4.4) -- (4.97,4.4) -- (4.97,5.55) arc (0:180:0.67) -- cycle;
\node[popdark, font=\footnotesize] at (4.3,5.0) {\ce{H2}};
% Tuyau de jonction
\draw[thick, popdark] (0.7,3.4) -- (0.7,4.1) -- (4.3,4.1) -- (4.3,3.2);
% Flamme
\draw[fill=poporange, poporange] (6.3,3.1) .. controls (6.05,3.7) and (6.15,4.0) .. (6.3,4.35)
  .. controls (6.45,4.0) and (6.55,3.7) .. (6.3,3.1);
\draw[fill=popgold, popgold] (6.3,3.15) .. controls (6.18,3.55) and (6.24,3.75) .. (6.3,3.95)
  .. controls (6.36,3.75) and (6.42,3.55) .. (6.3,3.15);
\draw[thick, popdark] (6.3,3.1) -- (6.3,2.5);
\node[popdark, font=\footnotesize, align=center] at (6.3,2.15) {flamme :\\ détonation};
\draw[-{Stealth}, poporange] (5.9,3.9) -- (5.05,4.6);
\end{tikzpicture}
\end{center}
\legende{Action du sodium sur l'éthanol : dégagement de dihydrogène, recueilli dans un tube retourné et identifié par la détonation au contact d'une flamme.}
\end{popfigure}

\courssub{Équation-bilan et généralisation}

L'équation-bilan de la réaction entre l'éthanol et le sodium s'écrit :

\[ \ce{2 C2H5OH + 2 Na -> 2 C2H5O^-Na+ + H2} \]

Commentons-la. Deux moles d'éthanol réagissent avec deux moles de sodium pour donner deux moles d'éthanolate de sodium et \emph{une} mole de dihydrogène. Chaque atome de sodium cède un électron (demi-équation d'oxydation : \ce{Na -> Na+ + e-}), et chaque hydrogène du groupe \ce{-OH} capte un électron (demi-équation de réduction : \ce{2 C2H5OH + 2 e- -> 2 C2H5O^- + H2}). C'est donc une réaction d'\emph{oxydoréduction}.

\begin{popfigure}
\begin{center}
\begin{tikzpicture}[every node/.style={font=\small}]
\node at (0,0) {\chemfig{CH_3-CH_2-O-H}};
\draw[-{Stealth}, very thick, poppink] (1.15,0.55) .. controls (1.9,1.3) and (2.9,1.3) .. (3.6,0.55);
\node[poppink, font=\footnotesize, align=center] at (2.4,1.45) {départ de \ce{H+}\\ (H mobile)};
\node at (5.6,0) {\chemfig{CH_3-CH_2-O^{-}}};
\node[popdark] at (7.15,0) {\ce{Na+}};
\node[popdark, font=\footnotesize, align=center] at (6.4,-0.9) {ion éthanolate\\ (base forte)};
\end{tikzpicture}
\end{center}
\legende{L'arrachage de l'hydrogène mobile du groupe hydroxyle de l'éthanol conduit à l'ion éthanolate \ce{C2H5O^-}, base conjuguée de l'alcool.}
\end{popfigure}

\begin{definition}[Ion alcoolate]
L'\cle{ion alcoolate} \ce{R-O^-} est la \emph{base conjuguée} de l'alcool \ce{R-OH}, obtenue par départ du proton du groupe hydroxyle. C'est une \textbf{base forte} : il réagit violemment et totalement avec l'eau selon \ce{R-O^- + H2O -> R-OH + HO^-}. Les alcoolates de sodium (méthanolate \ce{CH3O^-Na+}, éthanolate \ce{C2H5O^-Na+}...) sont des solides ioniques blancs, utilisés en synthèse organique comme bases et comme nucléophiles.
\end{definition}

La réaction se généralise à \emph{tous} les alcools, quelle que soit leur classe :

\[ \ce{2 R-OH + 2 Na -> 2 R-O^-Na+ + H2} \]

\begin{methode}[Équilibrer « alcool + sodium » sans se tromper]
\begin{enumerate}
\item Écris les réactifs : \ce{R-OH + Na}.
\item Souviens-toi que \emph{deux} hydrogènes mobiles sont nécessaires pour former \emph{une} molécule \ce{H2} : il faut donc \textbf{2 molécules d'alcool}.
\item Place le coefficient 2 devant l'alcool, devant le sodium et devant l'alcoolate ; le coefficient de \ce{H2} est 1.
\item Vérifie : \ce{2 R-OH + 2 Na -> 2 R-O^-Na+ + H2}. Bilan des atomes : 2 \ce{Na} de chaque côté, 2 \ce{O}, et les hydrogènes des groupes \ce{-OH} se retrouvent dans \ce{H2}.
\end{enumerate}
Astuce de mémorisation : \textbf{« deux alcools pour un dihydrogène »}.
\end{methode}

\begin{attention}[Les deux erreurs classiques]
\begin{itemize}
\item \textbf{Erreur n\textsuperscript{o} 1 : écrire \ce{H2O} comme produit.} Il ne se forme \emph{jamais} d'eau dans cette réaction : le produit gazeux est le dihydrogène \ce{H2}, et le produit solide est l'alcoolate \ce{R-O^-Na+} (et non l'hydroxyde \ce{NaOH} !).
\item \textbf{Erreur n\textsuperscript{o} 2 : faire réagir les hydrogènes de la chaîne carbonée.} Seul l'hydrogène du groupe \ce{-OH} est mobile, car seule la liaison \ce{O-H} est fortement polarisée. Les hydrogènes portés par les carbones ne participent pas. Conséquence quantitative capitale : \textbf{1 mol d'alcool dégage exactement $\dfrac{1}{2}$ mol de \ce{H2}}, quel que soit le nombre d'atomes d'hydrogène de la molécule.
\end{itemize}
\end{attention}

\courssub{Comparaison avec l'eau : l'alcool est un acide plus faible}

L'eau possède, elle aussi, une liaison \ce{O-H} polarisée, et réagit avec le sodium :

\[ \ce{2 H2O + 2 Na -> 2 Na+ + 2 HO^- + H2} \]

Mais l'expérience montre que la réaction est \emph{beaucoup plus vive} avec l'eau qu'avec l'éthanol. Pourquoi ? Le groupe alkyle R est \emph{donneur d'électrons} (effet inductif donneur) : il « pousse » de la densité électronique vers l'oxygène, ce qui diminue la polarisation de la liaison \ce{O-H} et rend l'hydrogène moins mobile. L'alcool est donc un \textbf{acide plus faible que l'eau} ; réciproquement, l'ion alcoolate \ce{R-O^-} est une \textbf{base plus forte} que l'ion hydroxyde \ce{HO^-}, ce qui explique pourquoi il est instantanément détruit par l'eau.

\begin{center}
\begin{tabularx}{\linewidth}{|l|X|X|}
\hline
\rowcolor{popblueL} \textbf{Critère} & \textbf{Eau \ce{H2O}} & \textbf{Éthanol \ce{C2H5OH}} \\
\hline
Hydrogène mobile & oui (2 par molécule) & oui (1 seul : celui du \ce{-OH}) \\
\hline
Vitesse avec \ce{Na} & très vive, bille de sodium qui fuse & modérée, effervescence douce \\
\hline
Force acide & plus acide & acide plus faible \\
\hline
Base conjuguée & \ce{HO^-} (base forte) & \ce{C2H5O^-} (base encore plus forte) \\
\hline
\end{tabularx}
\end{center}

\courssub{Mise en évidence quantitative de la mobilité de l'hydrogène}

La stœchiométrie de la réaction fournit un test élégant : en mesurant le volume de dihydrogène dégagé par une masse connue d'alcool, on vérifie qu'\emph{un seul} hydrogène par molécule est mobile. C'est historiquement l'une des preuves de la structure \ce{R-OH} des alcools.

\begin{exoresolu}[Combien de dihydrogène dégage l'éthanol ?]
\textbf{Énoncé.} On fait réagir un excès de sodium avec $m = \SI{4,6}{g}$ d'éthanol pur \ce{C2H5OH}. (1) Écris l'équation-bilan de la réaction. (2) Calcule la quantité de matière d'éthanol utilisée. (3) Déduis-en le volume de dihydrogène dégagé, mesuré dans les conditions où le volume molaire vaut $V_m = \SI{24,0}{L.mol^{-1}}$. (4) L'éthanol contient 6 atomes d'hydrogène par molécule : pourquoi n'obtient-on pas 3 fois plus de \ce{H2} ?

\emph{Donnée :} $M(\ce{C2H5OH}) = \SI{46,0}{g.mol^{-1}}$.
\tcblower
\textbf{Solution détaillée.}

(1) Équation-bilan :
\[ \ce{2 C2H5OH + 2 Na -> 2 C2H5O^-Na+ + H2} \]

(2) Quantité de matière d'éthanol :
\[ n(\ce{C2H5OH}) = \frac{m}{M} = \frac{4{,}6}{46{,}0} = \SI{0,10}{mol} \]

(3) Dressons le tableau d'avancement (le sodium est en excès, l'éthanol est le réactif limitant) :

\begin{center}
\begin{tabularx}{\linewidth}{|l|X|X|X|X|}
\hline
\rowcolor{popblueL} \textbf{État} & \ce{2 C2H5OH} & \ce{2 Na} & \ce{2 C2H5O^-Na+} & \ce{H2} \\
\hline
Initial (mol) & 0,10 & excès & 0 & 0 \\
\hline
Final (mol) & $0{,}10 - 2x_{max} = 0$ & excès & $2x_{max}$ & $x_{max}$ \\
\hline
\end{tabularx}
\end{center}

De $0{,}10 - 2x_{max} = 0$, on tire $x_{max} = \SI{0,050}{mol}$. D'où :
\[ n(\ce{H2}) = x_{max} = \SI{0,050}{mol} \qquad V(\ce{H2}) = n \times V_m = 0{,}050 \times 24{,}0 = \SI{1,2}{L} \]

(4) Si les 6 hydrogènes de la molécule étaient mobiles, on obtiendrait $3 \times 0{,}10 = \SI{0,30}{mol}$ de \ce{H2}, soit \SI{7,2}{L}. Or on ne recueille que \SI{1,2}{L} : cela prouve qu'\textbf{un seul hydrogène par molécule} — celui du groupe \ce{-OH} — est mobile. Les cinq hydrogènes liés aux carbones ne réagissent pas.
\end{exoresolu}

\begin{savaistu}[Pourquoi le sodium dort-il dans l'huile ?]
Le sodium est un métal si réducteur qu'il réagit violemment avec l'eau \emph{et} avec les alcools, en dégageant du dihydrogène inflammable : un vrai risque d'incendie ou d'explosion. C'est pourquoi, au laboratoire comme dans l'industrie, on le conserve \emph{immergé dans l'huile (pétrole lampant)}, un solvant hydrocarboné dont les molécules ne possèdent \emph{aucune} liaison \ce{O-H} : aucun hydrogène mobile, donc aucune réaction. Avant d'utiliser un morceau de sodium, on l'essuie soigneusement avec du papier filtre, et on ne le manipule jamais à mains nues — l'humidité de la peau suffirait à déclencher la réaction et à provoquer des brûlures. Et surtout : jamais de sodium à proximité d'une bouteille d'alcool à brûler !
\end{savaistu}

\begin{aretenir}[L'essentiel de la section]
\begin{itemize}
\item Un alcool \ce{R-OH} possède un groupe hydroxyle dont la liaison \ce{O-H} est polarisée : l'hydrogène fonctionnel est \cle{mobile} (caractère faiblement acide).
\item Le sodium arrache cet hydrogène : \ce{2 R-OH + 2 Na -> 2 R-O^-Na+ + H2}. Le gaz dégagé, \ce{H2}, est identifié par sa détonation à la flamme.
\item L'ion \cle{alcoolate} \ce{R-O^-} formé est une base forte, détruite par l'eau.
\item La réaction est moins vive qu'avec l'eau : l'alcool est un acide plus faible que l'eau.
\item Point stœchiométrique à retenir : \textbf{2 mol d'alcool $\rightarrow$ 1 mol de \ce{H2}}, car seul le H du \ce{-OH} réagit.
\end{itemize}
\end{aretenir}

\begin{exercice}[À toi de jouer]
On introduit un morceau de sodium en excès dans un échantillon de propan-1-ol pur. On recueille \SI{600}{mL} de dihydrogène dans des conditions où $V_m = \SI{24,0}{L.mol^{-1}}$. (1) Écris l'équation-bilan de la réaction et nomme le produit organique formé. (2) Calcule la masse de propan-1-ol ayant réagi. \emph{Donnée :} $M(\ce{C3H7OH}) = \SI{60,0}{g.mol^{-1}}$.
\end{exercice}

\begin{corrige}[Exercice 1]
(1) \ce{2 C3H7OH + 2 Na -> 2 C3H7O^-Na+ + H2} ; le produit organique est le \textbf{propan-1-olate de sodium} (ou propanolate de sodium).

(2) $n(\ce{H2}) = \dfrac{V}{V_m} = \dfrac{0{,}600}{24{,}0} = \SI{0,025}{mol}$. D'après la stœchiométrie, $n(\text{alcool}) = 2\,n(\ce{H2}) = \SI{0,050}{mol}$. D'où $m = n \times M = 0{,}050 \times 60{,}0 = \SI{3,0}{g}$ de propan-1-ol.
\end{corrige}

\coursec{Déshydratation des alcools}

Dans la section précédente, nous avons vu que l'hydrogène du groupe hydroxyle \ce{-OH} est \cle{mobile} : il peut être arraché par le sodium pour donner un alcoolate. Mais le groupe hydroxyle peut aussi quitter la molécule \emph{entièrement}, emportant avec lui un atome d'hydrogène voisin, sous forme d'une molécule d'eau. Cette transformation, appelée \cle{déshydratation}, est l'une des réactions les plus importantes des alcools : elle permet de fabriquer des alcènes (matières premières des plastiques) ou des étheroxydes (solvants). Comprendre comment une simple variation de température oriente la réaction vers l'un ou l'autre produit est un classique du Baccalauréat.

\courssub{Généralités sur la déshydratation}

\begin{definition}[Déshydratation d'un alcool]
La \cle{déshydratation} d'un alcool est une réaction d'\cle{élimination} au cours de laquelle une molécule d'eau \ce{H2O} est arrachée à partir du (ou des) groupe(s) hydroxyle. Elle nécessite :
\begin{itemize}
\item un \cle{déshydratant} : acide sulfurique concentré \ce{H2SO4}, acide phosphorique concentré \ce{H3PO4}, ou passage des vapeurs d'alcool sur de l'\cle{alumine} \ce{Al2O3} chauffée vers \SI{400}{\celsius} ;
\item un \cle{chauffage} dont la température détermine la nature du produit obtenu.
\end{itemize}
Selon les conditions, la déshydratation est \textbf{intramoléculaire} (une seule molécule d'alcool perd \ce{H2O} : on obtient un \textbf{alcène}) ou \textbf{intermoléculaire} (deux molécules d'alcool perdent ensemble \ce{H2O} : on obtient un \textbf{étheroxyde}).
\end{definition}

\textbf{Intuition.} Imagine le groupe \ce{-OH} comme une « poignée » par laquelle l'acide sulfurique (excellent capteur d'eau) va tirer. Si la molécule d'alcool est seule et très chaude, elle perd \ce{OH} plus un \ce{H} pris sur le carbone voisin : il reste une double liaison, c'est un alcène. Si au contraire deux molécules d'alcool sont proches et que la température est modérée, l'une fournit \ce{OH} et l'autre fournit \ce{H} : les deux squelettes carbonés restent soudés par un atome d'oxygène, c'est un étheroxyde.

\begin{attention}[Conditions opératoires obligatoires]
Au Baccalauréat, une équation de déshydratation \textbf{sans les conditions portées sur la flèche} (nature de l'acide, température) est considérée comme incomplète et coûte des points. Écris systématiquement : \ce{H2SO4} concentré (ou \ce{H3PO4}) et la température au-dessus de la flèche.
\end{attention}

\courssub{Déshydratation intramoléculaire : obtention d'un alcène}

\begin{definition}[Déshydratation intramoléculaire]
La déshydratation \cle{intramoléculaire} est l'élimination d'une molécule d'eau \emph{au sein d'une seule molécule d'alcool} : le groupe \ce{-OH} part avec un atome d'hydrogène porté par le carbone \emph{voisin} (carbone en $\alpha$). Il se forme une \textbf{liaison double \ce{C=C}} : le produit est un \cle{alcène}. Conditions : \ce{H2SO4} concentré (ou \ce{H3PO4}), température élevée, environ \SI{180}{\celsius}, avec \textbf{excès d'acide}.
\end{definition}

\textbf{Exemple fondamental : l'éthanol.} Chauffé vers \SI{180}{\celsius} en présence d'acide sulfurique concentré, l'éthanol donne l'éthène :

\[ \ce{CH3-CH2-OH} \xrightarrow[\text{180 °C}]{\text{H2SO4 conc.}} \ce{CH2=CH2 + H2O} \]

On vérifie le bilan : à gauche \ce{C2H6O}, à droite \ce{C2H4} + \ce{H2O} = \ce{C2H6O}. L'équation est équilibrée. L'éthène se dégage sous forme gazeuse ; on peut le recueillir et l'identifier par sa décoloration de l'eau de brome.

\textbf{Écriture en formules semi-développées :}

\begin{center}
\chemfig{CH_3-CH_2-OH} \quad $\xrightarrow[\text{180 °C}]{\text{H2SO4 conc.}}$ \quad \chemfig{CH_2=CH_2} \quad $+$ \quad \ce{H2O}
\end{center}

\textbf{Cas des alcools dissymétriques : règle de Zaïtsev.} Lorsque le carbone porteur du \ce{-OH} possède deux voisins différents, plusieurs alcènes peuvent se former. Exemple : le butan-2-ol \chemfig{CH_3-CH(-[2]OH)-CH_2-CH_3} peut perdre l'hydrogène du carbone 1 (donnant le but-1-ène) ou celui du carbone 3 (donnant le but-2-ène).

\begin{propriete}[Règle de Zaïtsev]
Lors de la déshydratation d'un alcool dissymétrique, l'alcène \textbf{majoritaire} est l'alcène le plus \cle{substitué}, c'est-à-dire celui dont la double liaison porte le \textbf{plus grand nombre de groupes alkyles}. Pour le butan-2-ol, le but-2-ène \chemfig{CH_3-CH=CH-CH_3} (double liaison disubstituée) est majoritaire devant le but-1-ène \chemfig{CH_2=CH-CH_2-CH_3} (monosubstitué).
\end{propriete}

\courssub{Déshydratation intermoléculaire : obtention d'un étheroxyde}

\begin{definition}[Déshydratation intermoléculaire]
La déshydratation \cle{intermoléculaire} est l'élimination d'une molécule d'eau \emph{entre deux molécules d'alcool} : l'une fournit le groupe \ce{-OH}, l'autre fournit l'hydrogène de son propre groupe hydroxyle. Les deux groupes alkyles restent reliés par un atome d'oxygène : le produit est un \cle{étheroxyde} de formule générale $R-O-R'$. Conditions : \ce{H2SO4} concentré, température modérée, environ \SI{140}{\celsius}, avec \textbf{excès d'alcool}.
\end{definition}

\textbf{Exemple fondamental : l'éthoxyéthane (éther diéthylique).}

\[ \ce{2 C2H5OH} \xrightarrow[\text{140 °C}]{\text{H2SO4 conc.}} \ce{C2H5-O-C2H5 + H2O} \]

soit, en formules semi-développées :

\begin{center}
\chemfig{CH_3-CH_2-OH} \; $+$ \; \chemfig{HO-CH_2-CH_3} \quad $\xrightarrow[\text{140 °C}]{\text{H2SO4 conc.}}$ \quad \chemfig{CH_3-CH_2-O-CH_2-CH_3} \; $+$ \; \ce{H2O}
\end{center}

\begin{attention}[Ne pas confondre les deux bilans]
Déshydratation intramoléculaire : \textbf{une} molécule d'alcool $\to$ alcène + eau (coefficient 1). Déshydratation intermoléculaire : \textbf{deux} molécules d'alcool $\to$ étheroxyde + eau (coefficient 2 devant l'alcool). Oublier ce coefficient 2 est l'erreur la plus fréquente ; elle fausse ensuite tous les calculs de quantités de matière.
\end{attention}

\courssub{Le rôle décisif de la température}

La température est le \cle{bouton de réglage} qui oriente la déshydratation :

\begin{center}
\begin{tabularx}{\linewidth}{|l|X|X|}
\hline
\rowcolor{popblueL} \textbf{Conditions} & \textbf{Produit majoritaire} & \textbf{Type de réaction} \\
\hline
\ce{H2SO4} conc., $\approx$ \SI{180}{\celsius}, excès d'acide & Alcène (ex. éthène) & Intramoléculaire \\
\hline
\ce{H2SO4} conc., $\approx$ \SI{140}{\celsius}, excès d'alcool & Étheroxyde (ex. éthoxyéthane) & Intermoléculaire \\
\hline
Vapeurs sur \ce{Al2O3}, $\approx$ \SI{400}{\celsius} & Alcène & Intramoléculaire \\
\hline
\end{tabularx}
\end{center}

\begin{methode}[Choisir le produit selon la température]
\begin{enumerate}
\item Lis la température indiquée dans l'énoncé.
\item Si $T \approx$ \SI{180}{\celsius} (ou alumine à \SI{400}{\celsius}) : écris \textbf{une} molécule d'alcool $\to$ alcène + \ce{H2O}. Cherche le carbone voisin le moins hydrogéné pour appliquer la règle de Zaïtsev si l'alcool est dissymétrique.
\item Si $T \approx$ \SI{140}{\celsius} : écris \textbf{deux} molécules d'alcool $\to$ étheroxyde \ce{R-O-R} + \ce{H2O}.
\item Porte toujours sur la flèche : \ce{H2SO4} concentré + température.
\item Vérifie la conservation des atomes (C, H, O) des deux côtés.
\end{enumerate}
\end{methode}

\begin{popfigure}
\begin{center}
\begin{tikzpicture}[scale=0.95, transform shape, >=Stealth]
% Nœud départ
\node[font=\large] (eth) at (0,0) {\chemfig{CH_3-CH_2-OH}};
% Flèche haute : 180 °C
\draw[->, line width=1.2pt, popink] (1.8,0.5) -- (5.2,2.3)
  node[midway, above, sloped, font=\small, text=popink] {\ce{H2SO4} conc., \SI{180}{\celsius}};
\node[font=\large] at (7.2,2.7) {\chemfig{CH_2=CH_2} \; $+$ \; \ce{H2O}};
\node[font=\small\itshape, text=popink] at (7.2,3.5) {Alcène (intramoléculaire)};
% Flèche basse : 140 °C
\draw[->, line width=1.2pt, popblue] (1.8,-0.5) -- (5.2,-2.3)
  node[midway, below, sloped, font=\small, text=popblue] {\ce{H2SO4} conc., \SI{140}{\celsius}, excès alcool};
\node[font=\large] at (7.5,-2.8) {\chemfig{CH_3-CH_2-O-CH_2-CH_3} \; $+$ \; \ce{H2O}};
\node[font=\small\itshape, text=popblue] at (7.5,-3.6) {Étheroxyde (intermoléculaire)};
\end{tikzpicture}
\end{center}
\legende{Les deux voies de déshydratation de l'éthanol : la température oriente la réaction vers l'éthène (\SI{180}{\celsius}) ou vers l'éthoxyéthane (\SI{140}{\celsius}).}
\end{popfigure}

\begin{popfigure}
\begin{center}
\begin{tikzpicture}[scale=0.9, transform shape, >=Stealth, every node/.style={font=\small}]
% --- Mécanisme intra ---
\node[font=\bfseries, text=popink] at (0,4.2) {Intramoléculaire : 1 molécule $\to$ Alcène};
\node (intra_start) at (0,3.0) {\chemfig{CH_2(-[2]H)-CH_2(-[6]OH)}};
% Flèches courbes pour départ H et OH
\draw[poporange, thick, ->] (0.7,3.4) to[bend left=45] (1.2,2.6);
\draw[poporange, thick, ->] (-0.5,2.5) to[bend right=40] (0.2,2.9);
\node[text=poporange, font=\small] at (2.0,3.0) {Élimination \ce{H2O}};
\draw[popdark, thick, ->] (3.2,3.0) -- (4.2,3.0);
\node at (5.5,3.0) {\chemfig{CH_2=CH_2} \; $+$ \; \ce{H2O}};
% --- Mécanisme inter ---
\node[font=\bfseries, text=popblue] at (0,1.0) {Intermoléculaire : 2 molécules $\to$ Étheroxyde};
\node (inter_A) at (-0.8,0.0) {\chemfig{CH_3-CH_2-O(-[2]H)}};
\node (inter_B) at (2.8,0.0) {\chemfig{HO-CH_2-CH_3}};
\draw[poporange, thick, ->] (0.2,0.4) to[bend left=30] (1.8,0.4);
\node[text=poporange, font=\small] at (1.0,0.8) {Élimination \ce{H2O}};
\draw[popdark, thick, ->] (4.2,0.0) -- (5.2,0.0);
\node at (6.8,0.0) {\chemfig{CH_3-CH_2-O-CH_2-CH_3} \; $+$ \; \ce{H2O}};
\end{tikzpicture}
\end{center}
\legende{Mécanismes simplifiés : en haut, une seule molécule perd \ce{H2O} (formation de la double liaison \ce{C=C}) ; en bas, deux molécules se condensent en éliminant \ce{H2O} (formation du pont oxygène \ce{-O-}).}
\end{popfigure}

\courssub{Exemple chiffré : rendement en éthoxyéthane}

\begin{exoresolu}[Déshydratation intermoléculaire de l'éthanol]
On chauffe \SI{9,2}{g} d'éthanol pur à \SI{140}{\celsius} en présence d'acide sulfurique concentré en excès. On suppose la réaction totale et exclusivement intermoléculaire. Déterminer la masse maximale d'éthoxyéthane \ce{C2H5-O-C2H5} obtenue. Données : $M(\ce{C}) = \SI{12}{g.mol^{-1}}$, $M(\ce{H}) = \SI{1}{g.mol^{-1}}$, $M(\ce{O}) = \SI{16}{g.mol^{-1}}$.
\tcblower
\textbf{Étape 1 : équation-bilan.}
\[ \ce{2 C2H5OH} \xrightarrow[\text{140 °C}]{\text{H2SO4}} \ce{C2H5-O-C2H5 + H2O} \]
\textbf{Étape 2 : quantité d'éthanol.}
$M(\ce{C2H5OH}) = 2 \times 12 + 6 \times 1 + 16 = \SI{46}{g.mol^{-1}}$.
\[ n(\text{éthanol}) = \frac{m}{M} = \frac{9,2}{46} = \SI{0,20}{mol}. \]
\textbf{Étape 3 : tableau d'avancement.}
\begin{center}
\begin{tabular}{l|c|c|c}
\rowcolor{popblueL} & \ce{2 C2H5OH} & \ce{C4H10O} & \ce{H2O} \\
\hline
État initial (mol) & 0,20 & 0 & 0 \\
État final (mol) & $0,20 - 2x_{\max}$ & $x_{\max}$ & $x_{\max}$ \\
\end{tabular}
\end{center}
L'éthanol est le réactif limitant : $0,20 - 2x_{\max} = 0 \Rightarrow x_{\max} = \SI{0,10}{mol}$.
\textbf{Étape 4 : masse d'éther.}
$M(\ce{C4H10O}) = 4 \times 12 + 10 + 16 = \SI{74}{g.mol^{-1}}$.
\[ m = n \times M = 0,10 \times 74 = \boxed{\SI{7,4}{g}}. \]
\textbf{Remarque :} il faut \textbf{2 moles} d'alcool pour \textbf{1 mole} d'éther ; oublier ce facteur 2 conduirait au double (\SI{14,8}{g}), résultat faux.
\end{exoresolu}

\courssub{Applications industrielles et au laboratoire}

\begin{itemize}
\item \textbf{Préparation des alcènes au laboratoire} : la déshydratation intramoléculaire est la voie classique d'obtention des alcènes ; l'éthène ainsi produit est la matière première du polyéthylène, plastique omniprésent (sachets, jerricans, tuyaux au Cameroun).
\item \textbf{L'éthoxyéthane} (éther diéthylique) est un excellent \cle{solvant} organique apolaire, utilisé pour les extractions liquide-liquide en laboratoire (ex. extraction de la caféine, des huiles essentielles) ; il est très volatil ($T_{\text{éb}} = \SI{34,5}{\celsius}$) et extrêmement inflammable : on le manipule loin de toute flamme, sous hotte.
\item \textbf{Distinction des fonctions} : la déshydratation permet de différencier expérimentalement alcools et étheroxydes, isomères de fonction (ex. \ce{C2H5OH} et \ce{CH3-O-CH3}, tous deux \ce{C2H6O}).
\end{itemize}

\begin{savaistu}[Le premier anesthésiant de l'histoire]
L'éther diéthylique a changé l'histoire de la médecine : le 16 octobre 1846, à Boston, le dentiste William Morton l'utilisa pour endormir un patient avant une opération chirurgicale — c'était la \textbf{première anesthésie générale} réussie et démontrée publiquement. Avant cette date, les opérations se pratiquaient sur des patients pleinement conscients, maintenus de force. L'éther a régné sur les blocs opératoires pendant près d'un siècle avant d'être remplacé par des anesthésiques moins inflammables (halothane, isoflurane). Un composé issu d'une simple déshydratation d'alcool a ainsi soulagé des millions de souffrances humaines.
\end{savaistu}

\begin{aretenir}[L'essentiel de la section]
\begin{itemize}
\item La déshydratation d'un alcool élimine \ce{H2O} grâce à \ce{H2SO4} ou \ce{H3PO4} concentré, ou sur alumine \ce{Al2O3}.
\item \textbf{\SI{180}{\celsius}} (excès d'acide) $\to$ \textbf{alcène} (intramoléculaire) : \ce{C2H5OH -> CH2=CH2 + H2O}.
\item \textbf{\SI{140}{\celsius}} (excès d'alcool) $\to$ \textbf{étheroxyde} (intermoléculaire) : \ce{2 C2H5OH -> C2H5-O-C2H5 + H2O}.
\item Alcool dissymétrique : l'alcène majoritaire est le plus substitué (\cle{règle de Zaïtsev}).
\item Les conditions (acide, température) figurent \textbf{toujours} sur la flèche de l'équation.
\end{itemize}
\end{aretenir}

\begin{exercice}[Pour t'entraîner]
On déshydrate le propan-2-ol \chemfig{H_3C-CH(-[2]OH)-CH_3} puis le butan-2-ol, chacun à \SI{180}{\celsius} en présence d'acide phosphorique.
\begin{enumerate}
\item Écrire l'équation de déshydratation du propan-2-ol et nommer le produit. Combien d'alcènes différents peut-on obtenir ?
\item Écrire les deux alcènes possibles issus du butan-2-ol et préciser, en justifiant, lequel est majoritaire.
\item Quelle masse de butan-2-ol faut-il déshydrater totalement pour obtenir \SI{5,6}{g} de but-2-ène ? ($M(\text{butan-2-ol}) = \SI{74}{g.mol^{-1}}$ ; $M(\text{but-2-ène}) = \SI{56}{g.mol^{-1}}$.)
\end{enumerate}
\end{exercice}

\begin{corrige}[Exercice]
\begin{enumerate}
\item $\ce{CH3-CH(OH)-CH3} \xrightarrow[\text{180 °C}]{\text{H3PO4}} \ce{CH3-CH=CH2 + H2O}$. Le produit est le \textbf{propène}. Les deux carbones voisins du carbone fonctionnel sont équivalents (deux groupes \ce{CH3}) : un \textbf{seul} alcène est possible.
\item Butan-2-ol : \ce{CH3-CH(OH)-CH2-CH3}. Perte de H sur C1 $\to$ but-1-ène \ce{CH2=CH-CH2-CH3} ; perte de H sur C3 $\to$ but-2-ène \ce{CH3-CH=CH-CH3}. D'après la règle de Zaïtsev, le \textbf{but-2-ène}, double liaison la plus substituée (deux groupes alkyles), est majoritaire.
\item $n(\text{but-2-ène}) = \dfrac{5,6}{56} = \SI{0,10}{mol}$. L'équation montre $n(\text{alcool}) = n(\text{alcène}) = \SI{0,10}{mol}$, d'où $m = 0,10 \times 74 = \SI{7,4}{g}$ de butan-2-ol.
\end{enumerate}
\end{corrige}

\coursec{Oxydation des alcools : combustion et oxydation ménagée}

Dans la section précédente, nous avons vu qu'un alcool peut \emph{perdre} une molécule d'eau : soit entre deux molécules pour donner un étheroxyde, soit au sein d'une même molécule pour donner un alcène. Dans les deux cas, l'atome d'oxygène \emph{quitte} la molécule. Mais le groupe hydroxyle \ce{-OH} peut aussi subir le sort inverse : au lieu de perdre de l'oxygène, la molécule peut en \emph{gagner}, ou perdre des hydrogènes. C'est le domaine de l'\cle{oxydation}, sans doute le chapitre le plus riche de la chimie des alcools — et le plus fréquent au Baccalauréat.

L'intuition de départ est simple : le feu de l'alcool à brûler dans un réchaud de cuisine, le contrôle d'alcoolémie au bord de la route entre Douala et Yaoundé, le vinaigre qui se forme quand une bouteille de vin de palme reste ouverte trop longtemps... Tous ces phénomènes quotidiens sont des \textbf{oxydations d'alcools}. Mais ils ne se ressemblent pas : l'une détruit totalement la molécule (combustion), les autres la transforment avec délicatesse (oxydation ménagée). Distinguons-les soigneusement.

\courssub{La combustion complète : une oxydation brutale}

\begin{definition}[Combustion complète d'un alcool]
La \cle{combustion complète} d'un alcool est sa réaction avec le dioxygène de l'air, qui détruit entièrement le squelette carboné. Les seuls produits sont le \textbf{dioxyde de carbone} \ce{CO2} et l'\textbf{eau} \ce{H2O}. C'est une réaction \textbf{vive} (elle s'accompagne d'une flamme) et fortement \textbf{exothermique} (elle libère de la chaleur).
\end{definition}

Pour un alcool saturé de formule brute \ce{C_{n}H_{2n+1}OH}, l'équation-bilan générale s'écrit :

\[ \ce{C_{n}H_{2n+1}OH + \frac{3n}{2} O2 -> n CO2 + (n+1) H2O} \]

\begin{exemplebox}[Combustion de l'éthanol]
Pour l'éthanol ($n = 2$), \ce{C2H5OH} :
\[ \ce{C2H5OH + 3 O2 -> 2 CO2 + 3 H2O} \]
C'est cette réaction qui chauffe la marmite lorsqu'on utilise de l'alcool à brûler (alcool dénaturé à \SI{90}{\degree}) dans un réchaud. Vérifions l'équilibre : à gauche $2$~C, $6$~H, $1 + 6 = 7$~O ; à droite $2$~C, $6$~H, $4 + 3 = 7$~O. L'équation est bien équilibrée.
\end{exemplebox}

\begin{attention}[Combustion incomplète]
Si le dioxygène manque (flamme mal réglée, espace confiné), la combustion devient \textbf{incomplète} : il se forme du monoxyde de carbone \ce{CO}, gaz incolore, inodore et \emph{mortel}, ainsi que du carbone (suie noire). C'est pourquoi on ne doit jamais utiliser un réchaud à alcool dans une pièce non ventilée. Au Baccalauréat, sauf indication contraire, « combustion » signifie toujours combustion \textbf{complète}.
\end{attention}

La combustion est énergétiquement précieuse, mais chimiquement « grossière » : elle ne nous apprend rien sur la structure de l'alcool, puisque tous les alcools finissent en \ce{CO2} et \ce{H2O}. Pour obtenir des produits organiques intéressants — et pour distinguer les classes d'alcools — il faut une oxydation plus douce.

\courssub{L'oxydation ménagée : définition et oxydants}

\begin{definition}[Oxydation ménagée]
Une \cle{oxydation ménagée} est une oxydation qui transforme le groupe fonctionnel d'un alcool \textbf{sans rompre le squelette carboné} : la chaîne carbonée est conservée, seul le carbone fonctionnel (celui qui porte le \ce{-OH}) voit son degré d'oxydation augmenter.
\end{definition}

Les oxydants ménagés usuels au programme sont :
\begin{itemize}
\item le \textbf{permanganate de potassium} \ce{KMnO4} en milieu acide : ion \ce{MnO4-} de couleur \textbf{violette} ;
\item le \textbf{dichromate de potassium} \ce{K2Cr2O7} en milieu acide : ion \ce{Cr2O7^2-} de couleur \textbf{orange} ;
\item le \textbf{dioxygène} \ce{O2} de l'air en présence d'un catalyseur (cuivre incandescent, par exemple — ce sera la lampe sans flamme de la section suivante).
\end{itemize}

L'oxydation ménagée d'un alcool consiste, sur le plan structural, à \textbf{retirer deux atomes d'hydrogène} : celui du groupe \ce{-OH} et celui porté par le carbone fonctionnel. Tout se joue donc sur une question unique : \emph{le carbone fonctionnel porte-t-il un hydrogène ? Et combien ?} C'est exactement ce qui distingue les trois classes d'alcools.

\courssub{Oxydation ménagée selon la classe de l'alcool}

\textbf{Alcool primaire : deux étapes successives.}\par
Un alcool primaire \ce{R-CH2OH} possède \textbf{deux} hydrogènes sur le carbone fonctionnel. L'oxydation se fait donc en \textbf{deux étapes} :
\begin{enumerate}
\item l'alcool donne d'abord un \textbf{aldéhyde} \ce{R-CHO} (perte des deux premiers hydrogènes) ;
\item si l'oxydant est en excès, l'aldéhyde est à son tour oxydé en \textbf{acide carboxylique} \ce{R-COOH}.
\end{enumerate}

\[ \ce{R-CH2OH ->[+\text{O}] R-CHO ->[+\text{O}] R-COOH} \]

Pour arrêter l'oxydation au stade aldéhyde, on utilise un oxydant en défaut et on distille l'aldéhyde au fur et à mesure de sa formation (il est plus volatil que l'alcool).

\textbf{Alcool secondaire : une seule étape.}\par
Un alcool secondaire \ce{R-CHOH-R$'$} ne possède qu'\textbf{un} hydrogène sur le carbone fonctionnel. Il est oxydé en \textbf{cétone} \ce{R-CO-R$'$} en une seule étape :

\[ \ce{R-CHOH-R$'$ ->[+\text{O}] R-CO-R$'$} \]

La cétone formée \textbf{n'est pas oxydable} en conditions ménagées : son carbone fonctionnel ne porte plus d'hydrogène. L'oxydation s'arrête donc naturellement là.

\textbf{Alcool tertiaire : aucune oxydation ménagée.}\par
Un alcool tertiaire ne possède \textbf{aucun} hydrogène sur son carbone fonctionnel : il ne subit \textbf{pas} d'oxydation ménagée. Si l'on insiste avec un oxydant puissant à chaud, la molécule se rompt (dégradation du squelette) : ce n'est plus une oxydation ménagée, et c'est hors programme.

\begin{popfigure}
\begin{center}
\begin{tabularx}{\linewidth}{|l|X|X|X|}
\hline
\rowcolor{popblueL} \textbf{Classe d'alcool} & \textbf{Exemple} & \textbf{Oxydant ménagé} & \textbf{Produit obtenu} \\
\hline
Primaire \ce{R-CH2OH} & éthanol \ce{C2H5OH} & \ce{MnO4-} ou \ce{Cr2O7^2-} (défaut) & aldéhyde \ce{R-CHO} \\
\hline
Primaire \ce{R-CH2OH} & éthanol \ce{C2H5OH} & \ce{MnO4-} ou \ce{Cr2O7^2-} (excès) & acide carboxylique \ce{R-COOH} \\
\hline
Secondaire \ce{R-CHOH-R$'$} & propan-2-ol & \ce{MnO4-} ou \ce{Cr2O7^2-} & cétone \ce{R-CO-R$'$} (étape unique) \\
\hline
\rowcolor{poporangeL} Tertiaire & 2-méthylpropan-2-ol & \ce{MnO4-} ou \ce{Cr2O7^2-} & \textbf{pas de réaction} \\
\hline
\end{tabularx}
\end{center}
\legende{Tableau récapitulatif : le produit de l'oxydation ménagée dépend de la classe de l'alcool et de la quantité d'oxydant.}
\end{popfigure}

\begin{attention}[Les deux erreurs classiques du Baccalauréat]
\begin{itemize}
\item \textbf{Ne jamais écrire} qu'un alcool tertiaire donne une cétone ou un acide : il ne réagit pas en oxydation ménagée. Dire « l'alcool tertiaire résiste à l'oxydation ménagée » est la formulation attendue.
\item \textbf{Ne jamais poursuivre} l'oxydation d'une cétone vers un « acide » : une cétone est un point d'arrivée. Seuls les aldéhydes (issus des alcools primaires) peuvent être oxydés davantage, en acides carboxyliques.
\end{itemize}
\end{attention}

\courssub{Suivi du degré d'oxydation : la cascade de l'éthanol}

Pour bien comprendre la progression, suivons le \textbf{nombre d'oxydation} (n.o.) du carbone fonctionnel de l'éthanol au fil de son oxydation : il passe de $-1$ dans l'alcool, à $+1$ dans l'aldéhyde, puis à $+3$ dans l'acide carboxylique. Chaque étape correspond à une perte de deux électrons.

\begin{popfigure}
\begin{center}
\begin{tikzpicture}[node distance=0.6cm and 1.6cm]
\node (alc) {\chemfig{CH_3-CH_2-OH}};
\node[below=0.15cm of alc, popblue] {\small éthanol, n.o. $= -1$};
\node[right=of alc] (ald) {\chemfig{CH_3-C(=[1]O)-[7]H}};
\node[below=0.15cm of ald, poporange] {\small éthanal, n.o. $= +1$};
\node[right=of ald] (ac) {\chemfig{CH_3-C(=[1]O)-[7]OH}};
\node[below=0.15cm of ac, popgreen] {\small acide éthanoïque, n.o. $= +3$};
\draw[-{Stealth}, very thick, popdark] (alc) -- (ald) node[midway, above, font=\small] {$-2\,e^-,\ -2\,$H};
\draw[-{Stealth}, very thick, popdark] (ald) -- (ac) node[midway, above, font=\small] {$-2\,e^-,\ +\,$O};
\end{tikzpicture}
\end{center}
\legende{Cascade d'oxydation de l'éthanol : alcool primaire $\rightarrow$ aldéhyde $\rightarrow$ acide carboxylique. Le nombre d'oxydation du carbone fonctionnel augmente de deux unités à chaque étape.}
\end{popfigure}

\begin{savaistu}[Du vin de palme au vinaigre]
Le vin de palme laissé à l'air libre « tourne » et devient acide : le dioxygène de l'air, aidé par des bactéries du genre \emph{Acetobacter}, oxyde l'éthanol en acide éthanoïque \ce{CH3COOH} — c'est exactement la cascade ci-dessus, réalisée par des micro-organismes. C'est d'ailleurs le principe industriel de la fabrication du vinaigre (du latin \emph{vinum acetum}, « vin aigre »).
\end{savaistu}

\courssub{Écriture des équations-bilans : demi-équations électroniques}

Les couples redox mis en jeu sont, pour les espèces organiques :
\begin{itemize}
\item couple aldéhyde/alcool primaire : \ce{R-CHO}/\ce{R-CH2OH} ;
\item couple acide carboxylique/aldéhyde : \ce{R-COOH}/\ce{R-CHO} ;
\item couple cétone/alcool secondaire : \ce{R-CO-R$'$}/\ce{R-CHOH-R$'$}.
\end{itemize}
Et pour les oxydants minéraux, en milieu acide :
\begin{itemize}
\item \ce{MnO4-}/\ce{Mn^2+} : \quad \ce{MnO4- + 8 H+ + 5 e- -> Mn^2+ + 4 H2O} ;
\item \ce{Cr2O7^2-}/\ce{Cr^3+} : \quad \ce{Cr2O7^2- + 14 H+ + 6 e- -> 2 Cr^3+ + 7 H2O}.
\end{itemize}

\begin{methode}[Équilibrer une équation redox en 4 étapes]
\begin{enumerate}
\item \textbf{Écrire les deux demi-équations.} Pour le couple organique, on équilibre d'abord les atomes autres que O et H, puis les O avec \ce{H2O}, puis les H avec \ce{H+}, puis les charges avec les électrons $e^-$.
\item \textbf{Multiplier} chaque demi-équation pour que le nombre d'électrons cédés égale le nombre d'électrons captés.
\item \textbf{Additionner} membre à membre et simplifier les \ce{H+} et \ce{H2O} qui apparaissent des deux côtés.
\item \textbf{Vérifier} la conservation des atomes \emph{et} des charges électriques.
\end{enumerate}
\end{methode}

\begin{exoresolu}[Oxydation de l'éthanol en acide éthanoïque par le permanganate]
Écrire l'équation-bilan de l'oxydation de l'éthanol en acide éthanoïque par les ions permanganate en milieu acide.
\tcblower
\textbf{Étape 1 — demi-équation organique} (couple \ce{CH3COOH}/\ce{C2H5OH}) :
\[ \ce{C2H5OH + H2O -> CH3COOH + 4 H+ + 4 e-} \]
Vérification : $2$~C de chaque côté ; O : $1+1=2$ à gauche, $2$ à droite ; H : $6+2=8$ à gauche, $4+4=8$ à droite ; charges : $0$ à gauche, $+4-4=0$ à droite.

\textbf{Étape 2 — demi-équation de l'oxydant} :
\[ \ce{MnO4- + 8 H+ + 5 e- -> Mn^2+ + 4 H2O} \]

\textbf{Étape 3 — combinaison.} On multiplie la première par $5$ et la seconde par $4$ (p.p.c.m. de $4$ et $5$ : $20$ électrons) :
\begin{align*}
\ce{5 C2H5OH + 5 H2O &-> 5 CH3COOH + 20 H+ + 20 e-}\\
\ce{4 MnO4- + 32 H+ + 20 e- &-> 4 Mn^2+ + 16 H2O}
\end{align*}
En additionnant et en simplifiant ($20$ \ce{H+} et $5$ \ce{H2O} se simplifient) :
\[ \boxed{\ce{5 C2H5OH + 4 MnO4- + 12 H+ -> 5 CH3COOH + 4 Mn^2+ + 11 H2O}} \]
\textbf{Étape 4 — vérification des charges} : à gauche $-4 + 12 = +8$ ; à droite $4 \times (+2) = +8$. L'équation est correcte. \textbf{Observation} : la solution violette de permanganate se \textbf{décolore} (les ions \ce{Mn^2+} sont quasiment incolores) : c'est le signe que la réaction a eu lieu.
\end{exoresolu}

\begin{exemplebox}[Exemple chiffré : oxydation du propan-2-ol par le dichromate]
Oxydons le propan-2-ol (alcool secondaire) par les ions dichromate en milieu acide. Demi-équation organique (couple \ce{CH3-CO-CH3}/\ce{CH3-CHOH-CH3}) :
\[ \ce{CH3-CHOH-CH3 -> CH3-CO-CH3 + 2 H+ + 2 e-} \]
Demi-équation du dichromate : \ce{Cr2O7^2- + 14 H+ + 6 e- -> 2 Cr^3+ + 7 H2O}. On multiplie la première par $3$ (pour échanger $6$ électrons) et on additionne ; les $6$ \ce{H+} produits se retranchent des $14$ \ce{H+} consommés :
\[ \ce{3 CH3-CHOH-CH3 + Cr2O7^2- + 8 H+ -> 3 CH3-CO-CH3 + 2 Cr^3+ + 7 H2O} \]
Vérification des charges : gauche $-2 + 8 = +6$ ; droite $2 \times (+3) = +6$. Correct. Le produit est la \textbf{propanone} (acétone), et l'oxydation s'arrête là.
\end{exemplebox}

\courssub{Observations expérimentales et application : l'éthylotest}

Ces réactions ne sont pas que des équations sur le papier : elles s'accompagnent de \textbf{changements de couleur spectaculaires}, faciles à observer en TP :
\begin{itemize}
\item le permanganate de potassium, \textbf{violet}, se \textbf{décolore} en donnant des ions \ce{Mn^2+} presque incolores ;
\item le dichromate de potassium, \textbf{orange}, vire au \textbf{vert} en donnant des ions chrome(III) \ce{Cr^3+}.
\end{itemize}

\begin{savaistu}[L'alcootest des forces de l'ordre]
Les anciens éthylotests (ballons chimiques) utilisés lors des contrôles routiers contiennent du dichromate de potassium acidifié, de couleur orange, déposé sur un gel de silice. L'automobiliste souffle : si son haleine contient des vapeurs d'éthanol, celui-ci est oxydé en acide éthanoïque tandis que le dichromate est réduit en ions \ce{Cr^3+} verts. Plus la tache verte avance loin dans le tube, plus la quantité d'alcool est élevée. La chimie redox des alcools veille ainsi sur la sécurité routière, de Yaoundé à Garoua !
\end{savaistu}

\begin{aretenir}[L'essentiel de la section]
\begin{itemize}
\item La \textbf{combustion complète} détruit le squelette : \ce{C_{n}H_{2n+1}OH + \frac{3n}{2} O2 -> n CO2 + (n+1) H2O}, réaction vive et exothermique.
\item L'\textbf{oxydation ménagée} conserve le squelette carboné et dépend de la \textbf{classe} de l'alcool : primaire $\rightarrow$ aldéhyde $\rightarrow$ acide carboxylique (oxydant en excès) ; secondaire $\rightarrow$ cétone (étape unique) ; tertiaire $\rightarrow$ \textbf{pas de réaction}.
\item Oxydants usuels : \ce{MnO4-}/\ce{H+} (violet $\rightarrow$ incolore) et \ce{Cr2O7^2-}/\ce{H+} (orange $\rightarrow$ vert, principe de l'éthylotest).
\item Toute équation-bilan s'obtient par la méthode des demi-équations : équilibrer O par \ce{H2O}, H par \ce{H+}, charges par $e^-$, puis combiner.
\end{itemize}
\end{aretenir}

Nous savons désormais oxyder un alcool « en solution », avec des oxydants minéraux. Mais il existe une expérience magnifique où l'oxydation se fait \emph{en phase gazeuse}, sur un fil de cuivre incandescent qui continue de rougeoyer sans aucune flamme : c'est la fameuse \textbf{lampe sans flamme}, que nous étudierons dans la section suivante, avec les tests d'identification des produits formés.

\coursec{Lampe sans flamme et identification des produits}

Dans la section précédente, nous avons vu que l'oxydation ménagée d'un alcool primaire par le permanganate ou le dichromate en milieu acide conduit d'abord à un aldéhyde, puis à un acide carboxylique. Ces oxydants sont des réactifs \emph{en solution}. Mais il existe une expérience spectaculaire, historique et très instructive, qui réalise cette oxydation \emph{en phase gazeuse}, sans flamme, grâce à un simple fil de cuivre chauffé : c'est l'expérience de la \cle{lampe sans flamme}. Elle va nous permettre de produire de l'éthanal à partir de l'éthanol, puis de répondre à une question fondamentale en chimie organique : \emph{comment identifier expérimentalement les produits formés ?} C'est tout l'objet de cette section : réaliser l'expérience, comprendre le rôle du cuivre, puis maîtriser les tests caractéristiques des aldéhydes et des cétones.

\courssub{Principe de l'oxydation catalytique sur cuivre}

\textbf{Intuition : un fil qui brille sans brûler}\par

Imagine la scène : un bécher contient un peu d'éthanol. On chauffe un fil de cuivre enroulé en spirale jusqu'à ce qu'il devienne rouge incandescent, puis on le suspend juste au-dessus de la surface du liquide, \emph{sans le toucher}. Résultat étonnant : le fil \textbf{reste incandescent}, alors qu'aucune flamme n'est visible, et une \textbf{odeur piquante} se dégage du bécher. Que se passe-t-il ?

L'éthanol s'évapore légèrement ; ses vapeurs se mélangent à l'air (donc au dioxygène \ce{O2}) et viennent au contact du cuivre chaud. Le cuivre joue alors le rôle de \cle{catalyseur} : il permet au dioxygène d'arracher deux atomes d'hydrogène à la molécule d'éthanol (celui du groupe hydroxyle et celui porté par le carbone fonctionnel), la transformant en \cle{éthanal} \ce{CH3CHO}, l'aldéhyde à l'odeur piquante caractéristique. La réaction est \textbf{exothermique} : c'est la chaleur qu'elle dégage qui maintient le fil au rouge, d'où le nom de \emph{lampe sans flamme} — la « combustion » a lieu, mais de façon contrôlée, à la surface du métal, sans flamme.

\begin{definition}[Catalyseur]
Un \cle{catalyseur} est une espèce chimique qui \textbf{accélère} une réaction chimique \textbf{sans figurer dans son équation-bilan} : il est consommé au cours d'une étape puis \textbf{régénéré} en fin de réaction. Il n'apparaît ni parmi les réactifs, ni parmi les produits du bilan global. Dans l'expérience de la lampe sans flamme, le catalyseur est le \textbf{cuivre} \ce{Cu} métallique. On parle de \emph{catalyse hétérogène} car le catalyseur (solide) et les réactifs (gazeux) ne sont pas dans la même phase.
\end{definition}

\begin{attention}[Ce qu'un catalyseur ne fait pas]
Un catalyseur ne rend \textbf{pas} possible une réaction impossible : il accélère une réaction déjà spontanée mais lente. Il ne modifie ni le bilan de matière, ni la quantité maximale de produit obtenue ; il permet seulement de l'atteindre plus vite. Le retrouver intact en fin d'expérience (le fil de cuivre est réutilisable) est d'ailleurs une preuve expérimentale de sa régénération.
\end{attention}

\textbf{Mécanisme en deux étapes et bilan}\par

Pour comprendre comment le cuivre « disparaît puis réapparaît », décomposons la transformation en deux étapes.

\textbf{Étape 1 — oxydation du cuivre par le dioxygène de l'air.} Le fil de cuivre chaud réagit avec \ce{O2} et se recouvre d'une couche noire d'oxyde de cuivre(II) \ce{CuO} :
\[ \ce{2 Cu + O2 -> 2 CuO} \]
C'est pourquoi le fil, brillant au départ, noircit lorsqu'on le chauffe à l'air.

\textbf{Étape 2 — oxydation de l'éthanol par l'oxyde de cuivre(II).} Les vapeurs d'éthanol arrachent l'oxygène de \ce{CuO}, qui est ainsi réduit en cuivre métallique (le fil redevient brillant) :
\[ \ce{C2H5OH + CuO -> CH3CHO + Cu + H2O} \]

En additionnant l'étape 1 et deux fois l'étape 2 (pour éliminer \ce{CuO} et \ce{Cu}, qui sont intermédiaires), on obtient l'\textbf{équation-bilan} de l'oxydation catalytique de l'éthanol :
\[ \ce{2 C2H5OH + O2 ->[Cu][chaud] 2 CH3CHO + 2 H2O} \]

Vérifie le bilan : 4 atomes de carbone, 12 d'hydrogène et 4 d'oxygène de chaque côté. Le cuivre n'y figure pas : il a bien été \textbf{régénéré}, conformément à la définition du catalyseur. Remarque aussi que cette réaction est une \emph{déshydrogénation} : l'alcool perd deux atomes d'hydrogène, ce qui est équivalent, du point de vue du degré d'oxydation, à une oxydation ménagée de l'alcool primaire en aldéhyde.

\begin{experience}[La lampe sans flamme]
\textbf{Dispositif et protocole.}
\begin{enumerate}
\item Verser environ \SI{5}{mL} d'éthanol dans un bécher (ou un erlenmeyer) propre et sec.
\item Enrouler un fil de cuivre décapé en spirale autour d'un crayon, le fixer à un support.
\item Chauffer la spirale à la flamme du bec Bunsen jusqu'à incandescence (rouge vif) : elle noircit (formation de \ce{CuO}).
\item Introduire immédiatement la spirale dans le bécher, \textbf{au-dessus} de la surface de l'éthanol, sans toucher le liquide.
\end{enumerate}
\textbf{Observations.} Le fil \textbf{reste incandescent} durablement, sans qu'aucune flamme n'apparaisse ; il alterne des phases noires et brillantes ; une \textbf{odeur piquante} (pomme verte âcre) se dégage : c'est l'éthanal. Le liquide recueilli trouble le réactif de Schiff et donne un précipité jaune avec la DNPH.
\textbf{Interprétation.} La réaction \ce{2 C2H5OH + O2 -> 2 CH3CHO + 2 H2O}, exothermique, fournit la chaleur qui maintient le fil au rouge. Le cuivre, alternativement oxydé en \ce{CuO} puis réduit en \ce{Cu}, joue le rôle de catalyseur.
\end{experience}

\begin{popfigure}
\begin{center}
\begin{tikzpicture}[scale=1.0, every node/.style={font=\small}]
% Bécher
\draw[thick, popdark] (-1.6,0) -- (-1.6,2.6);
\draw[thick, popdark] (1.6,0) -- (1.6,2.6);
\draw[thick, popdark] (-1.6,0) -- (1.6,0);
% Ethanol liquide
\fill[popblueL] (-1.55,0.05) rectangle (1.55,1.1);
\draw[popblue, thick] (-1.55,1.1) -- (1.55,1.1);
\node[popblue!70!black] at (0,0.55) {Éthanol liquide};
% Vapeurs
\foreach \x in {-1.1,-0.4,0.35,1.0}{
  \draw[popmuted, ->, decorate, decoration={snake, amplitude=0.6mm, segment length=2.5mm}] (\x,1.25) -- (\x,1.9);
}
\node[popmuted, align=center] at (0,2.15) {vapeurs d'éthanol + air (\ce{O2})};
% Fil de cuivre en spirale
\draw[poporange, very thick, decorate, decoration={coil, aspect=0.4, segment length=3.2mm, amplitude=2.2mm}] (0,5.6) -- (0,3.1);
% Halo incandescent
\fill[popgoldL, opacity=0.55] (0,3.35) ellipse (0.75 and 0.5);
\draw[poporange, very thick, decorate, decoration={coil, aspect=0.4, segment length=3.2mm, amplitude=2.2mm}] (0,5.6) -- (0,3.1);
% Support
\draw[thick, popdark] (0,5.6) -- (0,6.3);
\draw[thick, popdark] (-0.9,6.3) -- (0.9,6.3);
\fill[poppurple] (0,6.3) circle (1.6pt);
\node[poppurple, anchor=south] at (0,6.35) {support};
% Annotations
\node[poporange!80!black, anchor=west, align=left] at (2.0,4.4) {Fil de cuivre en spirale\\ \textbf{incandescent} (catalyseur)};
\draw[poporange!80!black, ->] (2.0,4.25) -- (0.55,3.6);
\node[popdark, anchor=west, align=left] at (2.0,2.6) {Aucune flamme :\\ réaction exothermique\\ en surface du métal};
\draw[popdark, ->] (2.0,2.45) -- (1.0,2.9);
\node[popgreen!70!black, anchor=west, align=left] at (-5.4,0.9) {Odeur piquante :\\ éthanal \ce{CH3CHO} + eau};
\draw[popgreen!70!black, ->] (-1.9,0.95) -- (-1.45,1.35);
\end{tikzpicture}
\end{center}
\legende{Montage de la lampe sans flamme : le fil de cuivre porté au rouge, placé au-dessus de l'éthanol, reste incandescent grâce à la chaleur dégagée par l'oxydation catalytique des vapeurs d'alcool.}
\end{popfigure}

\begin{savaistu}[Du laboratoire à l'industrie]
L'oxydation catalytique des alcools sur métal (cuivre, argent) est un procédé industriel majeur : le méthanal (formaldéhyde), obtenu à partir du méthanol sur catalyseur d'argent vers \SI{600}{\celsius}, sert à fabriquer les colles et résines des panneaux de bois aggloméré. Le même principe — un gaz, un métal chaud, pas de flamme — est à l'œuvre dans le pot catalytique des automobiles, où le platine et le palladium oxydent les vapeurs d'essence imbrûlées.
\end{savaistu}

\courssub{Identification des produits : les tests des composés carbonylés}

L'odeur piquante suggère la formation d'un aldéhyde, mais en science, une odeur n'est pas une preuve. Il faut des \textbf{tests d'identification} reproductibles, fondés sur des réactions chimiques caractéristiques. Tous ces tests ciblent le \cle{groupe carbonyle} \ce{C=O} et la propriété \emph{réductrice} des aldéhydes.

\textbf{Le test à la DNPH : détecter le groupe carbonyle}\par

La \cle{2,4-dinitrophénylhydrazine} (DNPH, ou réactif de Brady) réagit avec le groupe carbonyle \ce{C=O} pour former une 2,4-dinitrophénylhydrazone, solide \textbf{jaune-orangé} qui précipite.

\begin{propriete}[Test à la DNPH]
Tout composé carbonylé — \textbf{aldéhyde ou cétone} — donne avec la DNPH un \textbf{précipité jaune-orangé}. Ce test est donc positif pour les deux familles : il prouve la présence d'un groupe \ce{C=O}, mais ne permet \textbf{pas} de distinguer un aldéhyde d'une cétone.
\end{propriete}

\begin{attention}[Erreur fréquente au baccalauréat]
Conclure « le composé est un aldéhyde » sur le seul fondement d'un test DNPH positif est une \textbf{erreur} : les cétones réagissent aussi avec la DNPH ! Le test DNPH positif permet seulement d'affirmer : « le composé est un \emph{composé carbonylé} (aldéhyde \emph{ou} cétone) ». Il faut un second test pour trancher.
\end{attention}

\textbf{Distinguer aldéhyde et cétone : le caractère réducteur}\par

La différence essentielle entre les deux familles : un \textbf{aldéhyde est un réducteur} (il peut être oxydé en acide carboxylique), alors qu'une \textbf{cétone ne l'est pas} (elle résiste aux oxydants ménagers). Trois tests exploitent cette propriété : l'aldéhyde réduit l'ion cuivre(II) de la liqueur de Fehling, l'ion argent(I) du réactif de Tollens, et régénère la coloration du réactif de Schiff.

\begin{propriete}[Les trois tests des aldéhydes]
\begin{itemize}
\item \textbf{Liqueur de Fehling} (solution bleue contenant des ions \ce{Cu^2+} complexés, en milieu basique) : chauffée avec un aldéhyde, elle donne un \textbf{précipité rouge brique} d'oxyde de cuivre(I) \ce{Cu2O}. L'aldéhyde est oxydé en ion carboxylate.
\item \textbf{Réactif de Tollens} (nitrate d'argent ammoniacal, ions \ce{[Ag(NH3)2]+}) : chauffé doucement avec un aldéhyde, il donne un \textbf{miroir d'argent} (dépôt d'argent métallique \ce{Ag} sur les parois du tube).
\item \textbf{Réactif de Schiff} (fuchsine décolorée) : un aldéhyde lui redonne une \textbf{coloration rose-violette} à froid.
\end{itemize}
Une \textbf{cétone} ne réagit avec \emph{aucun} de ces trois réactifs : ses tests sont négatifs.
\end{propriete}

\begin{attention}[Sécurité et bon usage du réactif de Tollens]
Le réactif de Tollens doit être \textbf{préparé frais}, juste avant l'emploi : en vieillissant, il peut former un composé \emph{explosif} (fulminate d'argent). Le chauffage du tube se fait \textbf{au bain-marie} (vers \SI{60}{\celsius}), \textbf{jamais directement à la flamme}. Après le test, on détruit l'excès de réactif et on ne conserve jamais la solution.
\end{attention}

\begin{methode}[Identifier un composé carbonylé inconnu en deux temps]
\begin{enumerate}
\item \textbf{Premier test : DNPH.} Ajouter quelques gouttes de DNPH à l'échantillon.
   \begin{itemize}
   \item Précipité jaune-orangé $\Rightarrow$ le composé est un \textbf{carbonyle} (aldéhyde ou cétone) : passer à l'étape 2.
   \item Pas de précipité $\Rightarrow$ ce n'est ni un aldéhyde ni une cétone (c'est peut-être un alcool, un acide\ldots).
   \end{itemize}
\item \textbf{Second test : liqueur de Fehling (ou Tollens, ou Schiff).} Chauffer l'échantillon avec la liqueur de Fehling au bain-marie.
   \begin{itemize}
   \item Précipité rouge brique $\Rightarrow$ le composé est un \textbf{aldéhyde} (réducteur).
   \item La solution reste bleue $\Rightarrow$ le composé est une \textbf{cétone}.
   \end{itemize}
\end{enumerate}
\textbf{Conclusion logique :} un aldéhyde répond positivement à la DNPH \emph{et} aux tests réducteurs ; une cétone ne répond \emph{qu'}à la DNPH.
\end{methode}

\begin{popfigure}
\begin{center}
\begin{tikzpicture}[x=4.6cm, y=1.05cm, every node/.style={font=\small, align=center}]
% En-têtes
\fill[popblue] (0,3) rectangle (1,3.7);
\fill[poppurple] (1,3) rectangle (2,3.7);
\fill[popgreen] (2,3) rectangle (3,3.7);
\node[white, font=\small\bfseries] at (0.5,3.35) {Réactif};
\node[white, font=\small\bfseries] at (1.5,3.35) {Aldéhyde \ce{RCHO}};
\node[white, font=\small\bfseries] at (2.5,3.35) {Cétone \ce{R2CO}};
% Ligne DNPH
\fill[popgoldL] (0,2) rectangle (1,3);
\fill[popgoldL] (1,2) rectangle (2,3);
\fill[popgoldL] (2,2) rectangle (3,3);
\node[font=\small\bfseries] at (0.5,2.5) {DNPH};
\node at (1.5,2.5) {Précipité\\ \textcolor{poporange!85!black}{\textbf{jaune-orangé}}};
\node at (2.5,2.5) {Précipité\\ \textcolor{poporange!85!black}{\textbf{jaune-orangé}}};
% Ligne Fehling
\fill[popblueL] (0,1) rectangle (1,2);
\fill[popblueL] (1,1) rectangle (2,2);
\fill[popblueL] (2,1) rectangle (3,2);
\node[font=\small\bfseries] at (0.5,1.5) {Liqueur de Fehling\\ (à chaud)};
\node at (1.5,1.5) {Précipité\\ \textcolor{poporange!70!black}{\textbf{rouge brique}} \ce{Cu2O}};
\node at (2.5,1.5) {Solution\\ \textcolor{popblue}{\textbf{bleue inchangée}}};
% Ligne Tollens
\fill[poppurpleL] (0,0) rectangle (1,1);
\fill[poppurpleL] (1,0) rectangle (2,1);
\fill[poppurpleL] (2,0) rectangle (3,1);
\node[font=\small\bfseries] at (0.5,0.5) {Réactif de Tollens\\ (bain-marie)};
\node at (1.5,0.5) {\textcolor{popmuted}{\textbf{Miroir d'argent}}\\ \ce{Ag} métallique};
\node at (2.5,0.5) {Aucun\\ dépôt};
% Ligne Schiff
\fill[poppinkL] (0,-1) rectangle (1,0);
\fill[poppinkL] (1,-1) rectangle (2,0);
\fill[poppinkL] (2,-1) rectangle (3,0);
\node[font=\small\bfseries] at (0.5,-0.5) {Réactif de Schiff\\ (à froid)};
\node at (1.5,-0.5) {Coloration\\ \textcolor{poppink}{\textbf{rose-violette}}};
\node at (2.5,-0.5) {Reste\\ incolore};
% Contours
\draw[popdark, thick] (0,-1) rectangle (3,3.7);
\foreach \y in {-1,0,1,2,3}{\draw[popdark] (0,\y) -- (3,\y);}
\draw[popdark] (1,-1) -- (1,3.7);
\draw[popdark] (2,-1) -- (2,3.7);
\end{tikzpicture}
\end{center}
\legende{Bilan des tests d'identification : la DNPH détecte le groupe carbonyle (aldéhydes \emph{et} cétones) ; seuls les aldéhydes, réducteurs, répondent positivement à la liqueur de Fehling, au réactif de Tollens et au réactif de Schiff.}
\end{popfigure}

\begin{savaistu}[Le miroir d'argent, du tube à essai à la salle de bain]
Le dépôt d'argent métallique obtenu avec le réactif de Tollens — le fameux « miroir d'argent » — n'est pas qu'une curiosité de laboratoire : cette réaction a été utilisée \textbf{industriellement} pour \emph{argenter les miroirs}. On faisait réagir une solution de nitrate d'argent ammoniacal avec un aldéhyde (souvent le glucose, qui possède une fonction aldehyde) directement sur une surface de verre soigneusement nettoyée : l'argent se déposait en une couche mince, uniforme et réfléchissante. Ce procédé, mis au point au XIX\textsuperscript{e} siècle par le chimiste allemand Justus von Liebig, a révolutionné la fabrication des miroirs, auparavant réalisés au mercure, très toxique.
\end{savaistu}

\courssub{Oxydation ménagée des aldéhydes}

Pourquoi les aldéhydes sont-ils réducteurs ? Parce que leur groupe carbonyle porte encore un atome d'hydrogène \ce{-CHO} qu'un oxydant ménager peut transformer en \ce{-COOH}. L'oxydation ménagée d'un aldéhyde conduit donc à l'\textbf{acide carboxylique} correspondant, sans rompre la chaîne carbonée :
\[ \ce{RCHO + [O] -> RCOOH} \]

Écrivons la \textbf{demi-équation électronique} du couple \ce{RCOOH}/\ce{RCHO} en milieu acide. On équilibre d'abord l'oxygène avec des molécules d'eau, puis l'hydrogène avec des ions \ce{H+}, puis les charges avec des électrons :
\[ \ce{RCHO + H2O -> RCOOH + 2 H+ + 2 e-} \]
L'aldéhyde \emph{cède} deux électrons : il est bien \textbf{réducteur}. C'est exactement ce qui se passe dans le test de Tollens, où l'oxydant est l'ion diammineargent(I) :
\[ \ce{[Ag(NH3)2]+ + e- -> Ag + 2 NH3} \]
En combinant les deux demi-équations (multiplier celle de l'argent par 2), on obtient le bilan du miroir d'argent :
\[ \ce{RCHO + 2 [Ag(NH3)2]+ + H2O -> RCOOH + 2 Ag + 2 NH4+ + 2 NH3} \]
De même, avec la liqueur de Fehling en milieu basique, l'ion \ce{Cu^2+} est réduit en \ce{Cu2O} rouge brique tandis que l'aldéhyde devient un ion carboxylate \ce{RCOO-}.

\begin{aretenir}[L'essentiel de la section]
\begin{itemize}
\item La lampe sans flamme réalise l'\textbf{oxydation catalytique} de l'éthanol en éthanal : \ce{2 C2H5OH + O2 ->[Cu] 2 CH3CHO + 2 H2O}. Le cuivre, catalyseur, est régénéré ; la réaction est exothermique.
\item La \textbf{DNPH} détecte le groupe carbonyle : précipité jaune-orangé avec les aldéhydes \emph{et} les cétones.
\item Un \textbf{aldéhyde est réducteur} : précipité rouge brique avec la liqueur de Fehling, miroir d'argent avec le réactif de Tollens, coloration rose-violette avec le réactif de Schiff. Une \textbf{cétone} ne réagit avec aucun de ces trois réactifs.
\item L'oxydation ménagée d'un aldéhyde donne l'acide carboxylique : \ce{RCHO + H2O -> RCOOH + 2 H+ + 2 e-}.
\end{itemize}
\end{aretenir}

\begin{exoresolu}[Identifier les produits de la lampe sans flamme]
\textbf{Énoncé.} On réalise l'expérience de la lampe sans flamme avec un bécher contenant $m = \SI{4,6}{g}$ d'éthanol. Le liquide recueilli après l'expérience est divisé en deux échantillons. Avec le premier, la DNPH donne un précipité jaune-orangé. Avec le second, chauffé au bain-marie, la liqueur de Fehling donne un précipité rouge brique.
\begin{enumerate}
\item Interpréter chaque test et identifier le produit organique formé.
\item Écrire l'équation-bilan de la réaction sur le cuivre.
\item En supposant la réaction totale, calculer la masse maximale de produit organique obtenu.
\end{enumerate}
Données : $M(\ce{C2H5OH}) = \SI{46}{g.mol^{-1}}$ ; $M(\ce{CH3CHO}) = \SI{44}{g.mol^{-1}}$.
\tcblower
\textbf{Solution détaillée.}
\begin{enumerate}
\item Le test à la DNPH positif prouve la présence d'un \textbf{groupe carbonyle} \ce{C=O} : le produit est un aldéhyde \emph{ou} une cétone. Le test à la liqueur de Fehling positif (précipité rouge brique de \ce{Cu2O}) prouve que ce composé carbonylé est \textbf{réducteur} : c'est donc un \textbf{aldéhyde}. L'oxydation de l'éthanol, alcool primaire, conduit à l'\textbf{éthanal} \ce{CH3CHO}. (L'eau est l'autre produit.)
\item Les deux étapes sur le catalyseur :
\[ \ce{2 Cu + O2 -> 2 CuO} \qquad \ce{C2H5OH + CuO -> CH3CHO + Cu + H2O} \]
d'où le bilan, le cuivre étant régénéré :
\[ \ce{2 C2H5OH + O2 ->[Cu][chaud] 2 CH3CHO + 2 H2O} \]
\item Quantité d'éthanol introduite :
\[ n(\ce{C2H5OH}) = \frac{m}{M} = \frac{4,6}{46} = \SI{0,10}{mol}. \]
D'après l'équation-bilan, $2$ mol d'éthanol donnent $2$ mol d'éthanal, soit une proportion $1:1$ : $n(\ce{CH3CHO}) = \SI{0,10}{mol}$. La masse maximale d'éthanal est donc :
\[ m(\ce{CH3CHO}) = n \times M = 0,10 \times 44 = \SI{4,4}{g}. \]
\end{enumerate}
\textbf{Remarque.} En pratique, une partie de l'éthanal s'évapore (il bout à \SI{21}{\celsius}) ou est oxydée davantage : la masse réellement recueillie est inférieure à \SI{4,4}{g}.
\end{exoresolu}

\begin{exercice}[À toi de jouer]
On dispose de trois flacons étiquetés A, B, C contenant chacun l'un des composés suivants : propanal, propanone, propan-1-ol. Avec la DNPH, A et B donnent un précipité jaune-orangé, C ne réagit pas. Avec la liqueur de Fehling à chaud, seul A donne un précipité rouge brique. Identifier A, B et C, puis écrire la demi-équation électronique de l'oxydation ménagée de A et le nom du produit obtenu.
\end{exercice}

\begin{corrige}[Exercice — À toi de jouer]
C ne réagit pas avec la DNPH : ce n'est pas un carbonyle, c'est donc le \textbf{propan-1-ol}. A et B sont des carbonylés ; seul A réduit la liqueur de Fehling, donc A est l'aldéhyde, le \textbf{propanal} \ce{CH3CH2CHO}, et B est la cétone, la \textbf{propanone} \ce{CH3COCH3}. Oxydation ménagée du propanal :
\[ \ce{CH3CH2CHO + H2O -> CH3CH2COOH + 2 H+ + 2 e-} \]
Le produit est l'\textbf{acide propanoïque}.
\end{corrige}

En résumé, la lampe sans flamme illustre magnifiquement la notion de catalyse, et les tests d'identification (DNPH, Fehling, Tollens, Schiff) te donnent une boîte à outils expérimentale complète pour caractériser les produits d'oxydation des alcools. Retiens la logique en deux temps : \emph{carbonyle ?} puis \emph{réducteur ?}. Nous allons maintenant quitter les réactions d'oxydation pour découvrir une réaction d'un tout autre type, tout aussi caractéristique des alcools : l'\textbf{estérification}, qui transforme un alcool et un acide carboxylique en un composé aux parfums de fruits.

\coursec{Estérification}

Dans la section précédente, nous avons vu comment l'oxydation ménagée transforme un alcool primaire en aldéhyde puis en \cle{acide carboxylique} : l'éthanol donne ainsi l'éthanal, puis l'acide éthanoïque. Or l'acide carboxylique et l'alcool, ces deux familles que nous venons d'étudier séparément, peuvent réagir \emph{l'un avec l'autre}. Cette rencontre donne naissance à une nouvelle famille de composés, les \cle{esters}, dont les odeurs fruitées parfument les bonbons, les arômes alimentaires et même l'essence de banane vendue au marché. Cette réaction, appelée \cle{estérification}, est l'une des plus importantes du programme de Terminale : elle est systématiquement au menu du Baccalauréat.

\courssub{Définition et équation-bilan}

\textbf{Une situation concrète.} Verse dans un tube à essais un peu d'acide éthanoïque (l'acide du vinaigre, à l'odeur piquante) et un peu d'éthanol (l'alcool du vin de palme). Ajoute quelques gouttes d'acide sulfurique concentré et chauffe doucement au bain-marie. Après quelques minutes, refroidis et verse le contenu du tube dans un bécher d'eau salée : une nouvelle odeur se dégage, \emph{fruitée}, rappelant la colle ou certains fruits. Un nouveau composé s'est formé : l'\cle{éthanoate d'éthyle}. Ni l'acide ni l'alcool ne possèdent cette odeur : c'est bien le produit d'une réaction chimique entre les deux.

\begin{definition}[Ester et estérification]
Un \cle{ester} est un composé organique de formule générale \ce{R-COO-R^1}, où $R$ est un atome d'hydrogène ou un groupe alkyle et $R^1$ un groupe alkyle. Les esters possèdent le plus souvent une odeur agréable, fruitée.

L'\cle{estérification} est la réaction entre un \cle{acide carboxylique} \ce{R-COOH} et un \cle{alcool} \ce{R^1-OH}, qui donne un \cle{ester} \ce{R-COO-R^1} et de l'\cle{eau}.
\end{definition}

L'équation-bilan générale s'écrit, avec une \textbf{double flèche} obligatoire :

\[ \ce{R-COOH + R^1-OH <=> R-COO-R^1 + H2O} \]

\begin{attention}[La double flèche est obligatoire]
L'erreur la plus fréquente au Bac est d'écrire l'estérification avec une flèche simple \ce{->}. C'est \textbf{faux} : l'estérification est une réaction \cle{limitée} par un équilibre chimique. L'ester formé réagit avec l'eau pour régénérer l'acide et l'alcool : c'est la réaction inverse, appelée \cle{hydrolyse} de l'ester. On doit donc toujours écrire \ce{<=>}. Une flèche simple coûte des points !
\end{attention}

\textbf{Exemple de référence : synthèse de l'éthanoate d'éthyle.}

\[ \ce{CH3COOH + C2H5OH <=> CH3COOC2H5 + H2O} \]

Conditions expérimentales : quelques gouttes d'acide sulfurique concentré (\ce{H2SO4}, source d'ions \ce{H+} qui jouent le rôle de \cle{catalyseur}) et chauffage modéré (environ \SI{60}{\celsius} au bain-marie, ou chauffage à reflux).

\begin{popfigure}
\begin{center}
\chemfig{CH_3-C(=[1]O)(-[7]OH)} \hspace{0.4em} + \hspace{0.4em}
\chemfig{HO-CH_2-CH_3}
\hspace{0.6em} $\rightleftharpoons$ \hspace{0.6em}
\chemfig{CH_3-C(=[1]O)(-[7]O-CH_2-CH_3)} \hspace{0.4em} + \hspace{0.4em} \ce{H2O}
\end{center}
\legende{Estérification de l'acide éthanoïque par l'éthanol : formation de l'éthanoate d'éthyle et d'eau. La double flèche traduit l'équilibre chimique.}
\end{popfigure}

Remarque le mécanisme global : le groupe \ce{-OH} de l'acide et l'atome d'hydrogène du groupe hydroxyle de l'alcool s'éliminent ensemble pour former l'eau ; les deux fragments restants se lient par l'atome d'oxygène de l'alcool. On retient souvent la phrase mnémotechnique : « \emph{l'acide perd OH, l'alcool perd H} ».

\courssub{Nomenclature des esters}

Le nom d'un ester se construit en deux mots : \textbf{alkanoate d'alkyle}.

\begin{itemize}
\item Le premier mot, \emph{alkanoate}, désigne la partie venant de l'\textbf{acide carboxylique} : on remplace la terminaison « -oïque » par « -oate ». Ainsi, éthanoïque donne \emph{éthanoate}, méthanoïque donne \emph{méthanoate}, propanoïque donne \emph{propanoate}.
\item Le second mot, \emph{alkyle}, désigne le groupe $R'$ venant de l'\textbf{alcool} : méthanol donne \emph{méthyle}, éthanol donne \emph{éthyle}, propan-1-ol donne \emph{propyle}.
\end{itemize}

\begin{exemplebox}[Nommer quelques esters]
\begin{itemize}
\item \ce{CH3COOCH2CH3} : éthanoate d'éthyle (acide éthanoïque + éthanol) ;
\item \ce{HCOOCH3} : méthanoate de méthyle (acide méthanoïque + méthanol) ;
\item \ce{CH3CH2COOCH3} : propanoate de méthyle (acide propanoïque + méthanol) ;
\item \ce{CH3COOCH2CH2CH(CH3)2} : éthanoate d'isoamyle (ou éthanoate de 3-méthylbutyle), l'ester de la banane.
\end{itemize}
\textbf{Piège classique :} c'est le nom de l'\emph{alkyle} qui vient de l'alcool, et le nom de l'\emph{alkanoate} qui vient de l'acide. Ne les inverse pas !
\end{exemplebox}

\courssub{Caractéristiques de l'estérification}

\begin{propriete}[Les trois caractéristiques à connaître par cœur]
L'estérification est une réaction :
\begin{enumerate}
\item \textbf{lente} : l'équilibre n'est atteint qu'au bout de plusieurs heures, voire plusieurs jours à température ambiante ;
\item \textbf{limitée} (ou incomplète) : elle aboutit à un \cle{équilibre chimique} avec la réaction inverse, l'hydrolyse de l'ester ; les réactifs ne disparaissent jamais totalement ;
\item \textbf{athermique} : elle ne dégage ni n'absorbe de chaleur de façon notable ($\Delta_r H \approx 0$). Le chauffage ne sert donc pas à « fournir de l'énergie » à la réaction, mais uniquement à l'\emph{accélérer}.
\end{enumerate}
\end{propriete}

\textbf{Pourquoi un équilibre ?} Dès que l'ester et l'eau apparaissent dans le milieu, ils commencent à réagir entre eux selon la réaction inverse :

\[ \ce{R-COO-R^1 + H2O <=> R-COOH + R^1-OH} \quad \text{(hydrolyse)} \]

Au début, seule l'estérification a lieu (pas encore d'ester ni d'eau). Puis l'hydrolyse s'accélère au fur et à mesure que l'ester s'accumule. Lorsque les deux réactions se déroulent à la même vitesse, les quantités n'évoluent plus : le système a atteint l'\cle{état d'équilibre}. L'avancement se fige sur un \emph{plateau}, comme le montre la courbe suivante.

\begin{popfigure}
\begin{center}
\begin{tikzpicture}
\begin{axis}[
width=0.85\linewidth, height=6cm,
xlabel={Temps $t$ (h)}, ylabel={$n$(ester) (mol)},
xmin=0, xmax=10, ymin=0, ymax=0.12,
xtick={0,2,4,6,8,10},
ytick={0,0.033,0.066,0.10},
yticklabels={0, {0,033}, {0,066}, {0,10}},
grid=both, grid style={popline!40},
legend style={at={(0.97,0.35)}, anchor=east, draw=popline, fill=popcream}
]
\addplot[popcurve, domain=0:10, samples=200] {0.066*(1-exp(-0.6*x))};
\addlegendentry{Avancement réel}
\addplot[dashed, thick, poporange, domain=0:10] {0.066};
\addlegendentry{Limite d'équilibre ($n_{eq} = 0,066$ mol)}
\addplot[dotted, thick, popmuted, domain=0:10] {0.10};
\addlegendentry{Avancement théorique maximal ($n_{max} = 0,10$ mol)}
\end{axis}
\end{tikzpicture}
\end{center}
\legende{Évolution de la quantité d'ester au cours du temps pour un mélange équimolaire de \SI{0,10}{mol} d'acide éthanoïque et de \SI{0,10}{mol} d'éthanol. La réaction est rapide au début, puis ralentit et se stabilise sur un plateau : c'est l'état d'équilibre. La quantité d'ester reste inférieure à la quantité théorique maximale.}
\end{popfigure}

\courssub{Rendement de l'estérification}

Puisque la réaction est limitée, on n'obtient jamais la quantité d'ester prévue par un calcul stœchiométrique « idéal ». Pour quantifier cet écart, on définit le \cle{rendement}.

\begin{definition}[Rendement]
Le \cle{rendement} $\eta$ d'une estérification est le quotient de la quantité d'ester réellement obtenue par la quantité d'ester théoriquement attendue si la réaction était totale :

\[ \eta = \frac{n_{\text{ester obtenu}}}{n_{\text{ester théorique}}} \qquad \text{ou, en pourcentage :} \qquad \eta\,(\%) = \frac{n_{\text{ester obtenu}}}{n_{\text{ester théorique}}} \times 100 \]

La quantité théorique se calcule à partir du \cle{réactif limitant} en supposant la réaction totale.
\end{definition}

\begin{methode}[Calculer un rendement d'estérification]
\begin{enumerate}
\item Écrire l'équation-bilan équilibrée (avec la double flèche).
\item Identifier le \textbf{réactif limitant} en comparant les quantités initiales aux coefficients stœchiométriques (ici, tout est en proportion 1:1).
\item Calculer la quantité \textbf{théorique} d'ester : c'est la quantité du réactif limitant (coefficient 1).
\item Déterminer la quantité \textbf{réellement obtenue} (par dosage, pesée, ou donnée de l'énoncé).
\item Appliquer la formule $\eta = \dfrac{n_{\text{obtenu}}}{n_{\text{théorique}}}$ et exprimer le résultat en pourcentage.
\end{enumerate}
\end{methode}

\begin{exoresolu}[Rendement de la synthèse de l'éthanoate d'éthyle]
On chauffe un mélange de $n_1 = \SI{0,10}{mol}$ d'acide éthanoïque et de $n_2 = \SI{0,10}{mol}$ d'éthanol, en présence de quelques gouttes d'acide sulfurique concentré. À l'équilibre, un dosage montre qu'il s'est formé $n_{eq} = \SI{0,066}{mol}$ d'éthanoate d'éthyle.

Calculer le rendement de cette estérification.
\tcblower
\textbf{Étape 1 : équation-bilan.}
\[ \ce{CH3COOH + C2H5OH <=> CH3COOC2H5 + H2O} \]

\textbf{Étape 2 : réactif limitant.} Les deux réactifs sont introduits en proportions équimolaires ($0,10$ mol chacun) et l'équation est en proportion 1:1 : aucun n'est en excès. Les deux s'épuiseraient en même temps si la réaction était totale.

\textbf{Étape 3 : quantité théorique d'ester.} Si la réaction était totale, on obtiendrait :
\[ n_{\text{théorique}} = n_{\text{limitant}} = \SI{0,10}{mol} \]

\textbf{Étape 4 : quantité obtenue.} $n_{\text{obtenu}} = \SI{0,066}{mol}$ (donnée expérimentale).

\textbf{Étape 5 : rendement.}
\[ \eta = \frac{0,066}{0,10} \times 100 = 66\,\% \]

\textbf{Conclusion :} le rendement est de $\eta \approx 66\,\%$ (souvent arrondi à $67\,\%$). Ce résultat est \textbf{général} : pour un mélange équimolaire acide + alcool \emph{primaire}, le rendement à l'équilibre vaut environ $67\,\%$. Il tombe à environ $60\,\%$ pour un alcool secondaire et à peine $5\,\%$ pour un alcool tertiaire : retiens ces ordres de grandeur, ils sont parfois demandés au Bac.
\end{exoresolu}

\courssub{Comment améliorer une estérification ?}

Deux problèmes distincts se posent au chimiste : la réaction est \emph{lente}, et elle est \emph{limitée}. Les remèdes sont différents, et le Bac adore tester ta capacité à les distinguer.

\textbf{Accélérer la réaction (sans déplacer l'équilibre) :}
\begin{itemize}
\item \textbf{Chauffer} le mélange réactionnel (chauffage à reflux pour ne pas perdre les vapeurs) : la température augmente la vitesse de réaction. Comme la réaction est athermique, le chauffage ne déplace \emph{pas} l'équilibre.
\item \textbf{Ajouter un catalyseur} : quelques gouttes d'acide sulfurique concentré (ions \ce{H+}) ou d'acide paratoluènesulfonique. Le catalyseur accélère à la fois l'estérification et l'hydrolyse : l'équilibre est atteint plus vite, mais \emph{sa position ne change pas}.
\end{itemize}

\textbf{Déplacer l'équilibre pour augmenter le rendement} (loi de Le Chatelier : toute contrainte imposée à un équilibre le déplace dans le sens qui s'y oppose) :
\begin{itemize}
\item \textbf{Introduire l'un des réactifs en excès} (par exemple l'alcool, souvent le moins cher) : l'équilibre se déplace dans le sens de la formation de l'ester. Le rendement, calculé par rapport au réactif limitant, augmente.
\item \textbf{Éliminer l'eau} (ou l'ester) au fur et à mesure de sa formation, par distillation : privée d'un de ses produits, la réaction inverse ne peut plus se faire, et l'équilibre se déplace vers la droite.
\end{itemize}

\begin{attention}[Catalyseur et chauffage ne déplacent pas l'équilibre]
Piège récurrent du Baccalauréat : « l'acide sulfurique augmente-t-il le rendement ? » \textbf{Non !} Le catalyseur et le chauffage permettent seulement d'atteindre l'équilibre \emph{plus rapidement}. Le rendement final reste le même ($\approx 67\,\%$ pour un alcool primaire en proportions équimolaires). Seuls un excès de réactif ou l'élimination d'un produit améliorent réellement le rendement.
\end{attention}

\begin{aretenir}[L'essentiel sur l'estérification]
\begin{itemize}
\item Estérification : \ce{acide carboxylique + alcool <=> ester + eau} — \textbf{double flèche obligatoire}.
\item La réaction est \textbf{lente}, \textbf{limitée} (équilibre avec l'hydrolyse) et \textbf{athermique}.
\item Nomenclature : \textbf{alkanoate d'alkyle} (alkanoate $\leftarrow$ acide ; alkyle $\leftarrow$ alcool).
\item Rendement : $\eta = \dfrac{n_{\text{ester obtenu}}}{n_{\text{ester théorique}}}$ ; environ $67\,\%$ pour un alcool primaire en mélange équimolaire.
\item Pour aller plus vite : chauffer ou catalyser par \ce{H+}. Pour un meilleur rendement : excès d'un réactif ou élimination de l'eau.
\end{itemize}
\end{aretenir}

\begin{savaistu}[L'odeur de la banane... en bouteille]
L'odeur caractéristique de la banane bien mûre est due à un ester : l'\cle{éthanoate d'isoamyle} \ce{CH3COOCH2CH2CH(CH3)2}, formé naturellement dans le fruit lors de sa maturation. Ce même ester est synthétisé industriellement et utilisé comme arôme alimentaire dans les bonbons, les sirops et les boissons gazeuses vendus partout au Cameroun. De même, l'éthanoate d'éthyle sent la colle et les fruits, le butanoate d'éthyle rappelle l'ananas, et l'éthanoate d'octyle évoque l'orange. Les savonneries artisanales qui saponifient l'huile de palme réalisent d'ailleurs la réaction \emph{inverse} de l'estérification : l'hydrolyse (basique) des esters que sont les graisses !
\end{savaistu}

\begin{exercice}[Pour t'entraîner]
On fait réagir \SI{0,20}{mol} d'acide méthanoïque \ce{HCOOH} avec \SI{0,30}{mol} d'éthanol, en présence d'acide sulfurique concentré, jusqu'à l'équilibre. On recueille \SI{0,13}{mol} d'ester.
\begin{enumerate}
\item Écrire l'équation-bilan de la réaction et nommer l'ester formé.
\item Déterminer le réactif limitant et la quantité théorique d'ester.
\item Calculer le rendement de la réaction.
\item Proposer deux moyens d'augmenter ce rendement.
\end{enumerate}
\end{exercice}

\begin{corrige}[Exercice : synthèse du méthanoate d'éthyle]
\textbf{1.} Équation-bilan :
\[ \ce{HCOOH + C2H5OH <=> HCOOC2H5 + H2O} \]
L'ester formé est le \textbf{méthanoate d'éthyle} (alkanoate « méthanoate » venant de l'acide méthanoïque ; alkyle « éthyle » venant de l'éthanol).

\textbf{2.} Proportions stœchiométriques 1:1. L'acide méthanoïque (\SI{0,20}{mol}) est en plus faible quantité : c'est le \textbf{réactif limitant}. Si la réaction était totale : $n_{\text{théorique}} = \SI{0,20}{mol}$ d'ester.

\textbf{3.} Rendement :
\[ \eta = \frac{0,13}{0,20} \times 100 = 65\,\% \]

\textbf{4.} Pour augmenter le rendement : introduire un excès encore plus grand d'éthanol (déplacement d'équilibre vers la droite), ou éliminer l'eau au fur et à mesure de sa formation par distillation. (Chauffer ou ajouter davantage de catalyseur ne ferait qu'atteindre l'équilibre plus vite, sans améliorer le rendement.)
\end{corrige}

\coursec{Préparation des alcools}

Dans la section précédente, nous avons vu comment un alcool peut être \emph{transformé} : en réagissant avec un acide carboxylique, il donne un ester et de l'eau, dans une réaction lente et limitée. Une question se pose alors naturellement : d'où viennent les alcools eux-mêmes ? Comment l'industrie produit-elle les millions de tonnes d'éthanol utilisées chaque année, et comment le brasseur de \og bil-bil \fg{} ou le récolteur de vin de palme obtient-il de l'alcool sans aucune usine ? Il existe deux grandes voies de préparation des alcools : une voie \cle{chimique industrielle}, l'\cle{hydratation des alcènes}, et une voie \cle{biologique}, la \cle{fermentation alcoolique}. Étudions-les l'une après l'autre.

\courssub{Hydratation d'un alcène}

\textbf{Principe : addition d'eau sur une double liaison}\par

Tu as appris en classe de Première que les alcènes possèdent une \cle{double liaison} \ce{C=C}, riche en électrons et donc très réactive. L'idée de l'hydratation est simple : on \og casse \fg{} cette double liaison et on y \og accroche \fg{} une molécule d'eau. L'un des deux atomes de carbone reçoit un atome d'hydrogène \ce{H}, l'autre reçoit le groupe hydroxyle \ce{-OH}. La double liaison devient une simple liaison, et l'alcène devient un \textbf{alcool}.

\begin{definition}[Hydratation d'un alcène]
L'\cle{hydratation} d'un alcène est la réaction d'\textbf{addition} d'une molécule d'eau sur la double liaison \ce{C=C}, en présence d'un catalyseur acide (acide sulfurique \ce{H2SO4} au laboratoire, acide phosphorique \ce{H3PO4} supporté en industrie) et à chaud. Elle transforme un alcène en alcool :
\[ \ce{C_nH_{2n} + H2O ->[H+,\ \text{chaud}] C_nH_{2n+1}OH} \]
\end{definition}

L'exemple le plus important est l'hydratation de l'éthylène (éthène), qui fournit l'éthanol \og industriel \fg{} utilisé comme solvant, désinfectant ou intermédiaire de synthèse :
\[ \ce{CH2=CH2 + H2O ->[{H3PO4, $\approx$ \SI{300}{\celsius}, \SI{70}{bar}}] CH3-CH2OH} \]

\begin{exemplebox}[L'éthanol industriel]
L'éthanol de synthèse est obtenu à grande échelle par hydratation de l'éthène issu du pétrole (vapeur d'eau sur catalyseur acide phosphorique supporté, vers \SI{300}{\celsius} sous environ 70 bars). Cet alcool, très pur, sert à fabriquer des médicaments, des parfums et des solvants. Il est en général \textbf{dénaturé} (rendu imbuvable par addition de substances amères ou toxiques) pour les usages techniques, ce qui permet de l'exonérer de la taxe sur les boissons alcooliques.
\end{exemplebox}

\textbf{Le problème des alcènes dissymétriques : règle de Markovnikov}\par

Avec l'éthène \ce{CH2=CH2}, aucune hésitation : les deux atomes de carbone de la double liaison sont identiques, il ne peut se former qu'un seul alcool. Mais avec un alcène \cle{dissymétrique} comme le propène \ce{CH2=CH-CH3}, les deux carbones de la double liaison ne portent pas le même nombre d'atomes d'hydrogène. Deux produits sont alors \emph{possibles} :

\begin{itemize}
\item si \ce{H} se fixe sur le carbone 1 et \ce{-OH} sur le carbone 2 : on obtient le \textbf{propan-2-ol} (alcool secondaire) ;
\item si \ce{H} se fixe sur le carbone 2 et \ce{-OH} sur le carbone 1 : on obtient le \textbf{propan-1-ol} (alcool primaire).
\end{itemize}

L'expérience montre que ces deux produits ne se forment \textbf{pas en quantités égales}. Le chimiste russe Vladimir Markovnikov a énoncé en 1870 la règle qui permet de prévoir le produit majoritaire.

\begin{propriete}[Règle de Markovnikov]
Lors de l'addition d'une molécule \ce{H-OH} sur un alcène dissymétrique, l'atome d'\textbf{hydrogène se fixe préférentiellement sur l'atome de carbone de la double liaison le plus hydrogéné} (celui qui porte déjà le plus d'atomes \ce{H}). Le groupe \ce{-OH} se fixe donc sur l'autre carbone, le \textbf{plus substitué}. L'alcool majoritaire est l'alcool \textbf{le plus substitué} (de classe la plus élevée).
\end{propriete}

Pour le propène : le carbone \ce{CH2} porte 2 hydrogènes, le carbone \ce{CH} n'en porte qu'un. L'hydrogène va donc sur le \ce{CH2}, et le \ce{-OH} sur le carbone central :
\[ \ce{CH2=CH-CH3 + H2O ->[H+] CH3-CHOH-CH3} \quad \text{(majoritaire, } \approx 80\ \%\text{)} \]
\[ \ce{CH2=CH-CH3 + H2O ->[H+] CH3-CH2-CH2OH} \quad \text{(minoritaire)} \]

\begin{popfigure}
\begin{center}
\begin{tikzpicture}[scale=0.9]
% Reactants
\node (propene) at (0,0) {\chemfig{CH_2=CH-CH_3}};
\node[font=\small, below=2pt] at (propene.south) {propène};
\node (plus) at (2.5,0) {\Large $+$};
\node (water) at (3.8,0) {\chemfig{H_2O}};
% Arrow
\draw[-{Stealth[length=3mm]}, very thick, popblue] (5.0,0) -- (7.0,0);
\node[font=\small, above=2pt, popblue] at (6.0,0.3) {\ce{H2SO4}, chaud};
% Products
\node (major) at (9.5,1.2) {\chemfig{CH_3-CH(-[2]OH)-CH_3}};
\node[font=\small, popgreen, below=2pt] at (9.5,0.6) {\textbf{propan-2-ol} (majoritaire $\approx 80\ \%$)};
\node[font=\small, popmuted] at (8.5,-0.3) {$+$};
\node (minor) at (9.5,-1.5) {\chemfig{CH_3-CH_2-CH_2-OH}};
\node[font=\small, poporange, below=2pt] at (9.5,-2.1) {\textbf{propan-1-ol} (minoritaire $\approx 20\ \%$)};
\end{tikzpicture}
\end{center}
\legende{Hydratation acide du propène : le \ce{H} se fixe sur le carbone le plus hydrogéné (\ce{CH2}), le \ce{-OH} sur le carbone le plus substitué. Le propan-2-ol (alcool secondaire) est majoritaire, conformément à la règle de Markovnikov.}
\end{popfigure}

\begin{methode}[Appliquer la règle de Markovnikov]
Pour prévoir l'alcool majoritaire issu de l'hydratation d'un alcène dissymétrique :
\begin{enumerate}
\item \textbf{Repérer} les deux atomes de carbone de la double liaison \ce{C=C}.
\item \textbf{Compter} les atomes d'hydrogène portés par chacun de ces deux carbones.
\item \textbf{Placer} l'atome \ce{H} de l'eau sur le carbone \emph{le plus hydrogéné}.
\item \textbf{Placer} le groupe \ce{-OH} sur l'autre carbone (le plus substitué) : c'est l'alcool \textbf{majoritaire}.
\item Écrire l'autre alcool comme produit \textbf{minoritaire} si l'énoncé le demande.
\end{enumerate}
\end{methode}

\begin{attention}[Erreurs fréquentes]
\begin{itemize}
\item Ne pas inverser la règle : c'est bien \ce{H} qui va sur le carbone \emph{le plus hydrogéné}, et \ce{-OH} sur le carbone \emph{le moins hydrogéné}. Une phrase pour s'en souvenir : \og le riche s'enrichit \fg{} (le carbone riche en \ce{H} reçoit encore un \ce{H}).
\item L'hydratation du propène donne majoritairement un alcool \textbf{secondaire}, pas un alcool primaire. L'hydratation d'un alcène ne permet donc pas d'obtenir n'importe quel alcool : le produit majoritaire est imposé par la structure de l'alcène.
\item Ne pas oublier le catalyseur acide (\ce{H2SO4} ou \ce{H3PO4}) au-dessus de la flèche : sans lui, la réaction est extrêmement lente.
\item Attention à ne pas confondre \textbf{hydratation} (addition d'eau, donne un alcool) et \textbf{déshydratation} (élimination d'eau, donne un alcène), qui sont des réactions inverses selon les conditions.
\end{itemize}
\end{attention}

\courssub{Fermentation alcoolique}

\textbf{Une transformation biologique : l'intuition}\par

La seconde grande voie de préparation des alcools ne doit rien au pétrole : elle est aussi vieille que l'humanité. Quand on laisse du jus de bissap sucré, de la sève de palmier ou du moût de maïs reposer à l'abri de l'air, le liquide se met à \emph{bouillonner}, devient trouble et acquiert un goût piquant et alcoolisé. Des micro-organismes invisibles, les \cle{levures}, ont transformé le sucre en alcool. C'est la \cle{fermentation alcoolique}.

\begin{definition}[Fermentation alcoolique]
La \cle{fermentation alcoolique} est la dégradation \textbf{anaérobie} (c'est-à-dire \emph{en l'absence de dioxygène}) du glucose en \textbf{éthanol} et \textbf{dioxyde de carbone}, sous l'action d'\textbf{enzymes} produites par des levures (champignons microscopiques, par exemple \emph{Saccharomyces cerevisiae}).
\[ \ce{C6H12O6 ->[\text{enzymes des levures}] 2 C2H5OH + 2 CO2} \]
Conditions : absence d'air (anaérobiose stricte), température voisine de \SI{30}{\celsius}, milieu aqueux sucré, pH légèrement acide.
\end{definition}

L'équation-bilan mérite un commentaire attentif. Une mole de glucose (\ce{C6H12O6}) fournit \textbf{deux} moles d'éthanol et \textbf{deux} moles de dioxyde de carbone. Vérifions l'équilibre : à gauche, 6 C, 12 H, 6 O ; à droite, $2 \times 2 = 4$ C dans l'éthanol plus 2 C dans le \ce{CO2}, soit 6 C ; $2 \times 6 = 12$ H ; $2 \times 1 + 2 \times 2 = 6$ O. L'équation est bien équilibrée.

\begin{attention}[Le piège classique du Baccalauréat]
L'erreur la plus fréquente est d'écrire \ce{C6H12O6 -> C2H5OH + CO2} en oubliant les \textbf{coefficients 2}. Cette équation n'est \emph{pas équilibrée} (6 C à gauche, 3 C à droite). Retiens : \textbf{1 glucose $\rightarrow$ 2 éthanol + 2 \ce{CO2}}. Le dégagement gazeux abondant (les bulles du vin de palme !) est une preuve visible de la formation de \ce{CO2}.
\end{attention}

\textbf{Conditions et limites de la fermentation}\par

Trois points sont essentiels pour comprendre (et réussir) une fermentation alcoolique.

\textbf{1. L'anaérobiose est indispensable.} En présence de dioxygène, l'éthanol formé est oxydé par d'autres micro-organismes (bactéries acétiques, genre \emph{Acetobacter}) en \textbf{acide éthanoïque} \ce{CH3COOH} : c'est le \cle{vinaigrage}. Une boisson alcoolisée laissée à l'air libre tourne au vinaigre :
\[ \ce{CH3-CH2OH + O2 ->[\text{bactéries acétiques}] CH3-COOH + H2O} \]
C'est pourquoi le brasseur ferme soigneusement ses récipients (ou utilise une cheminée à eau) et pourquoi le vin de palme devient acide en un ou deux jours s'il n'est pas bu rapidement.

\textbf{2. La température doit rester modérée} (environ \SIrange{25}{35}{\celsius}). En dessous, les levures sont paresseuses ; au-dessus de \SI{40}{\celsius} environ, elles meurent (dénaturation des enzymes) et la fermentation s'arrête.

\textbf{3. La fermentation s'arrête d'elle-même vers 12 à 15 \% vol. d'alcool.} L'éthanol est \emph{toxique pour les levures elles-mêmes} : quand sa concentration atteint environ 12 à 15 degrés, les levures meurent et la transformation cesse, même s'il reste du sucre. C'est pourquoi aucune boisson \emph{fermentée seule} (vin, bière, vin de palme) ne dépasse cette teneur. Pour obtenir des boissons plus fortes (eau-de-vie, \og odontol \fg{} distillé), il faut \textbf{distiller} le produit de fermentation : l'éthanol, plus volatil que l'eau ($T_{\text{éb}} = \SI{78,4}{\celsius}$ contre \SI{100}{\celsius}), s'évapore d'abord et se concentre dans le distillat.

\begin{popfigure}
\begin{center}
\begin{tikzpicture}[scale=0.85]
% Cuve de fermentation
\draw[very thick, popdark] (0,0) -- (0,3.5) (4,0) -- (4,3.5);
\draw[very thick, popdark] (0,0) -- (4,0);
% Liquide
\fill[popgold!30] (0.1,0.1) rectangle (3.9,2.5);
\draw[popgold!70!black, thick] (0.1,2.5) -- (3.9,2.5);
\node[font=\small, popdark] at (2,1.4) {Glucose + Levures};
\node[font=\small, popdark] at (2,1.0) {($\approx \SI{30}{\celsius}$, anaérobiose)};
% Bulles CO2 dans la cuve
\foreach \x/\y/\r in {0.6/1.8/0.07, 1.2/2.1/0.05, 1.8/1.7/0.06, 2.5/2.0/0.07, 3.1/1.8/0.05, 3.3/2.2/0.06, 1.0/1.4/0.05, 2.8/1.3/0.05}
  \draw[popgold!80!black] (\x,\y) circle (\r);
% Couvercle et tuyau de dégagement
\draw[very thick, popdark] (0,3.5) -- (4,3.5);
\draw[very thick, popdark] (1.8,3.5) -- (1.8,4.3) -- (6.5,4.3) -- (6.5,2.8);
\node[font=\small, popblue, above] at (4.2,4.3) {\ce{CO2}};
% Bécher de barbotage (eau de chaux)
\draw[very thick, popdark] (5.5,0.6) -- (5.5,2.8) (7.5,0.6) -- (7.5,2.8);
\draw[very thick, popdark] (5.5,0.6) -- (7.5,0.6);
\fill[popblue!15] (5.55,0.65) rectangle (7.45,2.1);
\node[font=\small, popblue!70!black] at (6.5,1.1) {Eau de chaux};
\foreach \y in {2.4,2.15,1.9,1.65}
  \draw[popblue!60!black] (6.5,\y) circle (0.06);
\node[font=\small, popblue!70!black, right=2pt] at (7.6,1.4) {Trouble $\Rightarrow$ \ce{CO2}};
% Étiquettes
\node[font=\small\bfseries, popdark, below] at (2,-0.2) {Cuve de fermentation};
\node[font=\small, popgreen!60!black] at (2,2.9) {$\rightarrow$ Éthanol + \ce{CO2}\uparrow};
\end{tikzpicture}
\end{center}
\legende{Mise en évidence de la fermentation alcoolique : dans une cuve fermée (anaérobiose), les levures transforment le glucose en éthanol et dioxyde de carbone. Le \ce{CO2} dégagé barbote dans l'eau de chaux, qui se trouble (formation de \ce{CaCO3}) : c'est le test de reconnaissance du dioxyde de carbone.}
\end{popfigure}

\begin{exoresolu}[Fermentation de 180 g de glucose]
\textbf{Énoncé.} Un brasseur fait fermenter \SI{180}{g} de glucose pur en anaérobiose, à \SI{30}{\celsius}, jusqu'à épuisement du sucre. On donne $M(\ce{C}) = \SI{12}{g.mol^{-1}}$, $M(\ce{H}) = \SI{1}{g.mol^{-1}}$, $M(\ce{O}) = \SI{16}{g.mol^{-1}}$ et le volume molaire dans les CNTP $V_m = \SI{22,4}{L.mol^{-1}}$.
\begin{enumerate}
\item Calculer la quantité de glucose fermenté.
\item En déduire la masse d'éthanol obtenue.
\item Calculer le volume de dioxyde de carbone dégagé (CNTP).
\end{enumerate}
\tcblower
\textbf{Solution détaillée.}
\begin{enumerate}
\item Masse molaire du glucose : $M(\ce{C6H12O6}) = 6 \times 12 + 12 \times 1 + 6 \times 16 = \SI{180}{g.mol^{-1}}$.
\[ n(\ce{C6H12O6}) = \frac{m}{M} = \frac{180}{180} = \SI{1,0}{mol}. \]
\item D'après l'équation $\ce{C6H12O6 -> 2 C2H5OH + 2 CO2}$, une mole de glucose donne \textbf{deux} moles d'éthanol :
\[ n(\ce{C2H5OH}) = 2 \times 1{,}0 = \SI{2,0}{mol}. \]
Avec $M(\ce{C2H5OH}) = 2 \times 12 + 6 \times 1 + 16 = \SI{46}{g.mol^{-1}}$ :
\[ m(\ce{C2H5OH}) = n \times M = 2{,}0 \times 46 = \boxed{\SI{92}{g}}. \]
\item De même, $n(\ce{CO2}) = \SI{2,0}{mol}$, donc :
\[ V(\ce{CO2}) = n \times V_m = 2{,}0 \times 22{,}4 = \boxed{\SI{44,8}{L}}. \]
\end{enumerate}
\textbf{Remarque.} Près de 45 litres de gaz pour 180 g de sucre : on comprend pourquoi une cuve de fermentation mal fermée peut éclater, et pourquoi le vin de palme pétille autant !
\end{exoresolu}

\begin{experience}[Mise en évidence de la fermentation alcoolique au laboratoire]
\textbf{Dispositif et protocole.} Dans un erlenmeyer, dissoudre \SI{20}{g} de glucose (ou de saccharose) dans \SI{100}{mL} d'eau tiède (\SI{\approx 35}{\celsius}), ajouter \SI{5}{g} de levure de boulanger, puis boucher avec un bouchon traversé par un tube à dégagement plongeant dans un bécher contenant \SI{50}{mL} d'\textbf{eau de chaux} (\ce{Ca(OH)2} aqueuse). Placer l'ensemble dans une étuve ou près d'une source de chaleur douce à \SI{\approx 30}{\celsius}.

\textbf{Observations.} Au bout de 15 à 30 minutes, des bulles se forment dans l'erlenmeyer et barbotent dans l'eau de chaux, qui \textbf{se trouble} progressivement (formation d'un précipité blanc de carbonate de calcium \ce{CaCO3}). Une odeur caractéristique d'alcool (éthanol) se dégage du mélange réactionnel.

\textbf{Interprétation.} Le trouble de l'eau de chaux prouve le dégagement de \ce{CO2} (réaction $\ce{CO2 + Ca(OH)2 -> CaCO3 v + H2O}$) ; l'odeur et les tests usuels (réaction au dichromate, odeur) confirment la formation d'éthanol. La réaction globale est bien $\ce{C6H12O6 -> 2 C2H5OH + 2 CO2}$.
\end{experience}

\textbf{Applications au Cameroun}\par

\begin{savaistu}[Le secret des bulles du vin de palme]
Le \textbf{vin de palme}, récolté comme une sève sucrée et limpide à la cime du palmier raphia ou du palmier à huile, devient en quelques heures \textbf{trouble et pétillant}. Ce trouble et ces bulles ne sont pas un défaut : ce sont la signature du \ce{CO2} dégagé par la \textbf{fermentation naturelle} de la sève, déclenchée par des levures sauvages présentes dans l'air et sur l'écorce. Plus le vin \og travaille \fg{}, plus il est alcoolisé... et plus il risque de vinaigrer s'il est exposé à l'air ! Le même principe gouverne la fabrication du \og bil-bil \fg{} (bière de mil ou de maïs du Nord-Cameroun) et de l'odontol (eau-de-vie artisanale obtenue par distillation du vin de palme ou du bil-bil). À plus grande échelle, la fermentation de canne à sucre, de maïs ou de manioc produit le \textbf{bioéthanol}, un carburant renouvelable mélangé à l'essence (supercarburant sans plomb) dans de nombreux pays : la chimie des levures est aussi une affaire d'énergie et de développement durable.
\end{savaistu}

\begin{aretenir}[L'essentiel de la section]
\begin{itemize}
\item \textbf{Hydratation d'un alcène} : $\ce{C_nH_{2n} + H2O ->[H+,\ \text{chaud}] C_nH_{2n+1}OH}$. Exemple industriel : $\ce{CH2=CH2 + H2O ->[\text{H3PO4, 300 °C, 70 bar}] CH3-CH2OH}$.
\item \textbf{Règle de Markovnikov} : sur un alcène dissymétrique, \ce{H} se fixe sur le carbone de \ce{C=C} \emph{le plus hydrogéné} ; l'alcool majoritaire est le plus substitué. Propène $+$ eau $\rightarrow$ propan-2-ol (majoritaire) $+$ propan-1-ol (minoritaire).
\item \textbf{Fermentation alcoolique} : dégradation \emph{anaérobie} du glucose par les levures (\emph{Saccharomyces cerevisiae}), vers \SI{30}{\celsius}, pH acide :
\[ \ce{C6H12O6 -> 2 C2H5OH + 2 CO2} \quad \text{(ne pas oublier les coefficients 2 !)} \]
\item En présence d'air, l'éthanol s'oxyde en acide éthanoïque par action de bactéries acétiques : c'est le \textbf{vinaigrage}.
\item La fermentation s'arrête vers 12--15 \% vol. (toxicité de l'éthanol pour les levures) ; la \textbf{distillation} permet de concentrer l'alcool (odontol, eau-de-vie).
\item Bilan molaire type : 1 mol de glucose (\SI{180}{g}) $\rightarrow$ 2 mol d'éthanol (\SI{92}{g}) $+$ 2 mol de \ce{CO2} (\SI{44,8}{L} aux CNTP).
\end{itemize}
\end{aretenir}

\begin{exercice}[Alcène et fermentation]
\textbf{1.} On hydrate le but-1-ène \ce{CH2=CH-CH2-CH3} en présence d'acide sulfurique concentré et chaud. Écrire les formules semi-développées des deux alcools susceptibles de se former, les nommer (nomenclature IUPAC), et préciser lequel est majoritaire en justifiant par la règle de Markovnikov.

\textbf{2.} Une cuve de fermentation a fourni, après réaction complète, \SI{4,6}{kg} d'éthanol. Quelle masse de glucose a été consommée ? Quel volume de \ce{CO2} (aux CNTP, $V_m = \SI{22,4}{L.mol^{-1}}$) s'est dégagé ?
\end{exercice}

\begin{corrige}[Exercice : Alcène et fermentation]
\textbf{1.} Les deux carbones de la double liaison sont \ce{CH2} (carbone 1, 2 atomes \ce{H}) et \ce{CH} (carbone 2, 1 atome \ce{H}). Selon Markovnikov, \ce{H} se fixe sur le carbone 1 (\ce{CH2}) et \ce{-OH} sur le carbone 2 :
\begin{itemize}
\item \textbf{butan-2-ol} \chemfig{CH_3-CH(-[2]OH)-CH_2-CH_3} : alcool secondaire, \textbf{majoritaire} ;
\item \textbf{butan-1-ol} \chemfig{CH_3-CH_2-CH_2-CH_2-OH} : alcool primaire, minoritaire.
\end{itemize}

\textbf{2.} Données : $M(\ce{C2H5OH}) = \SI{46}{g.mol^{-1}}$ ; $M(\ce{C6H12O6}) = \SI{180}{g.mol^{-1}}$.
\[ n(\ce{C2H5OH}) = \frac{m}{M} = \frac{4600}{46} = \SI{100}{mol}. \]
D'après l'équation-bilan $\ce{C6H12O6 -> 2 C2H5OH + 2 CO2}$ :
\[ n(\text{glucose}) = \frac{n(\ce{C2H5OH})}{2} = \frac{100}{2} = \SI{50}{mol}. \]
\[ m(\text{glucose}) = n \times M = 50 \times 180 = \SI{9000}{g} = \boxed{\SI{9,0}{kg}}. \]
\[ n(\ce{CO2}) = n(\ce{C2H5OH}) = \SI{100}{mol} \quad (\text{coefficients stœchiométriques identiques}). \]
\[ V(\ce{CO2}) = n \times V_m = 100 \times 22{,}4 = \SI{2240}{L} = \boxed{\SI{2,24}{m^3}}. \]
\end{corrige}

\newpage

\coursec{Travaux pratiques}

\courssub{Objectifs du TP}

À l'issue de cette séance de travaux pratiques, tu dois être capable de :
\begin{itemize}
\item réaliser l'\cle{oxydation ménagée} des trois classes d'alcools (primaire, secondaire, tertiaire) par une solution acidifiée de dichromate de potassium ;
\item mettre en évidence expérimentalement l'influence de la \cle{classe de l'alcool} sur son aptitude à s'oxyder ;
\item identifier les produits formés (aldéhyde, cétone) à l'aide de \cle{tests caractéristiques} : DNPH, liqueur de Fehling, réactif de Schiff ;
\item observer l'oxydation catalytique de l'éthanol par l'expérience de la \cle{lampe sans flamme} ;
\item rédiger un compte rendu scientifique rigoureux : observations, équations-bilans, conclusion.
\end{itemize}

\begin{savaistu}[Pourquoi ce TP est-il important ?]
L'oxydation ménagée des alcools est au cœur de nombreuses situations camerounaises : la transformation du vin de palme en vinaigre (oxydation de l'éthanol en acide éthanoïque par l'oxygène de l'air, catalysée par des bactéries acétiques), le contrôle d'alcoolémie (les premiers alcootests utilisaient justement le virage orange $\rightarrow$ vert du dichromate de potassium !), ou encore la production d'acétone dans l'industrie. Comprendre ce TP, c'est comprendre la chimie qui t'entoure, du laboratoire de la SONARA aux brassins de bili-bili.
\end{savaistu}

\courssub{Matériel et produits}

\begin{center}
\begin{tabularx}{\linewidth}{|l|X|}
\hline
\rowcolor{popblue!40} \textbf{Matériel} & \textbf{Produits} \\
\hline
\rowcolor{popblueL} 4 tubes à essai avec portoir & Éthanol (alcool primaire) \\
\hline
Bécher de \SI{250}{mL} (bain-marie) & Propan-2-ol (alcool secondaire) \\
\hline
\rowcolor{popblueL} Pipettes graduées de \SI{1}{mL} et \SI{5}{mL} & 2-méthylpropan-2-ol (alcool tertiaire) \\
\hline
Fil de cuivre (environ \SI{15}{cm}) & Solution de dichromate de potassium \ce{K2Cr2O7} acidifiée à l'acide sulfurique dilué \\
\hline
\rowcolor{popblueL} Pince en bois ou en métal & Solution de DNPH (2,4-dinitrophénylhydrazine) \\
\hline
Bec Bunsen (ou lampe à alcool) & Liqueur de Fehling (ou réactif de Schiff) \\
\hline
\rowcolor{popblueL} Thermomètre & Eau distillée \\
\hline
\end{tabularx}
\end{center}

\textbf{Alternatives avec le matériel local :}
\begin{itemize}
\item Si le \cle{dichromate de potassium} manque, utiliser une solution diluée de \cle{permanganate de potassium} \ce{KMnO4} acidifiée (quelques cristaux dans l'eau distillée + quelques gouttes d'acide sulfurique dilué) : on observe alors la décoloration du violet.
\item Si le propan-2-ol est indisponible, l'alcool à brûler du commerce (souvent à base d'éthanol dénaturé) peut servir de témoin, mais la comparaison des classes sera alors limitée.
\item À défaut de bain-marie thermostaté, un bécher d'eau chauffée à environ \SI{60}{\celsius} sur bec Bunsen convient parfaitement.
\item Le fil de cuivre peut être récupéré sur un vieux câble électrique dénudé.
\end{itemize}

\courssub{Sécurité}

\begin{attention}[Consignes de sécurité impératives]
\begin{itemize}
\item \textbf{Port obligatoire} de la blouse, des lunettes de protection et, si possible, de gants.
\item \textbf{Produits inflammables} (pictogramme : flamme) : les alcools s'enflamment facilement. Ne jamais chauffer directement un tube contenant un alcool à la flamme : utiliser exclusivement le \cle{bain-marie}.
\item \textbf{Dichromate de potassium} (pictogrammes : toxique, cancérogène, mutagène, reprotoxique, corrosif, danger pour l'environnement) : éviter tout contact avec la peau ; en cas de contact, rincer abondamment à l'eau et prévenir l'enseignant. Ne jamais verser les résidus à l'évier : les récupérer dans un bidon de déchets chimiques spécifiques (métaux lourds).
\item \textbf{Acide sulfurique} (pictogramme : corrosif) : verser toujours l'acide dans l'eau, jamais l'inverse.
\item \textbf{Liqueur de Fehling} : solution très basique, corrosive ; manipuler avec précaution.
\item Tenir les tubes à essai avec une pince en bois, ouverture dirigée \textbf{loin de tout visage}.
\item L'éthanal dégagé lors de la lampe sans flamme a une odeur piquante : ne pas inhaler directement, travailler près d'une aération ou sous hotte.
\end{itemize}
\end{attention}

\courssub{Schéma du dispositif}

\begin{popfigure}
\begin{center}
\begin{tikzpicture}[scale=0.95, every node/.style={font=\small}]
% Bain-marie
\draw[thick, popdark] (-0.2,0) -- (-0.2,2.2) -- (6.2,2.2) -- (6.2,0);
\draw[thick, popdark] (-0.4,0) -- (6.4,0);
\fill[popblue!25] (-0.15,0.05) rectangle (6.15,1.6);
\draw[popblue, thick] (-0.15,1.6) -- (6.15,1.6);
\node[popblue, anchor=west] at (4.6,1.2) {eau à \SI{60}{\celsius}};
% Tube 1
\draw[thick, popdark] (0.7,1.0) -- (0.7,3.6) (1.5,1.0) -- (1.5,3.6);
\draw[thick, popdark] (0.7,1.0) arc (180:360:0.4);
\fill[popgreen!45] (0.73,0.65) -- (0.73,1.9) -- (1.47,1.9) -- (1.47,0.65) arc (360:180:0.37);
\node[popdark] at (1.1,4.0) {Tube 1 : éthanol};
\node[popgreen!60!black] at (1.1,1.3) {vert};
% Tube 2
\draw[thick, popdark] (2.7,1.0) -- (2.7,3.6) (3.5,1.0) -- (3.5,3.6);
\draw[thick, popdark] (2.7,1.0) arc (180:360:0.4);
\fill[popgreen!45] (2.73,0.65) -- (2.73,1.9) -- (3.47,1.9) -- (3.47,0.65) arc (360:180:0.37);
\node[popdark] at (3.1,4.0) {Tube 2 : propan-2-ol};
\node[popgreen!60!black] at (3.1,1.3) {vert};
% Tube 3
\draw[thick, popdark] (4.7,1.0) -- (4.7,3.6) (5.5,1.0) -- (5.5,3.6);
\draw[thick, popdark] (4.7,1.0) arc (180:360:0.4);
\fill[poporange!55] (4.73,0.65) -- (4.73,1.9) -- (5.47,1.9) -- (5.47,0.65) arc (360:180:0.37);
\node[popdark] at (5.1,4.0) {Tube 3 : tert-butanol};
\node[poporange!80!black] at (5.1,1.3) {orange};
% Flamme
\draw[poporange, thick, decorate, decoration={coil, aspect=0.5, segment length=4pt}] (2.6,-0.25) -- (3.4,-0.25);
\node[poporange!80!black] at (3.0,-0.7) {bec Bunsen};
% Lampe sans flamme (à droite)
\begin{scope}[xshift=9.2cm]
\draw[thick, popdark] (-1.0,0) -- (-1.0,2.6) -- (1.4,2.6) -- (1.4,0);
\draw[thick, popdark] (-1.2,0) -- (1.6,0);
\fill[popgold!35] (-0.95,0.05) rectangle (1.35,1.1);
\draw[popgold!80!black, thick] (-0.95,1.1) -- (1.35,1.1);
\node[popgold!60!black] at (0.2,0.55) {éthanol};
% vapeurs
\foreach \x in {-0.5,0.2,0.9}{\draw[popmuted, ->] (\x,1.2) -- (\x,1.9);}
\node[popmuted] at (0.2,2.3) {vapeurs d'éthanol};
% fil de cuivre
\draw[very thick, poppink] (0.2,4.2) -- (0.2,2.0);
\draw[very thick, poppink] (0.2,2.0) arc (90:-90:0.35);
\draw[very thick, poppink] (0.2,1.3) arc (90:270:0.35);
\fill[poppink!60, opacity=0.5] (0.2,1.65) circle (0.75);
\node[poppink!70!black, anchor=west] at (1.1,1.7) {spirale de cuivre};
\node[poppink!70!black, anchor=west] at (1.1,1.25) {incandescente};
\node[popdark] at (0.2,-0.5) {Lampe sans flamme};
\end{scope}
\end{tikzpicture}
\end{center}
\legende{À gauche : bain-marie contenant les trois tubes d'oxydation ménagée (éthanol, propan-2-ol, 2-méthylpropan-2-ol). À droite : dispositif de la lampe sans flamme (spirale de cuivre chauffée au rouge, maintenue au-dessus des vapeurs d'éthanol).}
\end{popfigure}

\courssub{Protocole}

\textbf{Partie A : Oxydation ménagée par le dichromate de potassium}

\begin{enumerate}
\item Numéroter trois tubes à essai 1, 2 et 3. Introduire dans chacun environ \SI{1}{mL} (soit 20 gouttes) respectivement d'éthanol (tube 1), de propan-2-ol (tube 2) et de 2-méthylpropan-2-ol (tube 3).
\item Ajouter dans chaque tube \SI{2}{mL} de solution acidifiée de dichromate de potassium. Noter la couleur initiale : orange (due à l'ion \ce{Cr2O7^2-}).
\item Placer les trois tubes dans le bain-marie maintenu à environ \SI{60}{\celsius} pendant 5 à 10 minutes. Agiter doucement de temps en temps.
\item Observer et noter les changements de couleur dans chaque tube.
\end{enumerate}

\textbf{Partie B : Identification des produits formés}

\begin{enumerate}
\setcounter{enumi}{4}
\item \textbf{Test à la DNPH} : prélever quelques gouttes du contenu des tubes 1 et 2, les déposer dans deux tubes propres, ajouter \SI{1}{mL} de solution de DNPH. Observer la formation éventuelle d'un précipité.
\item \textbf{Test à la liqueur de Fehling} : dans un tube, mélanger \SI{1}{mL} du contenu du tube 1 avec \SI{2}{mL} de liqueur de Fehling ; chauffer doucement au bain-marie. Observer. Refaire le même test avec le contenu du tube 2.
\item \textbf{Variante} : si la liqueur de Fehling manque, utiliser le réactif de Schiff (à froid) sur le contenu du tube 1 : une coloration rose-violacé apparaît en présence d'un aldéhyde.
\end{enumerate}

\textbf{Partie C : Lampe sans flamme (oxydation catalytique)}

\begin{enumerate}
\setcounter{enumi}{7}
\item Verser environ \SI{2}{cm} d'éthanol dans un bécher propre et sec.
\item À l'aide d'une pince, chauffer la spirale de fil de cuivre au rouge dans la flamme du bec Bunsen : le cuivre se recouvre d'oxyde de cuivre(II) \ce{CuO} noir.
\item Plonger immédiatement la spirale encore chaude juste au-dessus de la surface de l'éthanol, dans les vapeurs. Observer : la spirale reste incandescente et redevient brillante (cuivre régénéré).
\item Répéter l'opération plusieurs fois, puis sentir prudemment (en éventant) : odeur piquante caractéristique de l'éthanal.
\end{enumerate}

\courssub{Observations et résultats attendus}

\begin{center}
\begin{tabularx}{\linewidth}{|l|X|X|X|}
\hline
\rowcolor{poporange!50} \textbf{Tube} & \textbf{Alcool (classe)} & \textbf{Couleur après chauffage} & \textbf{Conclusion} \\
\hline
1 & \chemfig{CH_3-CH_2-OH} (primaire) & Orange $\rightarrow$ \textbf{vert} en 3--4 min & Oxydation : aldéhyde puis acide carboxylique \\
\hline
2 & \chemfig{H_3C-CH(-[2]OH)-CH_3} (secondaire) & Orange $\rightarrow$ \textbf{vert} en 4--5 min & Oxydation : cétone \\
\hline
3 & \chemfig{H_3C-C(-[2]CH_3)(-[6]OH)-CH_3} (tertiaire) & Reste \textbf{orange} & Pas d'oxydation ménagée \\
\hline
\end{tabularx}
\end{center}

\begin{center}
\begin{tabularx}{\linewidth}{|l|X|X|}
\hline
\rowcolor{popgreen!50} \textbf{Test d'identification} & \textbf{Contenu du tube 1 (éthanol oxydé)} & \textbf{Contenu du tube 2 (propan-2-ol oxydé)} \\
\hline
\rowcolor{popgreenL} DNPH & Précipité jaune-orangé (hydrazone) & Précipité jaune-orangé (hydrazone) \\
\hline
Liqueur de Fehling à chaud & Précipité \textbf{rouge brique} de \ce{Cu2O(s)} & Solution bleue inchangée \\
\hline
\rowcolor{popgreenL} Réactif de Schiff (variante) & Coloration rose-violacé & Aucune coloration \\
\hline
\end{tabularx}
\end{center}

Pour la lampe sans flamme : la spirale de cuivre, noircie par \ce{CuO}, redevient brillante et rougeoyante au contact des vapeurs d'éthanol ; une odeur piquante d'éthanal (pomme verte âcre) se dégage. La réaction est exothermique : elle entretient elle-même l'incandescence, d'où le nom de « lampe sans flamme ».

\courssub{Exploitation}

\begin{exoresolu}[Questions guidées et réponses]
\textbf{Question 1.} Pourquoi le dichromate de potassium change-t-il de couleur dans les tubes 1 et 2 ?\par
\textbf{Question 2.} Le test à la DNPH est positif dans les tubes 1 et 2, mais seul le tube 1 donne un précipité rouge brique avec la liqueur de Fehling. Interpréter.\par
\textbf{Question 3.} Écrire l'équation-bilan de l'oxydation ménagée de l'éthanol en éthanal par le dichromate en milieu acide, à l'aide des demi-équations électroniques.\par
\textbf{Question 4.} Pourquoi l'alcool tertiaire ne s'oxyde-t-il pas dans ces conditions ?\par
\textbf{Question 5.} Écrire l'équation-bilan de la réaction de la lampe sans flamme et préciser le rôle du cuivre.
\tcblower
\textbf{Réponse 1.} L'ion dichromate \ce{Cr2O7^2-_{(aq)}} (orange) est un oxydant puissant en milieu acide : il est réduit en ion chrome(III) \ce{Cr^3+_{(aq)}} (vert) pendant que l'alcool est oxydé. Le virage orange $\rightarrow$ vert prouve donc qu'une oxydation a eu lieu dans les tubes 1 et 2.

\textbf{Réponse 2.} La DNPH (2,4-dinitrophénylhydrazine) réagit avec tout composé carbonylé (aldéhyde \emph{ou} cétone) pour donner un précipité jaune-orangé d'hydrazone : le test positif dans les deux tubes confirme la présence d'un groupe \ce{C=O}. La liqueur de Fehling (ions \ce{Cu^2+} complexés en milieu basique) n'est réduite (en \ce{Cu2O} rouge brique) que par les aldéhydes : le précipité dans le tube 1 prouve la formation d'\cle{éthanal} \ce{CH3-CHO} (l'éthanol est un alcool primaire), tandis que le tube 2 contient une \cle{cétone}, la propanone \ce{CH3-CO-CH3}, qui ne réduit pas la liqueur de Fehling.

\textbf{Réponse 3.} Demi-équations électroniques :
\begin{align*}
\text{Oxydant :}\quad & \ce{Cr2O7^2-_{(aq)} + 14 H+_{(aq)} + 6 e- -> 2 Cr^3+_{(aq)} + 7 H2O_{(l)}} \\
\text{Réducteur :}\quad & \ce{CH3CH2OH_{(aq)} -> CH3CHO_{(aq)} + 2 H+_{(aq)} + 2 e-} \quad (\times 3)
\end{align*}
En additionnant (la seconde multipliée par 3 pour égaler les 6 électrons) :
\[ \ce{Cr2O7^2-_{(aq)} + 8 H+_{(aq)} + 3 CH3CH2OH_{(aq)} -> 2 Cr^3+_{(aq)} + 7 H2O_{(l)} + 3 CH3CHO_{(aq)}} \]
\vspace{2pt}
\textit{Remarque :} en excès de dichromate et chauffage prolongé, l'éthanal s'oxyde davantage en acide éthanoïque \ce{CH3COOH}.

\textbf{Réponse 4.} L'oxydation ménagée d'un alcool exige l'arrachage de l'atome d'hydrogène porté par le carbone fonctionnel (celui qui porte le groupe \ce{-OH}). Or, le carbone fonctionnel d'un alcool tertiaire ne porte \textbf{aucun atome d'hydrogène} (il est lié à trois atomes de carbone et à l'oxygène) : l'oxydation ménagée est donc impossible sans rompre la chaîne carbonée (ce qui demanderait des conditions drastiques, type combustion).

\textbf{Réponse 5.} Le cuivre chauffé à l'air s'oxyde en oxyde de cuivre(II) noir \ce{CuO} :
\[ \ce{2 Cu_{(s)} + O2_{(g)} ->[\Delta] 2 CuO_{(s)}} \]
L'oxyde de cuivre(II) oxyde alors les vapeurs d'éthanol en éthanal :
\[ \ce{CH3CH2OH_{(g)} + CuO_{(s)} -> CH3CHO_{(g)} + H2O_{(g)} + Cu_{(s)}} \]
Le cuivre métallique est régénéré à la fin de la réaction : il joue le rôle de \cle{catalyseur} (hétérogène). La chaleur dégagée par cette réaction exothermique entretient l'incandescence de la spirale.
\end{exoresolu}

\courssub{Conclusion}

Ce TP a permis de vérifier expérimentalement la règle fondamentale de l'oxydation ménagée des alcools :
\begin{itemize}
\item un alcool \cle{primaire} (éthanol) est oxydé en \textbf{aldéhyde} (éthanal), identifiable par la DNPH puis la liqueur de Fehling ou le réactif de Schiff ; en excès d'oxydant et chauffage prolongé, l'oxydation se poursuit jusqu'à l'\textbf{acide carboxylique} (acide éthanoïque) ;
\item un alcool \cle{secondaire} (propan-2-ol) est oxydé en \textbf{cétone} (propanone), positive à la DNPH mais négative à la liqueur de Fehling et au réactif de Schiff ;
\item un alcool \cle{tertiaire} (2-méthylpropan-2-ol) \textbf{ne subit pas} d'oxydation ménagée : la solution reste orange, l'ion \ce{Cr2O7^2-} n'est pas réduit.
\end{itemize}

L'expérience de la lampe sans flamme a en outre illustré l'oxydation catalytique de l'éthanol sur le cuivre, réaction exothermique produisant de l'éthanal. La classe d'un alcool conditionne donc entièrement le produit de son oxydation : c'est un résultat à retenir pour le Baccalauréat, et un principe utilisé chaque jour, du laboratoire de la SONARA au contrôle d'alcoolémie sur nos routes.

\newpage

\coursec{Méthodes et exercices résolus}

\courssub{Méthode 1 : Écrire l'équation-bilan de la réaction d'un alcool avec le sodium}

\begin{methode}[Alcool + sodium : le réflexe en 4 étapes]
\begin{enumerate}
\item \textbf{Identifier le proton mobile.} C'est l'hydrogène du groupe \cle{hydroxyle} \ce{-O-H}, jamais un hydrogène porté par le carbone. On écrit l'alcool sous la forme \ce{R-O-H} pour bien le visualiser.
\item \textbf{Écrire les produits.} Le sodium arrache ce proton : il se forme un \cle{alcoolate de sodium} (ou alcoolate) \ce{R-O^- Na+} et un dégagement de \cle{dihydrogène} \ce{H2}.
\item \textbf{Équilibrer.} Il faut \textbf{2 moles d'alcool pour 1 mole de \ce{H2}} :
\[ \ce{2 R-OH + 2 Na -> 2 R-O^- Na+ + H2 ^} \]
\item \textbf{Nommer l'alcoolate.} On remplace la terminaison \emph{-ol} de l'alcool par \emph{-anolate} : éthanol $\rightarrow$ éthanolate de sodium \ce{CH3-CH2-O^- Na+} ; propan-2-ol $\rightarrow$ propan-2-olate de sodium.
\end{enumerate}
\textbf{Réflexe de vérification :} le caractère ionique de l'alcoolate explique sa solubilité dans l'eau ; la solution obtenue est \textbf{basique} (l'ion alcoolate \ce{RO-} est une base forte : \ce{RO- + H2O -> ROH + HO-}).
\end{methode}

\begin{exoresolu}[Action du sodium sur le butan-1-ol]
On introduit un morceau de sodium de masse $m = \SI{0,46}{g}$ dans un excès de butan-1-ol \ce{C4H9OH}. 1) Écrire l'équation-bilan de la réaction. 2) Calculer le volume de dihydrogène dégagé, mesuré dans les CNTP ($V_m = \SI{22,4}{L.mol^{-1}}$). Donnée : $M(\ce{Na}) = \SI{23,0}{g.mol^{-1}}$.
\tcblower
\textbf{1) Équation-bilan.} Le proton mobile est celui du groupe hydroxyle. Le butan-1-ol s'écrit \ce{CH3-CH2-CH2-CH2-O-H} :
\[ \ce{2 C4H9OH + 2 Na -> 2 C4H9O^- Na+ + H2 ^} \]
Le produit organique est le \textbf{butan-1-olate de sodium}.

\textbf{2) Volume de dihydrogène.} Quantité de sodium introduite :
\[ n(\ce{Na}) = \frac{m}{M(\ce{Na})} = \frac{0,46}{23,0} = \SI{2,0e-2}{mol} \]
L'alcool est en excès, le sodium est donc le réactif limitant. D'après la stœchiométrie, $n(\ce{H2}) = \dfrac{n(\ce{Na})}{2} = \SI{1,0e-2}{mol}$.
\[ V(\ce{H2}) = n(\ce{H2}) \times V_m = 1,0 \times 10^{-2} \times 22,4 = \SI{0,224}{L} = \SI{224}{mL} \]
\textbf{Conclusion :} la réaction dégage \SI{224}{mL} de dihydrogène (détectable par la petite détonation à l'approche d'une flamme) et laisse une solution de butan-1-olate de sodium, basique.
\end{exoresolu}

\courssub{Méthode 2 : Écrire les équations de déshydratation d'un alcool}

\begin{methode}[Distinguer et écrire les deux déshydratations]
\begin{enumerate}
\item \textbf{Repérer le type demandé.} \textbf{Intramoléculaire} : une seule molécule d'alcool perd \ce{H2O} $\rightarrow$ produit un \cle{alcène}. \textbf{Intermoléculaire} : deux molécules d'alcool perdent ensemble une molécule d'eau $\rightarrow$ produit un \cle{étheroxyde} \ce{R^1-O-R^2}.
\item \textbf{Retenir les conditions expérimentales (très demandées au Bac) :}
\begin{center}
\begin{tabularx}{\linewidth}{|l|X|X|}
\hline
\rowcolor{popblueL} & \textbf{Intramoléculaire} & \textbf{Intermoléculaire} \\
\hline
Catalyseur & \ce{Al2O3} vers \SI{400}{\celsius}, ou \ce{H2SO4} / \ce{H3PO4} concentré & \ce{H2SO4} concentré \\
\hline
Température & élevée (vers \SI{180}{\celsius} avec \ce{H2SO4}) & modérée (vers \SI{140}{\celsius}) \\
\hline
Produit & alcène + eau & étheroxyde + eau \\
\hline
\end{tabularx}
\end{center}
\item \textbf{Écrire l'équation.} Intramoléculaire : \ce{C_{$n$}H_{$2n+1$}OH -> C_{$n$}H_{$2n$} + H2O}. Intermoléculaire : \ce{2 R-OH -> R-O-R + H2O}.
\item \textbf{Appliquer la règle de Zaïtsev} si l'alcool peut donner plusieurs alcènes : l'alcène \textbf{majoritaire} est le plus substitué sur la double liaison.
\end{enumerate}
\end{methode}

\begin{exoresolu}[Déshydratation de l'éthanol et du butan-2-ol]
1) On chauffe de l'éthanol en présence d'acide sulfurique concentré. Écrire les équations des deux déshydratations possibles et préciser les conditions qui favorisent chacune. 2) La déshydratation intramoléculaire du butan-2-ol donne deux alcènes isomères. Les écrire et préciser le majoritaire.
\tcblower
\textbf{1) Éthanol.}
\begin{itemize}
\item \textbf{Intermoléculaire} (vers \SI{140}{\celsius}, excès d'alcool) :
\[ \ce{2 CH3-CH2-OH ->[$H_2SO_4$, \SI{140}{\celsius}] CH3-CH2-O-CH2-CH3 + H2O} \]
Le produit est l'\textbf{éthoxyéthane} (éther diéthylique), ancien anesthésiant.
\item \textbf{Intramoléculaire} (vers \SI{180}{\celsius}, ou passage des vapeurs sur \ce{Al2O3} à \SI{400}{\celsius}) :
\[ \ce{CH3-CH2-OH ->[$Al_2O_3$, \SI{400}{\celsius}] CH2=CH2 + H2O} \]
Le produit est l'\textbf{éthène}, matière première des plastiques (polyéthylène).
\end{itemize}
\textbf{2) Butan-2-ol.} L'élimination de \ce{H2O} peut porter sur le carbone 1 ou le carbone 3 :
\[ \ce{CH3-CH2-CH(OH)-CH3 -> CH3-CH2-CH=CH2 + H2O} \quad \text{(but-1-ène)} \]
\[ \ce{CH3-CH2-CH(OH)-CH3 -> CH3-CH=CH-CH3 + H2O} \quad \text{(but-2-ène)} \]
D'après la \textbf{règle de Zaïtsev}, le \textbf{but-2-ène} est majoritaire car sa double liaison est la plus substituée (deux groupes alkyle contre un seul pour le but-1-ène).
\textbf{Conclusion :} la température et le catalyseur orientent la réaction vers l'éther ou vers l'alcène ; quand plusieurs alcènes sont possibles, le plus substitué prédomine.
\end{exoresolu}

\courssub{Méthode 3 : Conduire l'oxydation ménagée d'un alcool (demi-équations électroniques)}

\begin{methode}[Oxydation ménagée : la démarche complète]
\begin{enumerate}
\item \textbf{Identifier la classe de l'alcool} (compter les carbones liés au carbone fonctionnel) et \textbf{prévoir le produit} :
\begin{itemize}
\item Alcool \textbf{primaire} $\rightarrow$ \textbf{aldéhyde} $\rightarrow$ (si l'oxydant est en excès) \textbf{acide carboxylique} ;
\item Alcool \textbf{secondaire} $\rightarrow$ \textbf{cétone} (unique étape) ;
\item Alcool \textbf{tertiaire} $\rightarrow$ \textbf{pas d'oxydation ménagée} (pas d'hydrogène sur le carbone fonctionnel).
\end{itemize}
\item \textbf{Écrire la demi-équation de l'oxydant} en milieu acide :
\[ \ce{MnO4^- + 8 H+ + 5 e- -> Mn^2+ + 4 H2O} \qquad \ce{Cr2O7^2- + 14 H+ + 6 e- -> 2 Cr^3+ + 7 H2O} \]
\item \textbf{Écrire la demi-équation du couple de l'alcool} (sens oxydation, électrons à droite) :
\[ \ce{R-CH2-OH -> R-CHO + 2 H+ + 2 e-} \qquad \ce{R1-CHOH-R2 -> R1-CO-R2 + 2 H+ + 2 e-} \]
Pour l'oxydation de l'aldéhyde en acide : \ce{R-CHO + H2O -> R-COOH + 2 H+ + 2 e-}.
\item \textbf{Combiner} en multipliant pour égaliser les électrons (alcool : coefficient $\times 5$ avec \ce{MnO4^-}, $\times 3$ avec \ce{Cr2O7^2-}), puis simplifier \ce{H+} et \ce{H2O}.
\item \textbf{Se souvenir des couleurs} (indices observationnels au Bac) : \ce{MnO4^-} violet $\rightarrow$ \ce{Mn^2+} incolore ; \ce{Cr2O7^2-} orange $\rightarrow$ \ce{Cr^3+} vert.
\end{enumerate}
\end{methode}

\begin{exoresolu}[Oxydation du propan-1-ol par le dichromate de potassium]
On oxyde le propan-1-ol par une solution acidifiée de dichromate de potassium \ce{K2Cr2O7} en défaut, de sorte que l'oxydation s'arrête à l'aldéhyde. 1) Écrire les deux demi-équations électroniques puis l'équation-bilan. 2) Quel volume de solution de dichromate de concentration $C = \SI{0,10}{mol.L^{-1}}$ faut-il pour oxyder $n = \SI{6,0e-2}{mol}$ de propan-1-ol ?
\tcblower
\textbf{1) Équations.} Le propan-1-ol est un alcool primaire ; en défaut d'oxydant, il donne le \textbf{propanal} \ce{CH3-CH2-CHO}.

Demi-équation d'oxydation de l'alcool :
\[ \ce{CH3-CH2-CH2-OH -> CH3-CH2-CHO + 2 H+ + 2 e-} \]
Demi-équation de réduction du dichromate :
\[ \ce{Cr2O7^2- + 14 H+ + 6 e- -> 2 Cr^3+ + 7 H2O} \]
On multiplie la première par 3 pour échanger 6 électrons, puis on additionne :
\[ \ce{3 CH3-CH2-CH2-OH + Cr2O7^2- + 8 H+ -> 3 CH3-CH2-CHO + 2 Cr^3+ + 7 H2O} \]
(Vérification : $14\,\ce{H+} - 6\,\ce{H+} = 8\,\ce{H+}$ à gauche ; charges : $-2 + 8 = +6$ des deux côtés.)

\textbf{2) Volume de solution.} D'après la stœchiométrie :
\[ n(\ce{Cr2O7^2-}) = \frac{n(\text{alcool})}{3} = \frac{6,0 \times 10^{-2}}{3} = \SI{2,0e-2}{mol} \]
\[ V = \frac{n(\ce{Cr2O7^2-})}{C} = \frac{2,0 \times 10^{-2}}{0,10} = \SI{0,20}{L} = \SI{200}{mL} \]
\textbf{Conclusion :} il faut \SI{200}{mL} de solution de dichromate ; au cours de la réaction, la solution passe de l'orange au vert, et le propanal formé peut être caractérisé par la DNPH (précipité jaune) et le réactif de Schiff (coloration rose).
\end{exoresolu}

\courssub{Méthode 4 : Identifier aldéhydes et cétones (tests caractéristiques)}

\begin{methode}[Stratégie d'identification en deux temps]
\begin{enumerate}
\item \textbf{Premier test : la DNPH} (2,4-dinitrophénylhydrazine). Un \textbf{précipité jaune-orangé} prouve la présence du groupe \cle{carbonyle} \ce{C=O} : le composé est un aldéhyde \emph{ou} une cétone. Sans précipité : ni l'un ni l'autre.
\item \textbf{Deuxième test : discriminant aldéhyde/cétone.} Les aldéhydes sont des \textbf{réducteurs}, les cétones non. Trois tests positifs avec les aldéhydes seulement :
\begin{itemize}
\item \textbf{Liqueur de Fehling} (bleue) $\rightarrow$ \textbf{précipité rouge brique} de \ce{Cu2O} à chaud ;
\item \textbf{Réactif de Tollens} (nitrate d'argent ammoniacal) $\rightarrow$ \textbf{miroir d'argent} ;
\item \textbf{Réactif de Schiff} $\rightarrow$ coloration \textbf{rose-violette}.
\end{itemize}
\item \textbf{Conclure rigoureusement} : DNPH positive + Fehling positive $\Rightarrow$ aldéhyde ; DNPH positive + Fehling négative $\Rightarrow$ cétone.
\item \textbf{Application à la lampe sans flamme} : les vapeurs d'éthanol s'oxydent sur le fil de cuivre incandescent (catalyseur) en \textbf{éthanal} (odeur de pomme), identifié par les tests ci-dessus ; le fil reste incandescent car la réaction est \textbf{exothermique} : \ce{CH3-CH2-OH + 1/2 O2 ->[Cu] CH3-CHO + H2O}.
\end{enumerate}
\end{methode}

\begin{exoresolu}[Analyse d'un produit d'oxydation]
On réalise l'oxydation ménagée d'un alcool A de formule brute \ce{C3H8O} par le permanganate de potassium acidifié. Le produit organique B obtenu donne un précipité jaune avec la DNPH mais ne réagit ni avec la liqueur de Fehling ni avec le réactif de Tollens. 1) Donner la nature de B et sa formule semi-développée. 2) En déduire la classe et la formule semi-développée de A. 3) Écrire l'équation-bilan de l'oxydation de A en B par \ce{MnO4^-} en milieu acide.
\tcblower
\textbf{1) Nature de B.} Le test à la DNPH est positif : B possède un groupe carbonyle (aldéhyde ou cétone). Les tests de Fehling et de Tollens sont négatifs : B n'est \textbf{pas réducteur}, donc B est une \textbf{cétone}. Avec 3 carbones, une seule cétone existe : la \textbf{propanone} (acétone) \ce{CH3-CO-CH3}.

\textbf{2) Nature de A.} Seul un alcool \textbf{secondaire} donne une cétone par oxydation ménagée. A est donc le \textbf{propan-2-ol} : \ce{CH3-CH(OH)-CH3}. (L'isomère propan-1-ol, primaire, aurait donné un aldéhyde puis un acide, tous deux positifs à la DNPH, et l'aldéhyde positif au Fehling.)

\textbf{3) Équation-bilan.} Demi-équations :
\[ \ce{CH3-CH(OH)-CH3 -> CH3-CO-CH3 + 2 H+ + 2 e-} \quad (\times 5) \]
\[ \ce{MnO4^- + 8 H+ + 5 e- -> Mn^2+ + 4 H2O} \quad (\times 2) \]
En additionnant (10 électrons échangés) et en simplifiant les \ce{H+} ($16 - 10 = 6$) :
\[ \ce{5 CH3-CH(OH)-CH3 + 2 MnO4^- + 6 H+ -> 5 CH3-CO-CH3 + 2 Mn^2+ + 8 H2O} \]
\textbf{Conclusion :} A est le propan-2-ol, B est la propanone ; la solution violette de permanganate se décolore au cours de la réaction.
\end{exoresolu}

\courssub{Méthode 5 : Écrire une estérification et calculer son rendement}

\begin{methode}[Estérification : équation, caractéristiques, rendement]
\begin{enumerate}
\item \textbf{Écrire l'équation} : acide carboxylique + alcool $\rightleftharpoons$ \ce{<=>} obligatoire (réaction limitée) :
\[ \ce{R-C(=O)OH + R$'$-OH <=> R-C(=O)O-R$'$ + H2O} \]
\item \textbf{Citer les trois caractéristiques} : réaction \textbf{lente}, \textbf{limitée} (équilibre) et \textbf{athermique}. On l'accélère par chauffage (reflux) et catalyse acide (\ce{H+}), sans modifier la limite.
\item \textbf{Connaître les limites d'équilibre} (mélange équimolaire) : alcool primaire $\approx 67\,\%$ ; alcool secondaire $\approx 60\,\%$ ; alcool tertiaire $\approx 5\,\%$.
\item \textbf{Calculer le rendement} par rapport au réactif limitant :
\[ \eta = \frac{n_{\text{ester obtenu}}}{n_{\text{ester théorique (limitant)}}} \times 100 \]
\item \textbf{Améliorer le rendement} : mettre un réactif en excès ou éliminer l'eau (ou l'ester) au fur et à mesure (déplacement d'équilibre).
\end{enumerate}
\end{methode}

\begin{exoresolu}[Synthèse de l'éthanoate d'éthyle]
On chauffe à reflux un mélange de $n_a = \SI{0,30}{mol}$ d'acide éthanoïque et de $n_b = \SI{0,30}{mol}$ d'éthanol, en présence de quelques gouttes d'acide sulfurique concentré. À l'équilibre, on dose l'acide restant : il en reste \SI{0,10}{mol}. 1) Écrire l'équation de la réaction et nommer l'ester. 2) Dresser le tableau d'avancement et déterminer l'avancement final. 3) Calculer le rendement de la synthèse et le comparer à la limite théorique. 4) Proposer deux moyens d'améliorer ce rendement.
\tcblower
\textbf{1) Équation.}
\[ \ce{CH3-COOH + CH3-CH2-OH <=> CH3-COO-CH2-CH3 + H2O} \]
L'ester formé est l'\textbf{éthanoate d'éthyle} (odeur fruitée, solvant).

\textbf{2) Tableau d'avancement.}
\begin{center}
\begin{tabularx}{\linewidth}{|l|X|X|X|X|}
\hline
\rowcolor{popblueL} État (mol) & \ce{CH3COOH} & \ce{C2H5OH} & \ce{CH3COOC2H5} & \ce{H2O} \\
\hline
Initial & 0,30 & 0,30 & 0 & 0 \\
\hline
Final ($x_f$) & $0,30 - x_f$ & $0,30 - x_f$ & $x_f$ & $x_f$ \\
\hline
\end{tabularx}
\end{center}
Il reste \SI{0,10}{mol} d'acide : $0,30 - x_f = 0,10$ d'où $x_f = \SI{0,20}{mol}$. Il s'est donc formé \SI{0,20}{mol} d'ester.

\textbf{3) Rendement.} Le mélange est équimolaire et la réaction est équimolaire : l'avancement maximal serait $x_{\max} = \SI{0,30}{mol}$.
\[ \eta = \frac{x_f}{x_{\max}} \times 100 = \frac{0,20}{0,30} \times 100 \approx 67\,\% \]
Cette valeur correspond exactement à la \textbf{limite d'équilibre} d'une estérification avec un alcool primaire : la réaction a atteint son équilibre, le catalyseur n'a fait que l'accélérer.

\textbf{4) Amélioration du rendement.} On peut : (a) introduire l'alcool \textbf{en large excès} (réactif bon marché) ; (b) \textbf{éliminer l'eau} ou distiller l'ester au fur et à mesure de sa formation. Dans les deux cas, l'équilibre est déplacé vers la formation de l'ester (loi de Le Chatelier).
\textbf{Conclusion :} le rendement de $67\,\%$ n'est pas un échec expérimental mais la conséquence du caractère limité de l'estérification.
\end{exoresolu}

\courssub{Méthode 6 : Préparer un alcool — hydratation d'un alcène et fermentation alcoolique}

\begin{methode}[Deux voies de préparation à maîtriser]
\begin{enumerate}
\item \textbf{Hydratation d'un alcène} (réaction inverse de la déshydratation intramoléculaire) :
\[ \ce{C_nH_{2n} + H2O ->[H3PO4] C_nH_{2n+1}OH} \]
Si l'alcène est \textbf{asymétrique}, appliquer la \textbf{règle de Markovnikov} : l'hydrogène de l'eau se fixe sur le carbone de la double liaison \emph{le plus hydrogéné}, le groupe \ce{-OH} sur le \emph{plus substitué} $\rightarrow$ on obtient majoritairement l'alcool de classe la plus élevée.
\item \textbf{Fermentation alcoolique} : transformation du glucose en éthanol par les enzymes des levures, \textbf{en l'absence d'air} (anaérobie), vers \SI{30}{\celsius} :
\[ \ce{C6H12O6 ->[enzymes] 2 C2H5OH + 2 CO2 ^} \]
Signes de la fermentation : dégagement gazeux (\ce{CO2}, trouble l'eau de chaux), élévation de température (réaction exothermique).
\item \textbf{Choisir la voie adaptée} : l'hydratation est une voie industrielle (pétrochimie) ; la fermentation est la voie traditionnelle et agroalimentaire (vin de palme, bili-bili, bières).
\end{enumerate}
\end{methode}

\begin{exoresolu}[Du glucose au vin de palme]
Le jus sucré de palme, riche en glucose, fermente spontanément grâce aux levures présentes dans l'air. 1) Écrire l'équation de la fermentation alcoolique du glucose et préciser les conditions. 2) Quelle masse d'éthanol peut-on théoriquement obtenir à partir de $m = \SI{360}{g}$ de glucose ? 3) En réalité, on recueille \SI{138}{g} d'éthanol. Calculer le rendement de la fermentation. Données : $M(\ce{C6H12O6}) = \SI{180}{g.mol^{-1}}$ ; $M(\ce{C2H5OH}) = \SI{46}{g.mol^{-1}}$.
\tcblower
\textbf{1) Équation et conditions.}
\[ \ce{C6H12O6 ->[\text{enzymes des levures}] 2 CH3-CH2-OH + 2 CO2 ^} \]
Conditions : milieu \textbf{anaérobie} (à l'abri de l'air, sinon l'éthanol s'oxyde en acide éthanoïque — le vin de palme « tourne » au vinaigre), température voisine de \SI{30}{\celsius}.

\textbf{2) Masse théorique d'éthanol.}
\[ n(\text{glucose}) = \frac{m}{M} = \frac{360}{180} = \SI{2,0}{mol} \]
D'après la stœchiométrie, $n(\text{éthanol}) = 2 \times n(\text{glucose}) = \SI{4,0}{mol}$.
\[ m_{\text{th}}(\text{éthanol}) = n \times M = 4,0 \times 46 = \SI{184}{g} \]

\textbf{3) Rendement.}
\[ \eta = \frac{m_{\text{obtenue}}}{m_{\text{théorique}}} \times 100 = \frac{138}{184} \times 100 = 75\,\% \]
\textbf{Conclusion :} la fermentation de \SI{360}{g} de glucose fournit théoriquement \SI{184}{g} d'éthanol ; le rendement réel de $75\,\%$ s'explique par les pertes (évaporation, oxydation partielle en vinaigre, glucose consommé par la respiration des levures).
\end{exoresolu}

\begin{attention}[Les 5 erreurs qui coûtent des points au Bac]
\begin{enumerate}
\item \textbf{Écrire \ce{->} au lieu de \ce{<=>} pour l'estérification.} L'estérification est \emph{limitée} : la double flèche est exigée, ainsi que l'eau parmi les produits (souvent oubliée).
\item \textbf{Faire réagir un alcool tertiaire en oxydation ménagée.} Le 2-méthylpropan-2-ol ne réagit \emph{pas} avec \ce{MnO4^-} ou \ce{Cr2O7^2-} : le carbone fonctionnel ne porte pas d'hydrogène. Dire « pas de réaction » est la bonne réponse.
\item \textbf{Mal équilibrer les demi-équations redox.} Réflexes : équilibrer O avec \ce{H2O}, H avec \ce{H+}, les charges avec \ce{e-} ; vérifier que les électrons disparaissent de l'équation-bilan finale et que les charges sont conservées.
\item \textbf{Confondre les deux déshydratations.} Alcène = intramoléculaire (température élevée, \SI{180}{\celsius} ou \ce{Al2O3} à \SI{400}{\celsius}) ; étheroxyde = intermoléculaire (\SI{140}{\celsius}). Inverser les températures est une faute classique.
\item \textbf{Oublier les coefficients 2 dans les équations de fermentation et d'action du sodium.} \ce{C6H12O6 -> 2 C2H5OH + 2 CO2} et \ce{2 ROH + 2 Na -> 2 RO^- Na+ + H2} : sans ces coefficients, tous les calculs stœchiométriques qui suivent sont faux.
\end{enumerate}
\end{attention}

\newpage

\coursec{Exercices}

\courssub{Je vérifie mes connaissances}

\begin{exercice}[QCM : les bons réflexes]
Pour chaque question, choisir la (ou les) bonne(s) réponse(s) en justifiant brièvement.
\begin{enumerate}
\item Lorsqu'on verse un morceau de sodium dans de l'éthanol, on observe :
\begin{enumerate}
\item une explosion violente comparable à celle obtenue avec l'eau ;
\item un dégagement de dihydrogène ;
\item la formation d'éthanolate de sodium \ce{C2H5ONa} ;
\item un dégagement de dioxygène.
\end{enumerate}
\item La déshydratation \emph{intramoléculaire} de l'éthanol en présence d'acide sulfurique concentré vers \SI{180}{\celsius} donne :
\begin{enumerate}
\item un étheroxyde ; \item un alcène ; \item un aldéhyde ; \item une cétone.
\end{enumerate}
\item L'oxydation ménagée d'un alcool secondaire conduit à :
\begin{enumerate}
\item un aldéhyde ; \item une cétone ; \item un acide carboxylique ; \item un alcène.
\end{enumerate}
\item La réaction d'estérification est :
\begin{enumerate}
\item lente ; \item totale ; \item limitée par une réaction inverse ; \item athermique.
\end{enumerate}
\end{enumerate}
\end{exercice}

\begin{exercice}[Vrai ou faux, en justifiant]
Répondre par \textbf{vrai} ou \textbf{faux} et justifier chaque réponse par une phrase ou une équation-bilan.
\begin{enumerate}
\item Tous les alcools s'oxydent en acide carboxylique par action du permanganate de potassium en milieu acide.
\item Le cuivre, dans l'expérience de la lampe sans flamme, est un catalyseur : il n'apparaît pas dans l'équation-bilan globale.
\item La liqueur de Fehling permet de distinguer un aldéhyde d'une cétone.
\item L'hydratation du propène conduit majoritairement au propan-1-ol.
\item La fermentation alcoolique du glucose est une réaction qui se déroule en présence abondante de dioxygène.
\end{enumerate}
\end{exercice}

\begin{exercice}[Texte à trous : l'oxydation ménagée]
Recopier et compléter le texte suivant avec les mots et expressions de la liste ci-dessous (certains peuvent être utilisés plusieurs fois) : \emph{primaire ; secondaire ; tertiaire ; aldéhyde ; cétone ; acide carboxylique ; n'a pas d'hydrogène ; orange ; vert ; oxydant}.

L'oxydation ménagée d'un alcool \ldots{} conduit d'abord à un \ldots{}, qui peut être oxydé à son tour en \ldots{}. Celle d'un alcool \ldots{} conduit à une \ldots{}, étape finale car la \ldots{} ne peut pas être oxydée davantage. Un alcool \ldots{} ne subit pas d'oxydation ménagée car son carbone fonctionnel \ldots{} lié à celui-ci. Lors de l'oxydation par le dichromate de potassium en milieu acide, la solution passe de la couleur \ldots{} à la couleur \ldots{}, ce qui traduit la réduction de l'ion \ldots{} \ce{Cr2O7^2-} en ion \ce{Cr^3+}.
\end{exercice}

\begin{exercice}[Définitions et formules]
\begin{enumerate}
\item Définir : déshydratation intermoléculaire ; oxydation ménagée ; estérification ; fermentation alcoolique.
\item Donner la formule semi-développée et la classe de chacun des alcools suivants : butan-1-ol ; butan-2-ol ; 2-méthylpropan-2-ol ; éthanol.
\item Rappeler le rôle des réactifs suivants et le résultat positif attendu avec un aldéhyde : DNPH ; liqueur de Fehling ; réactif de Schiff ; nitrate d'argent ammoniacal (réactif de Tollens).
\end{enumerate}
\end{exercice}

\courssub{Je m'entraîne}

\begin{exercice}[Action du sodium sur l'éthanol]
On introduit un petit morceau de sodium dans un tube contenant de l'éthanol absolu.
\begin{enumerate}
\item Décrire les observations expérimentales et préciser le test qui permet d'identifier le gaz dégagé.
\item Écrire l'équation-bilan de la réaction et nommer le produit organique formé.
\item Quelle propriété du groupe hydroxyle cette réaction met-elle en évidence ?
\item On fait réagir \SI{0,10}{mol} de sodium avec un excès d'éthanol. Calculer le volume de gaz dégagé dans les CNTP ($V_m = \SI{22,4}{L.mol^{-1}}$).
\end{enumerate}
\end{exercice}

\begin{exercice}[Déshydratation du propan-1-ol]
On chauffe du propan-1-ol en présence d'acide sulfurique concentré.
\begin{enumerate}
\item À \SI{180}{\celsius} environ, on recueille un composé A gazeux qui décolore l'eau de brome. Donner la nature de la réaction, la formule semi-développée et le nom de A, puis écrire l'équation-bilan.
\item À \SI{140}{\celsius} environ, avec un excès d'alcool, on obtient un composé B de formule brute \ce{C6H14O}. Donner la nature de la réaction, la formule semi-développée et le nom de B, puis écrire l'équation-bilan.
\item Citer une application industrielle ou domestique d'un alcène et d'un étheroxyde.
\end{enumerate}
\end{exercice}

\begin{exercice}[Oxydation ménagée du butan-2-ol]
On oxyde le butan-2-ol par une solution acidifiée de dichromate de potassium. Le produit organique C obtenu donne un précipité jaune avec la DNPH mais ne réduit pas la liqueur de Fehling.
\begin{enumerate}
\item Quelle est la classe du butan-2-ol ? Justifier.
\item Que révèlent les deux tests réalisés sur C ? En déduire la nature, la formule semi-développée et le nom de C.
\item Écrire les demi-équations électroniques des couples \ce{Cr2O7^2-}/\ce{Cr^3+} et du couple de l'alcool, puis l'équation-bilan de l'oxydation.
\item Quel changement de couleur observe-t-on dans le milieu réactionnel ?
\end{enumerate}
\end{exercice}

\begin{exercice}[Estérification et rendement]
On chauffe un mélange équimolaire de \SI{0,20}{mol} d'acide éthanoïque et de \SI{0,20}{mol} d'éthanol en présence de quelques gouttes d'acide sulfurique concentré. Après équilibre, on isole \SI{11,7}{g} d'éthanoate d'éthyle.
\par\textbf{Données :} $M(\text{éthanoate d'éthyle}) = \SI{88}{g.mol^{-1}}$.
\begin{enumerate}
\item Écrire l'équation-bilan de la réaction et nommer tous les produits.
\item Citer les trois caractéristiques de la réaction d'estérification et préciser le rôle de l'acide sulfurique.
\item Calculer la quantité théorique d'ester attendue si la réaction était totale, puis le rendement de la synthèse.
\item Proposer deux moyens pratiques d'améliorer ce rendement.
\end{enumerate}
\end{exercice}

\courssub{Je me prépare au Bac}

\begin{exercice}[Type Bac — Le vin de palme et l'éthanol]
Le vin de palme, boisson très répandue dans les régions du Littoral et du Sud-Ouest, doit ses propriétés à l'éthanol produit par fermentation des sucres de la sève. On se propose d'étudier quelques propriétés chimiques de cet alcool.
\par\textbf{Données :} $M(\ce{Na}) = \SI{23}{g.mol^{-1}}$ ; $M(\ce{C2H5OH}) = \SI{46}{g.mol^{-1}}$ ; $M(\ce{C2H5ONa}) = \SI{68}{g.mol^{-1}}$ ; $V_m = \SI{22,4}{L.mol^{-1}}$ dans les CNTP.
\begin{enumerate}
\item \textbf{Action du sodium.} On fait réagir \SI{4,6}{g} de sodium avec un excès d'éthanol.
\begin{enumerate}
\item Écrire l'équation-bilan de la réaction et nommer les produits.
\item Calculer le volume de dihydrogène dégagé dans les CNTP.
\item Calculer la masse d'éthanolate de sodium formé.
\item Comparer la violence de cette réaction à celle du sodium avec l'eau et proposer une explication.
\end{enumerate}
\item \textbf{Lampe sans flamme.} On réalise l'expérience de la lampe sans flamme avec les vapeurs d'éthanol et un fil de cuivre chauffé au rouge.
\begin{enumerate}
\item Décrire le dispositif expérimental (schéma annoté exigé) et les observations.
\item Écrire l'équation-bilan de la réaction et préciser le rôle du cuivre.
\item Le produit formé est testé avec la DNPH puis avec le réactif de Schiff. Donner les résultats attendus et conclure sur sa nature.
\end{enumerate}
\item \textbf{Oxydation ménagée poussée.} L'éthanol est maintenant oxydé par un excès de permanganate de potassium en milieu acide jusqu'à l'acide éthanoïque.
\begin{enumerate}
\item Écrire les demi-équations électroniques des couples \ce{MnO4-}/\ce{Mn^2+} et \ce{CH3COOH}/\ce{CH3CH2OH}, puis l'équation-bilan.
\item Quelle observation colorée permet de suivre la réaction ?
\end{enumerate}
\end{enumerate}
\end{exercice}

\begin{exercice}[Type Bac — Identification d'un alcool et fermentation du bili-bili]
Dans une brasserie artisanale du Nord, le bili-bili est obtenu par fermentation de sucres issus du mil ou du sorgho. Un technicien de laboratoire dispose par ailleurs d'un alcool A de formule brute \ce{C4H10O}.
\par\textbf{Données :} $M(\ce{C6H12O6}) = \SI{180}{g.mol^{-1}}$ ; $M(\ce{C2H5OH}) = \SI{46}{g.mol^{-1}}$ ; $V_m = \SI{22,4}{L.mol^{-1}}$ dans les CNTP ; masse volumique de l'éthanol : $\rho = \SI{0,79}{g.mL^{-1}}$.
\begin{enumerate}
\item \textbf{Identification de A.} L'oxydation ménagée de A par le dichromate de potassium acidifié fournit un composé B. B donne un précipité jaune avec la DNPH mais ne réduit pas la liqueur de Fehling.
\begin{enumerate}
\item Que signifie chacun de ces deux tests ? En déduire la nature de B et la classe de A.
\item Donner les formules semi-développées possibles de A, puis identifier A sachant que sa chaîne carbonée est linéaire. Nommer A et B.
\item Écrire les demi-équations électroniques et l'équation-bilan de l'oxydation de A en B.
\item Un isomère de A, le 2-méthylpropan-2-ol, est soumis au même oxydant. Que se passe-t-il ? Justifier.
\end{enumerate}
\item \textbf{Fermentation alcoolique.} On réalise en anaérobiose la fermentation de \SI{450}{g} de glucose supposée totale.
\begin{enumerate}
\item Écrire l'équation-bilan de la fermentation alcoolique du glucose.
\item Calculer la masse d'éthanol et le volume de dioxyde de carbone produits (CNTP).
\item Le volume final de boisson obtenue est de \SI{2,5}{L}. Calculer le degré alcoolique approximatif de la boisson (volume d'éthanol pur pour \SI{100}{mL} de boisson).
\item Pourquoi la fermentation doit-elle être conduite à l'abri de l'air ? Que risque-t-il de se passer sinon pour le goût de la boisson ?
\end{enumerate}
\end{enumerate}
\end{exercice}

\begin{exercice}[Type Bac — Hydratation d'un alcène et synthèse d'un ester]
Une unité de traitement souhaite valoriser le but-1-ène en alcool, puis utiliser cet alcool pour préparer un ester à odeur fruitée destiné à l'agroalimentaire.
\par\textbf{Données :} $M(\text{acide éthanoïque}) = \SI{60}{g.mol^{-1}}$ ; $M(\text{butan-2-ol}) = \SI{74}{g.mol^{-1}}$ ; $M(\text{éthanoate de butyle}) = \SI{116}{g.mol^{-1}}$.
\begin{enumerate}
\item \textbf{Hydratation du but-1-ène.}
\begin{enumerate}
\item Écrire l'équation-bilan de l'hydratation du but-1-ène et nommer les deux alcools susceptibles de se former.
\item Énoncer la règle de Markovnikov et en déduire le produit majoritaire.
\item Donner la classe de chacun des deux alcools obtenus.
\item Le produit majoritaire est soumis à une oxydation ménagée par \ce{KMnO4} en milieu acide. Donner la nature, le nom et la formule semi-développée du produit obtenu, et proposer un test chimique pour l'identifier.
\end{enumerate}
\item \textbf{Déshydratation.} On chauffe ensuite le butan-2-ol avec de l'acide sulfurique concentré à \SI{180}{\celsius}.
\begin{enumerate}
\item Écrire l'équation-bilan de la réaction et nommer le(s) produit(s) organique(s).
\item Quel test simple permet de vérifier la présence d'une double liaison dans le produit ?
\end{enumerate}
\item \textbf{Synthèse de l'ester.} On fait réagir \SI{12,0}{g} d'acide éthanoïque avec \SI{14,8}{g} de butan-2-ol. À l'équilibre, on recueille \SI{13,9}{g} d'éthanoate de butyle.
\begin{enumerate}
\item Écrire l'équation-bilan de la réaction d'estérification.
\item Vérifier que le mélange initial est équimolaire, puis construire le tableau d'avancement.
\item Calculer le rendement de la synthèse.
\item Rappeler les trois caractéristiques de cette réaction et proposer deux méthodes pour déplacer l'équilibre vers la formation de l'ester.
\end{enumerate}
\end{enumerate}
\end{exercice}

\newpage

\coursec{Corrigés des exercices}

\begin{corrige}[Exercice 1 — QCM : les bons réflexes]
\begin{enumerate}
\item \textbf{Bonnes réponses : b et c.} Le sodium réagit avec l'éthanol selon :
\[ \ce{2 C2H5OH + 2 Na -> 2 C2H5ONa + H2} \]
Il y a bien dégagement de \cle{dihydrogène} (b) et formation d'\cle{éthanolate de sodium} (c). La réaction est vive mais beaucoup moins violente qu'avec l'eau (a faux) ; le gaz dégagé est \ce{H2} et non \ce{O2} (d faux).
\item \textbf{Bonne réponse : b.} Vers \SI{180}{\celsius} avec \ce{H2SO4} concentré, la déshydratation est \emph{intramoléculaire} : elle fournit un \cle{alcène}, l'éthylène (éthène) \ce{CH2=CH2}. L'étheroxyde est obtenu vers \SI{140}{\celsius} (déshydratation intermoléculaire).
\item \textbf{Bonne réponse : b.} Un alcool \cle{secondaire} s'oxyde en \cle{cétone} (le carbone fonctionnel ne porte qu'un seul hydrogène : une seule étape d'oxydation est possible).
\item \textbf{Bonnes réponses : a, c et d.} L'estérification est \cle{lente}, \cle{limitée} par la réaction inverse (hydrolyse de l'ester) et \cle{athermique}. Elle n'est jamais totale (b faux) : on aboutit à un équilibre.
\end{enumerate}
\end{corrige}

\begin{corrige}[Exercice 2 — Vrai ou faux]
\begin{enumerate}
\item \textbf{Faux.} Seuls les alcools \emph{primaires} conduisent à un acide carboxylique (via l'aldéhyde). Un alcool secondaire donne une cétone, et un alcool tertiaire ne subit pas d'oxydation ménagée.
\item \textbf{Vrai.} Le cuivre est oxydé en \ce{CuO} puis régénéré : il accélère la réaction sans apparaître dans le bilan global $\ce{CH3CH2OH} + \frac{1}{2}\ce{O2} \ce{->} \ce{CH3CHO} + \ce{H2O}$. C'est la définition d'un \cle{catalyseur}.
\item \textbf{Vrai.} La liqueur de Fehling est réduite par les aldéhydes (précipité rouge brique d'oxyde de cuivre(I) \ce{Cu2O}) alors que les cétones ne réagissent pas : elle distingue donc un aldéhyde d'une cétone.
\item \textbf{Faux.} D'après la \cle{règle de Markovnikov}, l'hydrogène se fixe sur le carbone de la double liaison le plus hydrogéné : l'hydratation du propène donne majoritairement le \textbf{propan-2-ol} (alcool secondaire), et non le propan-1-ol.
\item \textbf{Faux.} La fermentation alcoolique est une réaction \cle{anaérobie} : elle se déroule \emph{en l'absence} de dioxygène, sous l'action des enzymes des levures : $\ce{C6H12O6 -> 2 C2H5OH + 2 CO2}$.
\end{enumerate}
\end{corrige}

\begin{corrige}[Exercice 3 — Texte à trous : l'oxydation ménagée]
L'oxydation ménagée d'un alcool \textbf{primaire} conduit d'abord à un \textbf{aldéhyde}, qui peut être oxydé à son tour en \textbf{acide carboxylique}. Celle d'un alcool \textbf{secondaire} conduit à une \textbf{cétone}, étape finale car la \textbf{cétone} ne peut pas être oxydée davantage. Un alcool \textbf{tertiaire} ne subit pas d'oxydation ménagée car son carbone fonctionnel \textbf{n'a pas d'hydrogène} lié à celui-ci. Lors de l'oxydation par le dichromate de potassium en milieu acide, la solution passe de la couleur \textbf{orange} à la couleur \textbf{verte}, ce qui traduit la réduction de l'ion \textbf{oxydant} \ce{Cr2O7^2-} en ion \ce{Cr^3+}.
\end{corrige}

\begin{corrige}[Exercice 4 — Définitions et formules]
\begin{enumerate}
\item \textbf{Définitions.}
\begin{itemize}
\item \textbf{Déshydratation intermoléculaire} : élimination d'une molécule d'eau entre \emph{deux} molécules d'alcool, conduisant à un étheroxyde (vers \SI{140}{\celsius}, excès d'alcool, \ce{H2SO4} concentré).
\item \textbf{Oxydation ménagée} : oxydation qui conserve la chaîne carbonée et transforme uniquement le groupe fonctionnel (alcool $\to$ aldéhyde ou cétone $\to$ acide carboxylique), à l'aide d'un oxydant comme \ce{KMnO4} ou \ce{K2Cr2O7} en milieu acide.
\item \textbf{Estérification} : réaction entre un acide carboxylique et un alcool donnant un \emph{ester} et de l'eau ; elle est lente, limitée et athermique.
\item \textbf{Fermentation alcoolique} : transformation, en l'absence de dioxygène (anaérobiose), du glucose en éthanol et dioxyde de carbone sous l'action d'enzymes des levures.
\end{itemize}
\item \textbf{Formules et classes.}
\begin{center}
\begin{tabularx}{\linewidth}{|l|X|l|}
\hline
\rowcolor{popblueL} \textbf{Nom} & \textbf{Formule semi-développée} & \textbf{Classe} \\
\hline
Butan-1-ol & \chemfig{CH_3-CH_2-CH_2-CH_2-OH} & primaire \\
\hline
Butan-2-ol & \chemfig{CH_3-CH_2-CH(-[2]OH)-CH_3} & secondaire \\
\hline
2-Méthylpropan-2-ol & \chemfig{C(-[2]CH_3)(-[4]CH_3)(-[6]CH_3)-OH} & tertiaire \\
\hline
Éthanol & \chemfig{CH_3-CH_2-OH} & primaire \\
\hline
\end{tabularx}
\end{center}
\item \textbf{Réactifs d'identification.}
\begin{center}
\begin{tabularx}{\linewidth}{|l|X|X|}
\hline
\rowcolor{popblueL} \textbf{Réactif} & \textbf{Rôle} & \textbf{Test positif (aldéhyde)} \\
\hline
DNPH & Caractérise le groupe carbonyle \ce{C=O} (aldéhydes \emph{et} cétones) & Précipité jaune-orangé \\
\hline
Liqueur de Fehling & Distingue les aldéhydes des cétones & Précipité rouge brique de \ce{Cu2O} (à chaud) \\
\hline
Réactif de Schiff & Distingue les aldéhydes des cétones & Coloration rose-violette \\
\hline
Réactif de Tollens & Distingue les aldéhydes des cétones & Miroir d'argent sur les parois du tube \\
\hline
\end{tabularx}
\end{center}
\end{enumerate}
\end{corrige}

\begin{corrige}[Exercice 5 — Action du sodium sur l'éthanol]
\begin{enumerate}
\item \textbf{Observations :} le sodium fond en une bille qui se déplace à la surface de l'alcool, avec une vive effervescence (dégagement gazeux) ; la réaction est exothermique mais nettement moins violente qu'avec l'eau. \textbf{Test du gaz :} on approche une buchette incandescente de l'ouverture du tube : une légère détonation (« aboiement ») caractérise le \cle{dihydrogène} \ce{H2}.
\item \textbf{Équation-bilan :}
\[ \ce{2 CH3CH2OH + 2 Na -> 2 CH3CH2ONa + H2} \]
Le produit organique est l'\cle{éthanolate de sodium} \ce{C2H5ONa} (un alcoolate, ionique : \ce{C2H5O-} et \ce{Na+}).
\item Cette réaction met en évidence la \cle{mobilité de l'hydrogène} du groupe hydroxyle \ce{-OH} : cet hydrogène est « acide » (faiblement) et peut être arraché par un métal réducteur comme le sodium.
\item D'après l'équation : $n(\ce{H2}) = \dfrac{n(\ce{Na})}{2} = \dfrac{0{,}10}{2} = \SI{0,050}{mol}$.
\[ V(\ce{H2}) = n \times V_m = 0{,}050 \times 22{,}4 = \SI{1,12}{L} \]
Le volume de dihydrogène dégagé dans les CNTP est $\boxed{V = \SI{1,12}{L}}$.
\end{enumerate}
\end{corrige}

\begin{corrige}[Exercice 6 — Déshydratation du propan-1-ol]
\begin{enumerate}
\item À \SI{180}{\celsius} : \textbf{déshydratation intramoléculaire}. A décolore l'eau de brome : c'est un \cle{alcène}, le \textbf{propène} \ce{CH2=CH-CH3}.
\[ \ce{CH3-CH2-CH2OH ->[H2SO4,\ 180\,^\circ\mathrm{C}] CH2=CH-CH3 + H2O} \]
\item À \SI{140}{\celsius} avec excès d'alcool : \textbf{déshydratation intermoléculaire}. B est un \cle{étheroxyde} : le \textbf{1-propoxypropane} (éther dipropylique) \chemfig{CH_3-CH_2-CH_2-O-CH_2-CH_2-CH_3}, de formule brute \ce{C6H14O}, conforme à l'énoncé.
\[ \ce{2 CH3CH2CH2OH ->[H2SO4,\ 140\,^\circ\mathrm{C}] CH3CH2CH2-O-CH2CH2CH3 + H2O} \]
\item \textbf{Applications :} un alcène comme l'éthylène sert à fabriquer les matières plastiques (polyéthylène des sachets et bidons) ; un étheroxyde comme l'éthoxyéthane (éther) est utilisé comme solvant et fut le premier anesthésique chirurgical.
\end{enumerate}
\end{corrige}

\begin{corrige}[Exercice 7 — Oxydation ménagée du butan-2-ol]
\begin{enumerate}
\item Le butan-2-ol \chemfig{CH_3-CH_2-CH(-[2]OH)-CH_3} est un alcool \textbf{secondaire} : le carbone fonctionnel (porteur du \ce{-OH}) est lié à \emph{deux} autres atomes de carbone.
\item Le précipité jaune avec la DNPH prouve la présence d'un \cle{groupe carbonyle} \ce{C=O} (aldéhyde ou cétone). L'absence de réduction de la liqueur de Fehling exclut l'aldéhyde : C est donc une \cle{cétone}, la \textbf{butanone} \chemfig{CH_3-CH_2-C(=[1]O)-CH_3}. C'est cohérent avec l'oxydation d'un alcool secondaire.
\item \textbf{Demi-équations :}
\begin{align*}
\text{Réduction : } & \ce{Cr2O7^2- + 14 H+ + 6 e- -> 2 Cr^3+ + 7 H2O} \\
\text{Oxydation : } & \ce{CH3-CH2-CH(OH)-CH3 -> CH3-CH2-CO-CH3 + 2 H+ + 2 e-} \quad (\times 3)
\end{align*}
\textbf{Équation-bilan :}
\[ \ce{3 CH3CH2CH(OH)CH3 + Cr2O7^2- + 8 H+ -> 3 CH3CH2COCH3 + 2 Cr^3+ + 7 H2O} \]
\item La solution passe de la couleur \textbf{orange} (ion dichromate \ce{Cr2O7^2-}) à la couleur \textbf{verte} (ion chrome(III) \ce{Cr^3+}).
\end{enumerate}
\end{corrige}

\begin{corrige}[Exercice 8 — Estérification et rendement]
\begin{enumerate}
\item \textbf{Équation-bilan :}
\[ \ce{CH3COOH + CH3CH2OH <=> CH3COOCH2CH3 + H2O} \]
Produits : l'\textbf{éthanoate d'éthyle} (ester) et l'\textbf{eau}.
\item L'estérification est \cle{lente}, \cle{limitée} (équilibre avec l'hydrolyse inverse) et \cle{athermique}. L'acide sulfurique joue le rôle de \cle{catalyseur} : il accélère l'atteinte de l'équilibre sans le déplacer.
\item Si la réaction était totale, le mélange étant équimolaire : $n_{\text{th}}(\text{ester}) = \SI{0,20}{mol}$.
Quantité réellement obtenue : $n(\text{ester}) = \dfrac{m}{M} = \dfrac{11{,}7}{88} = \SI{0,133}{mol}$.
\[ \eta = \frac{n_{\text{obtenue}}}{n_{\text{th}}} = \frac{0{,}133}{0{,}20} = 0{,}664 \approx \boxed{66\ \%} \]
(Valeur proche de la limite classique de $67\ \%$ pour un mélange équimolaire acide + alcool primaire.)
\item Pour améliorer le rendement : \textbf{introduire un réactif en excès} (alcool ou acide bon marché) ou \textbf{éliminer l'eau} (ou l'ester) au fur et à mesure de sa formation (distillation) : on déplace ainsi l'équilibre vers la formation de l'ester.
\end{enumerate}
\end{corrige}

\begin{corrige}[Exercice 9 — Type Bac : le vin de palme et l'éthanol]
\emph{Barème indicatif : question 1 (6 pts) ; question 2 (7 pts) ; question 3 (7 pts).}

\begin{enumerate}
\item \textbf{Action du sodium.}
\begin{enumerate}
\item Équation-bilan :
\[ \ce{2 CH3CH2OH + 2 Na -> 2 CH3CH2ONa + H2} \]
Produits : l'\textbf{éthanolate de sodium} et le \textbf{dihydrogène}.
\item $n(\ce{Na}) = \dfrac{m}{M} = \dfrac{4{,}6}{23} = \SI{0,20}{mol}$ ; d'où $n(\ce{H2}) = \dfrac{0{,}20}{2} = \SI{0,10}{mol}$.
\[ V(\ce{H2}) = 0{,}10 \times 22{,}4 = \boxed{\SI{2,24}{L}} \]
\item $n(\ce{C2H5ONa}) = n(\ce{Na}) = \SI{0,20}{mol}$ ; donc :
\[ m(\ce{C2H5ONa}) = 0{,}20 \times 68 = \boxed{\SI{13,6}{g}} \]
\item La réaction avec l'éthanol est \textbf{beaucoup moins violente} qu'avec l'eau (pas d'inflammation spontanée du dihydrogène). Explication : l'hydrogène du groupe \ce{-OH} de l'alcool est \emph{moins mobile} (moins « acide ») que celui de l'eau, car le groupe éthyle \ce{C2H5-} est donneur d'électrons, ce qui diminue la polarisation de la liaison \ce{O-H}.
\end{enumerate}
\item \textbf{Lampe sans flamme.}
\begin{enumerate}
\item \textbf{Dispositif :} un erlenmeyer contenant un peu d'éthanol est chauffé doucement ; les vapeurs d'éthanol passent sur un fil de cuivre en spirale, préalablement chauffé au rouge, suspendu au-dessus du liquide. \textbf{Observations :} le fil de cuivre \emph{reste incandescent} (la réaction est exothermique et s'auto-entretient, d'où le nom « lampe sans flamme ») ; il alterne noircissement (\ce{CuO}) et re-rougeoiement (\ce{Cu} régénéré) ; une odeur piquante (éthanal) se dégage.
\begin{popfigure}
\begin{tikzpicture}[scale=0.9, >=Stealth]
\draw[thick, fill=popblue!10] (-1.5,0) -- (-1.5,1.5) -- (-1.0,2.0) -- (1.0,2.0) -- (1.5,1.5) -- (1.5,0) -- cycle;
\node at (0,0.5) {\small éthanol liquide};
\draw[thick, poporange, decorate, decoration={coil, aspect=0.4, segment length=3pt, amplitude=3pt}] (0,2.0) -- (0,3.5);
\node[poporange, right] at (0.3,2.7) {\small fil de cuivre incandescent};
\draw[->, thick, popmuted] (-0.8,1.0) -- (-0.5,1.8) node[midway, left, font=\scriptsize] {vapeurs \ce{C2H5OH}};
\draw[thick, popmuted] (-1.8,-0.3) -- (1.8,-0.3);
\node[below, font=\scriptsize] at (0,-0.35) {chauffage doux (bec Bunsen)};
\end{tikzpicture}
\legende{Schéma de la lampe sans flamme : oxydation catalytique des vapeurs d'éthanol sur cuivre.}
\end{popfigure}
\item Équation-bilan :
\[ \ce{CH3CH2OH + \frac{1}{2} O2 ->[Cu] CH3CHO + H2O} \]
Le cuivre est un \cle{catalyseur} : il participe (cycle \ce{Cu}/\ce{CuO}) mais est régénéré et n'apparaît pas dans le bilan.
\item DNPH : \textbf{précipité jaune-orangé} $\Rightarrow$ présence d'un groupe carbonyle. Réactif de Schiff : \textbf{coloration rose-violette} $\Rightarrow$ le carbonyle est celui d'un \textbf{aldéhyde}. Conclusion : le produit est l'\textbf{éthanal} \ce{CH3CHO}.
\end{enumerate}
\item \textbf{Oxydation ménagée poussée.}
\begin{enumerate}
\item Demi-équations :
\begin{align*}
\text{Réduction : } & \ce{MnO4- + 8 H+ + 5 e- -> Mn^2+ + 4 H2O} \quad (\times 4)\\
\text{Oxydation : } & \ce{CH3CH2OH + H2O -> CH3COOH + 4 H+ + 4 e-} \quad (\times 5)
\end{align*}
Équation-bilan :
\[ \ce{5 CH3CH2OH + 4 MnO4- + 12 H+ -> 5 CH3COOH + 4 Mn^2+ + 11 H2O} \]
\item La solution violette du permanganate \ce{MnO4-} se \textbf{décolore} progressivement (formation d'ions \ce{Mn^2+} quasi incolores) tant que l'éthanol est en excès ; la décoloration cesse lorsque tout l'alcool est consommé.
\end{enumerate}
\end{enumerate}
\textbf{Points de vigilance :} ne pas oublier le coefficient 2 devant le sodium ; bien équilibrer les demi-équations en milieu acide avec \ce{H+} et \ce{H2O} ; le cuivre ne doit jamais figurer dans l'équation-bilan globale ; distinguer soigneusement DNPH (carbonyle) et Schiff/Fehling (aldéhyde).
\end{corrige}

\begin{corrige}[Exercice 10 — Type Bac : identification d'un alcool et fermentation du bili-bili]
\emph{Barème indicatif : partie 1 (10 pts) ; partie 2 (10 pts).}

\begin{enumerate}
\item \textbf{Identification de A.}
\begin{enumerate}
\item DNPH : précipité jaune $\Rightarrow$ B possède un \cle{groupe carbonyle} (aldéhyde ou cétone). Fehling négatif $\Rightarrow$ B n'est \emph{pas} un aldéhyde : B est une \cle{cétone}. A est donc un alcool \textbf{secondaire}.
\item Alcools secondaires de formule \ce{C4H10O} : seul le \textbf{butan-2-ol} \chemfig{CH_3-CH_2-CH(-[2]OH)-CH_3} convient (sa chaîne est d'ailleurs linéaire, ce qui confirme l'identification). A est le \textbf{butan-2-ol} ; B est la \textbf{butanone} \chemfig{CH_3-CH_2-C(=[1]O)-CH_3}.
\item Demi-équations et bilan :
\begin{align*}
& \ce{Cr2O7^2- + 14 H+ + 6 e- -> 2 Cr^3+ + 7 H2O}\\
& \ce{CH3CH2CH(OH)CH3 -> CH3CH2COCH3 + 2 H+ + 2 e-} \ (\times 3)
\end{align*}
\[ \ce{3 CH3CH2CH(OH)CH3 + Cr2O7^2- + 8 H+ -> 3 CH3CH2COCH3 + 2 Cr^3+ + 7 H2O} \]
\item Le 2-méthylpropan-2-ol \chemfig{C(-[2]CH_3)(-[4]CH_3)(-[6]CH_3)-OH} est un alcool \textbf{tertiaire} : son carbone fonctionnel ne porte \emph{aucun hydrogène}. Il \textbf{ne subit pas d'oxydation ménagée} : la solution de dichromate reste orange.
\end{enumerate}
\item \textbf{Fermentation alcoolique.}
\begin{enumerate}
\item Équation-bilan :
\[ \ce{C6H12O6 ->[enzymes] 2 C2H5OH + 2 CO2} \]
\item $n(\ce{C6H12O6}) = \dfrac{450}{180} = \SI{2,5}{mol}$.
\begin{align*}
n(\ce{C2H5OH}) &= 2 \times 2{,}5 = \SI{5,0}{mol} \Rightarrow m = 5{,}0 \times 46 = \boxed{\SI{230}{g}} \\
n(\ce{CO2}) &= \SI{5,0}{mol} \Rightarrow V = 5{,}0 \times 22{,}4 = \boxed{\SI{112}{L}}
\end{align*}
\item Volume d'éthanol pur : $V = \dfrac{m}{\rho} = \dfrac{230}{0{,}79} \approx \SI{291}{mL}$.
Degré alcoolique : $\dfrac{291}{2500} \times 100 \approx \boxed{11{,}6\ \%}$ en volume, soit environ $11{,}6$ degrés.
\item La fermentation alcoolique est \cle{anaérobie} : en présence de dioxygène, l'éthanol serait oxydé en \textbf{acide éthanoïque} (acétification, par des bactéries acétiques) : la boisson s'acidulerait et prendrait un goût de \textbf{vinaigre}.
\end{enumerate}
\end{enumerate}
\textbf{Points de vigilance :} un test DNPH positif seul ne suffit pas à conclure (aldéhyde \emph{ou} cétone) ; la stoechiométrie $1:2:2$ de la fermentation est une source fréquente d'erreur ; le degré alcoolique s'exprime en \emph{volume} d'éthanol, d'où l'usage de la masse volumique.
\end{corrige}

\begin{corrige}[Exercice 11 — Type Bac : hydratation d'un alcène et synthèse d'un ester]
\emph{Barème indicatif : partie 1 (7 pts) ; partie 2 (4 pts) ; partie 3 (9 pts).}

\begin{enumerate}
\item \textbf{Hydratation du but-1-ène.}
\begin{enumerate}
\item Équation-bilan :
\[ \ce{CH2=CH-CH2-CH3 + H2O ->[H+] CH3-CH(OH)-CH2-CH3} \quad \text{(majoritaire)} \]
Les deux alcools possibles sont le \textbf{butan-2-ol} et le \textbf{butan-1-ol} \chemfig{CH_3-CH_2-CH_2-CH_2-OH}.
\item \textbf{Règle de Markovnikov :} lors de l'addition d'un composé hydrogéné sur un alcène asymétrique, l'hydrogène se fixe préférentiellement sur le carbone de la double liaison \emph{le plus hydrogéné}. Ici, H se fixe sur \ce{CH2} et \ce{OH} sur le carbone n°2 : le produit majoritaire est le \textbf{butan-2-ol}.
\item Butan-2-ol : alcool \textbf{secondaire} ; butan-1-ol : alcool \textbf{primaire}.
\item L'oxydation ménagée du butan-2-ol (secondaire) donne une \cle{cétone} : la \textbf{butanone} \chemfig{CH_3-CH_2-C(=[1]O)-CH_3}. \textbf{Test :} précipité jaune avec la DNPH (carbonyle) et \emph{absence} de réaction avec la liqueur de Fehling (ce n'est pas un aldéhyde).
\end{enumerate}
\item \textbf{Déshydratation.}
\begin{enumerate}
\item À \SI{180}{\celsius}, déshydratation intramoléculaire :
\[ \ce{CH3-CH(OH)-CH2-CH3 ->[H2SO4, \SI{180}{\celsius}] CH3-CH=CH-CH3 + H2O} \]
Produit principal : le \textbf{but-2-ène} (avec du but-1-ène minoritaire).
\item \textbf{Test :} l'\textbf{eau de brome} (rouge-orangée) est \textbf{décolorée} par addition sur la double liaison \ce{C=C}.
\end{enumerate}
\item \textbf{Synthèse de l'ester.}
\begin{enumerate}
\item Équation-bilan :
\[ \ce{CH3COOH + CH3-CH(OH)-CH2-CH3 <=> CH3COO-CH(CH3)-CH2-CH3 + H2O} \]
L'ester formé est l'\textbf{éthanoate de 1-méthylpropyle} (éthanoate de \emph{sec}-butyle).
\item Quantités initiales :
\[ n(\text{acide}) = \frac{12{,}0}{60} = \SI{0,20}{mol} \qquad n(\text{alcool}) = \frac{14{,}8}{74} = \SI{0,20}{mol} \]
Le mélange est bien \textbf{équimolaire}.
\begin{center}
\begin{tabular}{|c|c|c|c|c|}
\hline
\rowcolor{popblueL} & \ce{CH3COOH} & butan-2-ol & ester & \ce{H2O} \\
\hline
État initial (mol) & 0,20 & 0,20 & 0 & 0 \\
\hline
En cours (mol) & $0{,}20 - x$ & $0{,}20 - x$ & $x$ & $x$ \\
\hline
Si total : $x_{\max}$ (mol) & 0 & 0 & 0,20 & 0,20 \\
\hline
\end{tabular}
\end{center}
\item Quantité d'ester obtenue : $n = \dfrac{13{,}9}{116} \approx \SI{0,120}{mol}$.
\[ \eta = \frac{0{,}120}{0{,}20} \approx \boxed{60\ \%} \]
(Valeur cohérente : pour un alcool \emph{secondaire}, la limite d'équilibre est d'environ $60\ \%$, inférieure aux $67\ \%$ d'un alcool primaire.)
\item Caractéristiques : réaction \cle{lente}, \cle{limitée} (équilibre) et \cle{athermique}. Pour déplacer l'équilibre vers l'ester : \textbf{mettre un réactif en excès} (l'acide éthanoïque, peu coûteux) et/ou \textbf{éliminer l'eau ou l'ester} par distillation au fur et à mesure de sa formation.
\end{enumerate}
\end{enumerate}
\textbf{Points de vigilance :} bien appliquer Markovnikov (l'alcool majoritaire est le plus substitué) ; ne pas confondre éthanoate de butyle (butan-1-ol) et éthanoate de 1-méthylpropyle (butan-2-ol) ; le catalyseur \ce{H2SO4} accélère la réaction mais \emph{ne déplace pas} l'équilibre, donc n'améliore pas le rendement à l'équilibre.
\end{corrige}

\newpage

\coursec{Fiche bilan}

\begin{popfigure}
\begin{center}
\begin{tikzpicture}[
  mindmap/.style={},
  central/.style={circle, draw=popdark, fill=popblue, text=white, font=\bfseries\small, align=center, minimum size=2.6cm, inner sep=1pt},
  branche/.style={rounded corners=4pt, draw=#1, fill=#1!18, text=popdark, font=\scriptsize, align=center, inner sep=3pt, minimum height=0.9cm},
  lien/.style={draw=#1, line width=1.4pt}
]
\node[central] (C) at (0,0) {Propriétés\\chimiques\\des alcools};
\node[branche=poporange, anchor=west] (Na) at (2.6,2.6) {\textbf{Sodium}\\ \ce{R-OH + Na -> R-O^-Na^+ + 1/2 H2}\\mobilité de H du groupe $-$OH};
\node[branche=popgreen, anchor=west] (Desh) at (2.6,0.6) {\textbf{Déshydratation}\\ intra : alcène (\SI{180}{\celsius}, \ce{H2SO4} conc.)\\ inter : éther (\SI{140}{\celsius})};
\node[branche=poppurple, anchor=west] (Oxy) at (2.6,-1.8) {\textbf{Oxydation}\\ combustion : \ce{CO2 + H2O}\\ ménagée : \ce{MnO4^-}, \ce{Cr2O7^2-} / \ce{H+}};
\node[branche=poppink, anchor=east] (Lampe) at (-2.6,2.2) {\textbf{Lampe sans flamme}\\ vapeurs d'éthanol sur Cu\\ incandescent $\to$ éthanal};
\node[branche=popgold, anchor=east] (Tests) at (-2.6,0.4) {\textbf{Tests d'identification}\\ DNPH, Fehling, Schiff,\\ Tollens};
\node[branche=popblue, anchor=east] (Ester) at (-2.6,-1.4) {\textbf{Estérification}\\ \ce{acide + alcool <=> ester + eau}\\ lente, limitée, athermique};
\node[branche=popgreen!70!popblue, anchor=south] (Prep) at (0,-3.3) {\textbf{Préparation}\\ hydratation d'alcène (Markovnikov)\\ fermentation : \ce{C6H12O6 -> 2 C2H5OH + 2 CO2}};
\draw[lien=poporange] (C) -- (Na);
\draw[lien=popgreen] (C) -- (Desh);
\draw[lien=poppurple] (C) -- (Oxy);
\draw[lien=poppink] (C) -- (Lampe);
\draw[lien=popgold] (C) -- (Tests);
\draw[lien=popblue] (C) -- (Ester);
\draw[lien=popgreen!70!popblue] (C) -- (Prep);
\end{tikzpicture}
\end{center}
\legende{Carte mentale de la leçon : les six grandes réactions des alcools et leurs conditions.}
\end{popfigure}

\begin{bilan}[L'essentiel de la leçon]
\textbf{Définitions clés}
\begin{itemize}
  \item \cle{Déshydratation intramoléculaire} : départ de \ce{H2O} au sein d'une même molécule d'alcool $\to$ \textbf{alcène} (acide sulfurique concentré, vers \SI{180}{\celsius}).
  \item \cle{Déshydratation intermoléculaire} : départ de \ce{H2O} entre deux molécules d'alcool $\to$ \textbf{étheroxyde} (vers \SI{140}{\celsius}).
  \item \cle{Oxydation ménagée} : oxydation qui conserve la chaîne carbonée (oxydants : \ce{MnO4^-} ou \ce{Cr2O7^2-} en milieu acide, ou \ce{O2}/Cu).
  \item \cle{Estérification} : réaction entre un acide carboxylique et un alcool donnant un \textbf{ester} et de l'eau ; elle est \textbf{lente, limitée} (équilibre) et \textbf{athermique}.
  \item \cle{Fermentation alcoolique} : transformation anaérobie du glucose en éthanol et dioxyde de carbone (levures, vers \SIrange{25}{30}{\celsius}).
\end{itemize}

\textbf{Équations-types à connaître par cœur}
\begin{align*}
&\ce{R-OH + Na -> R-O^- + Na^+ + 1/2 H2 ^} &&\text{(alcoolate de sodium)}\\
&\ce{C2H5OH ->[{\ce{H2SO4}, \SI{180}{\celsius}}] CH2=CH2 + H2O} &&\text{(déshydratation intra)}\\
&\ce{2 C2H5OH ->[{\ce{H2SO4}, \SI{140}{\celsius}}] C2H5-O-C2H5 + H2O} &&\text{(déshydratation inter)}\\
&\ce{C_nH_{2n+1}OH + \tfrac{3n}{2} O2 -> n CO2 + (n+1) H2O} &&\text{(combustion complète)}\\
&\ce{CH3-CH2-OH + O2 ->[Cu] CH3-CHO + H2O} &&\text{(lampe sans flamme)}\\
&\ce{CH3COOH + C2H5OH <=> CH3COO-C2H5 + H2O} &&\text{(estérification)}\\
&\ce{CH2=CH2 + H2O ->[{\ce{H+}}] CH3-CH2-OH} &&\text{(hydratation d'alcène)}\\
&\ce{C6H12O6 -> 2 C2H5OH + 2 CO2} &&\text{(fermentation alcoolique)}
\end{align*}

\textbf{Oxydation ménagée selon la classe de l'alcool}
\begin{center}
\begin{tabularx}{\linewidth}{|l|X|X|X|}
\hline
\rowcolor{popblueL} \textbf{Classe} & \textbf{1\ier{} produit} & \textbf{Si excès d'oxydant} & \textbf{Tests} \\
\hline
Primaire \ce{$R$-CH2-OH} & aldéhyde \ce{$R$-CHO} & acide carboxylique \ce{$R$-COOH} & DNPH $+$ ; Fehling $+$ ; Schiff rose ; Tollens $+$ \\
\hline
Secondaire \ce{$R$-CHOH-$R'$} & cétone \ce{$R$-CO-$R'$} & (s'arrête à la cétone) & DNPH $+$ ; Fehling $-$ ; Schiff $-$ ; Tollens $-$ \\
\hline
Tertiaire \ce{$R$-C($R'$)($R''$)-OH} & \multicolumn{3}{X|}{pas d'oxydation ménagée (pas de H sur le carbone fonctionnel)} \\
\hline
\end{tabularx}
\end{center}

\textbf{Demi-équations électroniques de référence}
\begin{align*}
&\ce{Cr2O7^2- + 14 H+ + 6 e- -> 2 Cr^3+ + 7 H2O}\\
&\ce{MnO4^- + 8 H+ + 5 e- -> Mn^2+ + 4 H2O}\\
&\ce{R-CHO + 2 H+ + 2 e- -> R-CH2-OH} &&\text{(couple aldéhyde/alcool primaire)}\\
&\ce{R-COOH + 2 H+ + 2 e- -> R-CHO + H2O} &&\text{(couple acide/aldéhyde)}
\end{align*}

\textbf{Méthodes express}
\begin{itemize}
  \item \cle{Équilibrer une oxydation ménagée} : écrire les deux demi-équations $\to$ multiplier pour égaliser les électrons $\to$ additionner et simplifier \ce{H+} et \ce{H2O}.
  \item \cle{Reconnaître la classe d'un alcool} : compter le nombre de groupes alkyles portés par le carbone fonctionnel (1 : primaire ; 2 : secondaire ; 3 : tertiaire).
  \item \cle{Règle de Markovnikov} : lors de l'hydratation d'un alcène dissymétrique, H se fixe sur le carbone de la double liaison le \emph{plus hydrogéné} ; l'alcool majoritaire est donc le plus substitué.
  \item \cle{Rendement d'estérification} : $\eta = \dfrac{n_{\text{ester obtenu}}}{n_{\text{ester théorique}}}\times 100$ ; pour un mélange équimolaire acide $+$ alcool primaire, $\eta \approx \SI{67}{\percent}$ (alcool secondaire $\approx \SI{60}{\percent}$, tertiaire $\approx \SI{5}{\percent}$).
\end{itemize}
\end{bilan}

\begin{aretenir}[Checklist avant le Bac]
\begin{itemize}
  \item[$\square$] Je sais écrire et équilibrer la réaction d'un alcool avec le sodium, avec dégagement de \ce{H2}.
  \item[$\square$] Je distingue déshydratation intra (alcène, \SI{180}{\celsius}) et inter (éther, \SI{140}{\celsius}) et j'en écris les équations.
  \item[$\square$] Je connais l'équation générale de combustion complète d'un alcool.
  \item[$\square$] Je décris l'expérience de la lampe sans flamme : fil de cuivre incandescent, vapeurs d'éthanol, odeur de pomme verte (éthanal).
  \item[$\square$] Je sais identifier aldéhyde et cétone avec la DNPH, la liqueur de Fehling, le réactif de Schiff et le réactif de Tollens.
  \item[$\square$] Je prédis les produits d'oxydation ménagée selon la classe de l'alcool (primaire $\to$ aldéhyde $\to$ acide ; secondaire $\to$ cétone ; tertiaire $\to$ rien).
  \item[$\square$] Je sais écrire les demi-équations de \ce{MnO4^-}/\ce{Mn^2+} et \ce{Cr2O7^2-}/\ce{Cr^3+} et équilibrer une équation-bilan redox complète.
  \item[$\square$] Je sais que l'estérification est lente, limitée et athermique, et je calcule son rendement.
  \item[$\square$] J'applique la règle de Markovnikov pour l'hydratation d'un alcène dissymétrique.
  \item[$\square$] Je connais l'équation de la fermentation alcoolique du glucose et ses conditions (anaérobie, levures).
  \item[$\square$] Je vérifie toujours mes équations : conservation des atomes \emph{et} des charges.
\end{itemize}
\end{aretenir}

\findecours

\end{document}
